\quellseite{305}
\part[IV. The Material for Religious Instruction]{IV. The Material for Religious Instruction and Its Distribution Over the Several Educational Agencies}
\quellseite{307}
\chapter[§ 24. THE CATECHETICAL MATERIAL]{§ 24. THE CATECHETICAL MATERIAL AS SUGGESTED BY THE AIM OF RELIGIOUS INSTRUCTION}

*Zezschwitz ii 1, 1864. *Buchrucker, \emph{Grundlinien d. kirchl. Katechetik}, 1889, pp. 116–174. Sachsse, pp. 328–368. *F. Zange, \emph{Evang. Religionsunterricht}, 1897, pp. 73–142. A. Eckert, \emph{Der erziehende Religionsunterricht}, 1899, pp. 30–78. *O. Baumgarten, \emph{Neue Bahnen}, 1903, pp. 46–107. J. Gottschick, \emph{Homiletik u. Katechetik}, 1908, pp. 125–159. J. Berndt, \emph{Methodik i. Unterricht d. ev. Religion}, 1909, pp. 33–72. *R. Kabisch, \emph{Wie lehren wir Religion?} 1910, pp. 117–181. Achelis, pp. 365–410. Chr. Burkstuemmer, \emph{Der Religionsunterricht i. d. Volksschule}, 1913, pp. 138–174. Steinbeck, pp. 147–165.

The aim of catechetical instruction, as stated above, is (1) faithfully to imbed and anchor in the intellect all those \emph{sacred truths} upon which the mature congregation’s life fundamentally is based and by which alone it is constantly renewed and without a knowledge of which one can not possibly participate in its life; (2) to stir the emotions to a vital interest in those truths; and (3) to bend the will so that it may run in the paths in which the Holy Spirit in His own time will lead to personal faith and to participation in the life of the mature congregation. If this is true, then the question as to the material required for religious instruction has, at least in outline, already been answered. Those sacred truths which \emph{nourish, strengthen, and uphold} the congregation of Jesus Christ, form the material for religious instruction.

The life of the mature congregation is primarily a life of faith. Through faith it has grasped Jesus Christ and has become a Christian congregation; in faith it must ever embrace Christ anew if it is to maintain its character as such; and all of its activities must grow out of such faith
\quellseite{308}
to be acceptable to God. The briefest and most succinct summary of the faith living in the Christian congregation we have in Luther’s \emph{Small Catechism}. If this Catechism presented the theology of any one period of the Church’s history, it would, indeed, lose its usefulness for any other period; but among its chief merits (cf. §§ 14 and 26) we named the fact that it voices the fundamental thoughts of the Christian congregation’s life of faith which are the same in every period. And we also noted that it does this not in an abstract way but in the notes of a confession coming from the heart; nor is it ever forgotten that such faith is to be demonstrated in the daily life. When the adolescent youth has come to feel at home in the truths enunciated in the Small Catechism, it will always be at home in that world of thought which is peculiar to the mature congregation. The Small Catechism of Dr. Luther constitutes for this reason an important part of the material for religious instruction.

But the life of faith which animates the evangelical congregation since apostolic times with power occasionally hidden but ever breaking forth anew, is based upon a \emph{sacred history}. A life of faith and a Christian congregation have become possible only on basis of this sacred history involving God and men: the life of faith of the Christian congregation is the result of a history ranging over more than a thousand years, the history of human sin and grace divine. The rising generation would be incapable of understanding and appreciating the life of faith of the congregation in the present except as it is introduced to this sacred history. To this must be added that, though Christian faith is by no means mere \emph{fides historica}, facts of sacred history are nevertheless integral elements of the true Christian faith, as may be seen most clearly from the Second Article of our Christian Confession. Accordingly a summary of sacred history, such
\quellseite{309}
as we find in our Bible Histories, is an indispensable part of the material for religious instruction.

Both Catechism and Biblical History point back to \emph{Holy Scripture} from which they have been derived. While the Catechism deals with the present in that it gives expression to the life of faith in the Christian congregation, while Biblical History deals with the past in that it discloses the historical basis upon which the former rests, the Bible combines both aspects. It shows us the paths which God followed in order to establish fellowship between Himself and men, and the paths upon which today He seeks to lead men into this fellowship. It is the Bible in which the mature congregation is to search, from which it is to draw power and guidance, light and solace for all conditions of life. How can the rising generation be expected to make intelligent and independent use of the Scriptures unless it first receive suitable guidance and training? Hence, in addition to the Catechism and Biblical History, also the Bible must be an integral element of religious instruction.

The life of faith of the mature congregation is revealed particularly through the common public service. The communion of believers feels the urgent need of gathering in the public service for the purpose of renewing and quickening its life by the use of the means of grace and of expressing it in united confession, singing and praying. A type of religious instruction that fails to train the rising generation for intelligent and active participation in the worship of the mature congregation, overlooks a most important point. Everything, however, which is needed to equip the adolescent for such participation we find collected in the Hymnal, or \emph{Church-Book}, which, for that reason, is likewise to be accounted as part of the material for religious instruction.

The life in which the younger generation is later to take an independent part, is the life of the Evangelical, of the
\quellseite{310}
Evangelical-Lutheran Church. Now our youth sees itself surrounded by other denominations which likewise profess to be part of the great Church of Christ upon earth. Under the circumstances the rising generation will not consciously join our Church nor gladly advocate her interests unless it has been made acquainted with her character and her right of existence. This object is best accomplished by introducing the pupil into the chief periods of the \emph{history of the Christian Church}, and especially into the history of one’s own particular Church. One loves and appreciates only what one knows. That is the reason that church history in some form should be a part of the instruction given by the Church. — The mature congregation also has its own peculiar task and \emph{work}. If the rising generation, at a later day, is to take part in it with readiness and cheerfulness, it must know the why and wherefore, the spheres and occasions of such efforts. Only thus can an active laity, so much needed by the Church, be established. Hence an introduction to the several spheres of activity and, finally, also to the fundamental features of the constitution of the Church will be part of the material for instruction.

That the material enumerated does not all occupy the same level of importance, is readily seen. It is naturally divided into primary and secondary elements, instruction in Biblical History, Catechism, Holy Scriptures, and, in a measure, the Hymnal, belonging to the former. While stressing the primary elements, however, that which is secondary dare not be relegated too much to the background or ignored altogether. Preparation for participation in the whole life of the mature congregation being the aim, we dare not omit any part of the material enumerated.
\quellseite{311}
\chapter{§ 25. BIBLICAL HISTORY}

*G. v. Zezschwitz ii 2: \emph{Der biblische Unterricht in der Volksschule}, 1869, pp. 62–187. *K. Buchrucker, pp. 135–154. E. Sachsse, pp. 330–333. *F. Zange, pp. 73–142. A. Eckert, pp. 30–35. O. Baumgarten, pp. 27–39; 60–80. J. Gottschick, pp. 126–137. *J. Berndt, pp. 36–53. R. Kabisch, pp. 117–128. *A. Rude, \emph{Methodik d. Gesinnungs- u. Sprachunterrichts}, ¹¹912, pp. 14–84. Chr. Burkstuemmer, pp. 138 f.; 144–157. J. Steinbeck, pp. 171–185. L. Wangemann, \emph{Handreichung beim Unterricht d. Kleinen in d. Gotteserkenntnis}, 1861, pp. 6–43. *J. H. Schueren, \emph{Gedanken ueber d. Religionsunterricht i. d. christl. Volksschule}, ⁵1900. *K. Buchrucker, \emph{Der biblische Geschichtsunterricht}, ²1880. F. H. Kahle, \emph{Die Geschichte d. Reiches Gottes}, \emph{Handbuch zur unterrichtl. Behandlg. d. Biblischen Geschichte}, ⁷1892. *E. and O. Zurhellen, \emph{Wie erzaehlen wir den Kindern die biblischen Geschichten}, ²1910. M. Paul, \emph{Fuer Herz u. Gemuet der Kleinen}, ⁶1911. M. Reu, \emph{Unsere Erziehungsaufgabe im Licht von Joh. 17,6} in Kirch. Zeitschr. 1911. F. Rienecker, \emph{Gottes Wort bleibt in Ewigkeit}, 1912. *G. Fankhauser, \emph{Die biblische Geschichte in Sonntagsschule u. Religionsstunde}, ⁹1915. Religious Instruction by Means of Tales: Reukauf, \emph{Zur Lehrplantheorie d. geschichtl. Stoffe im Religionsunterricht d. Volksschule}, 1901. Landmann, \emph{Maerchenunterricht in Rein’s Handbuch}, ²1902. Just, Maerchenunterricht, 1906. Hiemesch, \emph{Der Gesinnungsunterricht}. \emph{Volksmaerchen als Gesinnungsstoff im I. Schuljahr}, 1910. Eick, \emph{Das Maerchen u. s. Stellung im Unterricht}, 1911. Schroecke, \emph{Maerchen und Kind}, 1911. Fuchs, \emph{Robinson als Stoff e. erzieh. Unterrichts}, 1893. Redecker and Putz, \emph{Der Gesinnungsunterricht im I. u. 2. Schuljahr}, ²1903. Hiemesch, \emph{Die Robinson-erzaehlung als Gesinnungsstoff der Volksschule}, 1907. *Rein, Pickel, Scheller, \emph{Das erste Schuljahr}, ⁹1908. H. H. Meyer, \emph{The Graded Sunday School in Principle and Practice}, 1910. Instruction in the Old Testament: E. Katzer, \emph{Das Judentchristentum i. d. religioesen Volkserziehung d. deutschen Protestantismus}, 1893. Floering, \emph{Das A. Test. i. evang. Religionsunterricht}, 1895. Lietz, \emph{Neue Aufgaben auf d. Gebiet d. Religionsunterrichts} in Rein’s \emph{Mitteilungen aus d. paedagog. Universitaetsseminar} in Jena vi, 1895. Hoffmeister, \emph{Die paedagogische Erziehg. d. christl. Jugend}, 1897. Meltzer, \emph{Das A. Test. im christl. Unterricht}, 1899. Reukauf, \emph{Grundlegung f. Auswahl u. Anordnung d. Unterrichtsstoffe},
\quellseite{312}
1900. Brammer, \emph{Neue Bahnen f. d. Religionsunterr.} i, 1900. E. Katzer, \emph{Judenchristentum} in Rein’s Handbuch, ²1902. Dreydorff, \emph{Quousque Tandem. Ein ernstes Wort wider d. alttestl. Geschichtsunterr.}, 1902. Kautzsch, \emph{Die bleibende Bedeutung d. A. Tests.}, ²1903. Kautzsch, \emph{Bibelwissenschaft u. Religionsunterricht}, ²1903. Staude, \emph{Das Alte Test. i. Licht d. Neuen}, 1905. Auffahrt, \emph{Die religiöse Frage und die Schule}, No. 2, 1905. *J. Koeberle, \emph{Heilsgeschichte u. religionsgeschichtliche Betrachtung d. Alten Testaments} in Neue kirchl. Zeitschr. 1906. C. F. Kent, \emph{The Origin and Permanent Value of the Old Test.}, 1906. *J. Orr, \emph{Problem of the Old Test.}, 1906. Geo. Hodges, \emph{The Training of Children in Religion}, 1911, pp. 167–217. A. J. W. Myers, \emph{The Old Testament and the Sunday School}, 1913. *E. Ziemer, \emph{Die voelkische, religioese und paedagogische Bedeutung d. Alten Testaments}, 1925.

ON THE AIM OF INSTRUCTION IN BIBLICAL HISTORY:

Grube, \emph{Charakterbilder aus d. H. Schrift}, 1854. Kolbe, \emph{Die Biblische Geschichte in Lebensbildern}. \emph{Katechesen f. d. Oberstufe}, 1900. Falcke, \emph{Die Heilige Geschichte in Lebensbildern}, ²1908. Thraendorf, \emph{Paedagogische Studien}, part ii, 1893. Bang, \emph{Das Leben Jesu}, Seine unterrichtl. Behandlg. i. d. Volksschuloberklasse. ⁴1902. Bang, \emph{Das Leben Jesu in historisch-pragmatischer Darstellung}, 1898. Bang, \emph{Sind die gegen e. historisch-pragmatische Behandlg. d. Lebens Jesu erhobenen Bedenken gerechtfertigt?} 1908. Niebergall, \emph{Jesus im Unterricht}, 1909. T. G. Soares, \emph{Heroes of Israel}. C. F. Kent, \emph{Historical Bible}. E. D. Burton and S. H. Mathews, \emph{Life of Christ}.

From the history of catechetical instruction it is evident that, from the beginning, it was customary in one way or another to give instruction in Biblical History; but that it did not receive attention in the same measure as instruction in the Catechism; that such instruction was almost entirely forgotten during the Middle Ages; and that even in the Church of the Reformation it was only by slow degrees accorded wider reception. From these facts the conclusion might be drawn that Biblical History, even though rightfully a part of catechetical instruction, was deserving of secondary instead of primary place. Moreover, when we bear in mind that Biblical History leads the pupil into the past while the rising generation needs primarily to be made acquainted with the life of faith of the mature congregation in the present,
\quellseite{313}
we appear to find an additional reason to give first place to the Catechism. Notwithstanding, dogmatical and, still more, pedagogic grounds require that preference be given to Biblical History. The dogmatical reasons are (a) that God’s revelation by acts, presented in Biblical History, has been the fundamental and creative factor in a much higher degree than his revelation by words which finds more extensive expression in the Catechism, and (b) that Christ, the center of our faith, is not understood aright unless He be viewed as starting-point and creator of the whole New Testament period of salvation and, likewise, as the goal of a more than millennial history of God’s relation to man. The pedagogic reasons are (a) that the principle of intuition as determined by the mental development (pp. 204, 217, 226) requires plastic images such as are offered nowhere else so copiously as in Biblical History and that it especially requires such plastic images as the indispensable basis for instruction in the Catechism, and (b) that the gradual unfolding of the soul of the child plainly suggests the use of historical material as the primary means of instruction (p. 250). So much is certain, however, that such instruction must not bury the youth of the Church in the past: always it must be viewed and turned to account in its bearings upon the present life of faith. Where this is neglected, the knowledge imparted is dead. While it may prove of value in the future, it is certain that it offers the pupil nothing for the present.

Instruction in Biblical History being a fundamental factor should be given in every grade. This has been denied by the school of Ziller (Ziller 1817–1882) or by the Neo-Herbartians (p. 151) who appointed in its place a dozen tales for the lower grade, and the story of Robinson for the upper. Also in America the attempt was made to introduce such material. It was thought that even the most simple
\quellseite{314}
Bible story was too difficult for the capacity of children six or seven years of age; that its concepts were too remote from the concept sphere of children; that the language of Scripture which was to be retained to the fullest possible extent, was not that of the children; also that such historic material as there furnished failed to stimulate the phantasy sufficiently and left the mind atrophied. But these arguments fit rather the faulty handling of the matter as found here and there than the matter itself. They remind us, indeed, that for the lower grade only such sections should be selected which neither transcend the capacity of the children nor are too remote from their concept sphere — that is, family stories; and that they should be couched in such language as the children can understand, even though it should differ from the language of Scripture; likewise, that the children’s phantasy should be duly stimulated. If these requirements are fulfilled, it is quite appropriate to let religious instruction at the age of from six to seven years begin with Bible stories (pp. 250 f.). The use of tales in religious instruction involves the risk, by no means to be underestimated, that the child loses confidence in the truth of religious material as taught in the class period, and that he becomes inclined to consider the subsequent biblical stories as mere tales, or at least to view with skepticism the miraculous element in them. Possibly this may not take place during the school period, but at a later day the thought might easily suggest itself to the youth that, as he formerly was led from tales to biblical stories, so it is now his privilege, or duty, to mount above this stage, to view historical elements as symbols, and thus to penetrate to the right “spirit” of Scripture and Christianity. We maintain, then, that Biblical History is to be taught already in the lower grade. Of course, every story must form a unit by itself; it dare not be coupled with others. Nor is it necessary to connect one
\quellseite{315}
story with the other in point of time or logic; the mental development of the child precludes such treatment during these years almost entirely.— A different situation is presented by still younger children, who are either instructed by the mother at home or in the “Beginners’ Classes” in Sunday School; or by children of from six to eight years who are mentally backward or who have grown up in homes where the idea of God is an unknown quantity. In such cases it is the teacher’s duty to arouse and to strengthen the consciousness of God by a preliminary course, as it is found, for instance, in the first course of the \emph{Wartburg Lesson Helps} (For Beginners). But even this preliminary course may be largely joined to subjects of a biblical character.

Ziller proposed the following twelve tales from the collection of the Grimm brothers: (1) The Star Dollars; (2) The Three Sluggards; (3) The Three Spinners; (4) The Straw, the Coal, and the Bean; (5) The Seven Goslings; (6) The Little Hen and the Little Rooster, or as alternate, The Death of the Little Hen; (7) Wolf and Fox; (8) The Rabble; (9) The Town Musicians; (10) Wren and Bear; (11) “Fundevogel”; (12) Poor Man and Rich Man. Two of these tales, viz., The Three Sluggards and The Three Spinners, were soon removed from this list as of doubtful educational value, “Snowwhite and Rosed” and “Mother Frost” being put in their place. Rein, of Jena, expunged also “The Rabble” and “Wren and Bear”, replacing them with “Red Ridinghood” and “The Ear of Rye”. He also received into the collection both “The Little Hen and the Little Rooster” and ”The Death of the Little Hen”, and added “The Sweet Broth”, thus increasing the number to fourteen. In Das erste Schuljahr, by Rein, Pickel, and Scheller (Leipzig, ⁸1908), we find didactic sketches on these tales (pp. 193–245) which in The Progressive Educator (1898–1899) were translated by Dr. Abbott into English. The fundamental principles on which these projects and many of the numerous other liberal Graded Series of the last two or three decades are based, are totally different from those of the preparatory course of the Wartburg Lesson Helps. This appears at once from the titles of the first twenty-two lessons: (1) God Made All Things; (2) God Made Adam, the First Man; (3) What God Has Given Me; (4) God Cares For Me; (5 and 6)
\quellseite{316}
God Protects Us by His Angel; (7) How God Has Made Me His Child; (8 and 9) How We May Speak With Our Heavenly Father; (10) God Hears Our Prayers; (11) Jesus is Our Good Shepherd; (12 and 13) God Sees and Knows Everything; (14 and 15) God Punishes Us if We Do Evil; (16) God Keeps His Promise; (17) We Should Thank God When He Has Helped Us; (18) God Looks Upon the Heart; (19) What We Should Do on Sunday; (20) We Should Help All Those Who Are in Need; (21) The Bible Is the Best Book; (22) How Beautiful It Is in Heaven.

Not rarely the attempt is made to exclude the Biblical History of the Old Testament where the necessity of teaching New Testament Biblical History is conceded. And, indeed, if modern criticism were right, according to which this part of the Bible is highly legendary and mythical or at best a mere “Jewish Chronicle,” then instruction in Old Testament history would be less important than national history and certainly would deserve no place in the curriculum of the church school. But let us not forget that the Old Testament was the Bible of Jesus; that through its connection with the history of the New Testament it has become part and parcel of the history of salvation; that it is preparatory to the New Testament history of revelation to such an extent that many single facts of the New Testament simply can not be understood without the Old; and that even Jesus Christ and His life work can not possibly be fully appreciated save as He is recognized as the goal toward which for more than a thousand years God directed the course of human history, and as the Saviour who in the fulness of the time should still the longing of the whole pre-Christian world (Gal. 4:4). These facts indicate the necessity for an introduction into Old Testament history, and, at the same time, the duty of relating everything in it to Christ and of interpreting everything with a view to Him as its goal. Moreover, it is just the men of the Old Testament—tainted with sin and infirmity though they were—that stand before us as sublime types of an evangelical life of faith: their
\quellseite{317}
faith in Yahweh unshaken, their hope unswervingly turned toward the promises of God, their courage persevering in spite of bitter adversity. The teacher of religion who knows that his object in all of his instructions must be to train his pupil, could not wish for better examples. And finally, God’s holiness, His grace, and likewise, a number of fundamental moral principles, as, for instance, “Righteousness exalted a nation, but sin is a reproach to any people,” can be most clearly exhibited in the light of the history of the Old Covenant, and thus be impressed upon the adolescent youth.

Therewith the question as to the purpose of instruction in Biblical History has already been answered in part. Naturally such purpose dare not be in conflict with the object of all the training efforts of the Church. Instruction in Biblical History contributes its share to the achievement of the general object. But if it is to be an integral factor in the process of imbedding and anchoring in the threefold inner life of the adolescent youth the sacred concept material with which the mature congregation nourishes its life, such instruction should indeed impress upon the mind the several features of the more important stories of the Bible; but, that attained, it should go further. If the catechist were satisfied when his pupils can tolerably well recite the contents of the Bible stories, he would have accomplished no more than a mechanical, memoriter, appropriation of the material; it would not have been digested through intensive thinking, it would not have stirred the emotions nor set the will in motion, in other words, the inner life would not have reaped any benefit whatever, at least not so far as the present time is concerned.

For this reason the teacher must take a further step. He should draw forth and impress upon the children’s hearts the holy divine thoughts of permanent value which are contained
\quellseite{318}
in the individual Bible story, whether they have reference to the life of faith in its God-ward or in its man-ward aspects. A particularly welcome opportunity to do this presents itself when the teacher joins a number of stories for the delineation of the character of such outstanding persons as Abraham, Moses, or David, and subsequently in a higher sense, even Christ Himself; he must carefully draw forth the typical traits that to this day continue to be essential to a truly evangelical life and must cast light upon God’s relation to these persons and upon their attitude toward God and their fellow-men. The life of such types must be made an object lesson for children; they must understand and vividly realize the seriousness and grievousness of sin, and recognize that sin is a most momentous reality; they should learn once for all that sin and faith are incompatible; and they should acquire an abiding horror for sin. When the Old Testament is under consideration and occasion demands it, the teacher will point out the difference between Old Testament and New Testament morality in order to obviate subsequent misunderstandings, and to show the progress of the divine revelation; for the life in which our catechumens are to join with full consciousness of what they do, is not that of the Old but of the New Testament, not of the congregation under the Law, but of that under the gospel.

Let us illustrate the purpose of instruction in Biblical History by a few examples. That the children know the order of the several creative acts of God is commendable: the evident gradation from the lower to the higher clearly demonstrates that the creation of the world was preliminary to that of man; more important, however, than even this truth, and of greater moment for their whole life, is that they learn from this story that God is the creator, and therefore the Lord, of the whole world so that they need not fear anything in heaven or on earth provided they have Him for a friend. It is well and good that the children should be conversant with the details of man’s fall; but it is more important that this story should serve to impress upon their minds beyond peradventure of doubt that “our God is not a God that
\quellseite{319}
hath pleasure in wickedness and that the evil man shall not sojourn with Him.” It is well if the class is able to state the dimensions of the ark or even the dates connected with the deluge; but it is of more importance that every member of the class should carry home the vivid impression: \emph{God is not mocked}; for \emph{whatsoever a man soweth, that shall he also reap}. The catechist may ask the children to recite all the events that preceded the dispensation of the Law on Mt. Sinai, especially since these are closely connected with \emph{the main point—the revelation of God’s holiness}; but he has failed to do justice to his task if he does not present the story of the giving of the Law \emph{with such clearness and vividness, with such impressiveness and earnestness} that Law and divine holiness are inextricably blended in the juvenile understanding and phantasy, and that the divine command, “Ye shall be holy; for I, the Lord your God, am holy,” is \emph{indelibly engraved upon their consciences}.—The treatment of the several events in the life of Abraham, the careful development of the involved moral and religious truths, and the application of these to the life of the child, are necessary and \emph{most important duties of the catechist}; but when all the stories, one by one, have been considered and turned to account, let the whole be summarized for the purpose of showing, with the life of “the father of the faithful” as object lesson, what is really meant by a life of faith; how God trains us for it; and how it is evidenced in the various conditions of life. “The obedience of faith drew Abraham into a foreign country; in the humility of faith he yielded to Lot, his cousin; in the strength of faith he smote with 318 men four heathen kings; in the perseverance of faith, reason and nature to the contrary, he rested in the word of promise; in the boldness of faith he bartered with God and implored Him to spare Sodom; in the joy of faith he received, named, and circumcised the son of promise; in the fidelity of faith he conformed, at the behest of God, to the will of Sarah and cast out Ishmael with his mother Hagar; in the gratitude of faith he planted, at the spot where Abimelech had solicited his friendship and accepted his gifts, a tamarisk tree to the honor of the ever faithful God” (Delitzsch); in the steadfastness and conquering power of faith he showed himself equal to the supreme test, the sacrifice of his son; in the assurance of faith he secured, for himself a family burial place; in the fidelity of faith he refused to take a wife for his son from the women of Canaan; rich in faith he is gathered with his fathers—a true “father of the faithful.”—It is a beautiful sight, and one that gives joy to the heart when the children know the details of the passion
\quellseite{320}
and death of our Saviour, when every feature of this “story of stories” is \emph{impressed upon their memory}; but in spite of such commendable exactness and accuracy, the catechist has totally failed if he has not aroused something better in his children than \emph{a horror of “those mean Jews.”} The result of the contemplation of all the stations on the way of suffering should be \emph{a “revelation” of both the holiness and the love of God}. It should reveal the holiness of God; the scenes enacted in Gethsemane and on Calvary must indelibly impress upon the children that God is a holy God who does not reduce his demand but exacts the full penalty, that “it is a fearful thing to fall into the hands of the living God,” that \emph{God’s wrath and judgment are terrible realities}. And \emph{it should reveal the love of God}; by his instruction the catechist should withdraw, as it were, \emph{the veil and disclose the fatherly love of God} to His human children that the souls of the children may be grasped and won by this love.—Compare Wartburg Lesson Helps, especially the Intermediate and Junior Grades.

Even now, however, we have not yet fully set forth the purpose of instruction in Biblical History although prevailing conditions may restrict the teacher to what has been dealt with thus far. After a careful treatment of the individual Bible stories; after combining them, to the extent necessary, for the purpose of portraying Bible characters; after forming them into groups (History of Creation, History of the Patriarchs, History of Moses, Establishment of the Nation of Israel, etc.), the groups themselves will have to be combined into a final unit in order to construct—not a history of the kingdom of God, which would be beyond the scope of juvenile instruction, but a history of salvation. This would show the historic development of salvation in which the congregation of the faithful rejoices at this day and which should become a conscious possession of the adolescent youth as well. Accorded such treatment as here described, the several stories will lose their incoherent character: they will rather appear as links of one unbroken chain which leads to our salvation. In the eyes of the young the grace of God will be magnified more and more (cf. Micah 7:18; Isa. 40:9) when they behold how for thousands of
\quellseite{321}
years it never swerved from its one purpose, i.e., our salvation, how it overcame all the obstacles put in its way by the sin of man, how it won the grand victory on the cross, and how since that day it has without wearying sought both Jews and Gentiles to grant them salvation. The children will prize more highly this salvation, and they will begin to realize that he who rejects Christ rejects Him who is the center of all history, the only mediator between God and man, the one true Saviour. No one can deny that such a treatment of the history of salvation is of great educating value. If properly carried out, the important truth forces itself upon the youth that only the Evangelical Church, which places sin and grace, grace and trust, \emph{“sola gratia,”} \emph{“solus Christus,”} \emph{“sola fides,”} into the center of faith and life, is the orthodox church of God upon earth; for, ever since the fall, the relation between God and man has turned on these truths. Where the teacher has faithfully taught the several stories and paid some attention to the main lessons involved, this final object of instruction in Biblical History can be attained even where there is no regular parochial school, as the author knows from experience.

A few hints in regard to such treatment may be found acceptable. The Old Testament period is divided into two large sections: the history of the human race, and the history of the chosen people. We deal first with the human race as a whole. It springs from the one man Adam. Created in the image of God, he as well as Eve enjoyed fellowship with their Creator; they rejoiced at the prospect of His coming, and in His presence they felt blessed; that is, they had all they needed. Death and evil did not exist, and labor was known only as a joy. When the time had come that their fellowship with God was to become a permanent blessing, they decided against God. Nature of sin. Thus the fellowship with God was lost: they are ashamed; they are filled with fear (opposite of blessing); they lie. God must punish sin, for He is the Holy One. Punishment is inflicted as previously announced; words to the serpent, to the woman, to the man: pain and submission; the ground cursed; bodily death; expulsion from paradise, the place
\quellseite{322}
of God’s presence; subjection to Satan’s power, lost, condemned. But \emph{God is also gracious}: \emph{in grace He seeks the sinner}; \emph{in grace He directs him to a blessed future}; the word of \emph{the “seed of woman”} is broad enough to bear in its bosom all future salvation. Satan thought that he had gotten humanity into his power forever; but the descendants of woman shall resist him. Conflict arises. It would be a hopeless battle, however, were it not for the grace of God which provides ultimate victory through a descendant of woman.

Without murmuring and complaint, i.e., in repentance, \emph{Adam and Eve submit to the inflicted penalty}, believing the word of promise (cf. the naming of the first-born son). This was their solace amid the toil and grief on earth, now a vale of tears. Thus in the midst of the material life, \emph{a life of faith} took its beginning: Adam and Eve typify for us \emph{the laying hold of the grace of God}. But this faith coupled with repentance does not at once pass over to the children; they rather \emph{inherit sin}. Proof: Cain and \emph{his family}. From such a generation the one who was to bruise the serpent’s head could not arise. While Seth is godly, also here sin increases, and Lamech longs for the Redeemer with great intensity. Of both branches of the human race it is true: in their going astray they are flesh. From such a race redemption can not be expected: it is \emph{ripe for destruction}. Deluge.

Yet, the Redeemer was to arise from among the descendants of woman. For this reason, God, \emph{eternally faithful}, preserves Noah and his family. For a second time the whole human race is to develop from one man. Will it be more godly and trust in the promise? So it would seem: Noah consecrates the renewed earth with a sacrifice. But a second fall dispels the illusion. In one of his sons the sin of the destroyed race has maintained itself to a horrifying degree; sin is transplanted to the new earth. What is to become of the race when the mocking of the father becomes a practise and unchastity passes unrebuked? Noah understands the situation and pronounces a curse upon Ham and his family. Now the Redeemer could come forth only from Shem or Japheth. Shem was destined to be the ancestor of the Redeemer; but Japheth was to partake of His blessings. As from the beginning a word of grace comforted Adam’s race, so here, for God is exceedingly rich in mercy. Once more sin spreads. In the old world the family of Cain was its nucleus, in the new it is Ham’s: Tower of Babel; Nimrod; Ashur; after the dispersion ingratitude, indifference to God, idolatry. In spite of the divine judgment that had come upon the world, such is the consummation of the history of the race descended
\quellseite{323}
from Noah. It is clear that \emph{such a race cannot be an active redemptive factor}. It does not even allow itself to be trained to look for a Redeemer. If salvation is to come to the world, it can come forth from \emph{God alone}.

But will God once more be gracious? Oh, great is His mercy; \emph{He is faithful and keeps His promise}. As in the first generation he preserved Noah as His seed, He now makes use of Abraham. From him the chosen race shall spring through which \emph{the blessing shall come} upon all peoples. But why just Abraham? Not because he had deserved it; in unmerited love God elected him in order to make him, through a \emph{steady discipline of faith}, an instrument fit for His hands. Here, again, the \emph{divine promise of grace, and human faith} stand correlated, impressing upon us upon what our salvation rests and how we appropriate it. And what a faith! Abraham is justly called the father of the faithful. The promise passes from Abraham to Isaac who at his sacrifice already had disclosed the fundamental disposition of his heart; from Isaac to Jacob who had been disciplined for his task in a hard school of suffering and had learned to rely less and less upon his own power and wisdom and desired at length but one staff on which to lean: the grace of God. The promises are being fulfilled, but not in a straight line, not according to man's thoughts. The promised land indeed becomes the first home for the bearers of the promise: trusting in the promise of God, Abraham purchases a family burial plot and thereby lays his hand upon that country as his own. But the last patriarch is driven by famine to Egypt where his habitation was prepared by a son whom his brethren had sold but whom God highly exalted. His gaze fixed upon the promised land, i.e., the promise, he dies with the words: “I have waited for Thy salvation, O Yahveh,” after pointing out Judah as the one from whom should come the Redeemer, the Prince of Peace.

Jacob’s family was led to Egypt in order to develop into a nation. But lest prosperity should cause it to forget the promise and to look upon Egypt as its home, the favor of the kings of Egypt turns into hostility; beneath dire oppression it is to become mindful of the God of its fathers and to long for the deliverer. And behold, God appoints one of Jacob’s offspring as the deliverer of the people. He saves him from death in a miraculous manner, trains him at the royal court and in the wilderness to become His instrument, and appoints him solemnly as the deliverer of the people. Amid mighty signs and wonders which are to prove that the God of the fathers
\quellseite{324}
and of the promise, though silent for four centuries, is the only true God, Moses leads Israel out of the house of bondage: rescued, it chants a hymn of praise on the other side of the sea. Indeed, \emph{Moses was a deliverer of his people: not the true one}, however, but merely a type. Israel needs to be trained for \emph{the true deliverer}. Accordingly God reveals to it His Law so that it may behold His holiness and its own sin. And lest it should be driven to despair under the Law, God, through a system of sacrifices, directs it to Him who is to offer the true sacrifice. As a type of a future reality, fellowship is made between God and this chosen people. But the generation first delivered will not permit itself to be trained, and perishes in the wilderness. The new generation, though it sins often and grievously, is on the whole a different people. For this reason Moses, as Balaam before him, is permitted to pronounce a blessing upon it and to proclaim anew the coming of Him of whom he was merely a type. Joshua is allowed to lead Israel out of the wilderness into the promised land, a type of Him who is to lead into the land of rest.

As Moses before him, Joshua exacted the promise from the people not to mingle its blood with that of the pagan nations; but Israel did not keep its promise; oppressed now by this nation now by that, it had to realize that the true deliverer had not yet appeared. Whenever Israel repented, God gave it a temporary deliverer. Thus it happened that \emph{the Judges without exception became types of the true deliverer}. All these had been set over Israel by God; but when eventually it desires a ruler after its own heart, it receives Saul who contrary to early appearances does not prove the right man: he does not persevere in \emph{faith and obedience} to the end. Then God elects a man after His own heart, obscure and lowly, the descendant of shepherd folk, but godly and with a heart emboldened by faith. Like all His servants, \emph{God trains him in the school of suffering}, making him, the training completed, a mighty king in Israel, who saves it from all its enemies, widens its boundaries, fills its treasury, receives tribute from neighboring kings, and transfers the tabernacle, the habitation of God, to Jerusalem. If at any time, it is now that the thought might be entertained in Israel: the promised One has come! \emph{David is the true son of Judah and the seed of the woman!} But his deep fall discloses the fact that he is not; yet, Nathan is permitted to tell him that the deliverer is to be one of his descendants. Is it perhaps Solomon whose very name reminds one of the Prince
\quellseite{325}
of Peace, whose ships traverse the seven seas, whose wisdom is known even in far away Arabia, who is privileged to build a temple for the Lord? Doubtless some believed this; but when Solomon oppressed the people and took “strange women” as wives, it became clear that \emph{the true son of David} was still to come.

When the kingdom was divided, and calf-worship and, afterward, in defiance of Elijah’s warning, the worship of Baal was countenanced in the land; and when even Judah did not follow the Lord faithfully, it became evident that Israel’s nature had only been kept in check under David and Solomon, but not really renewed. The prophets warn with great zeal and make an effort to induce the people to renew their obedience to the law. The less it shows itself inclined to take to heart the warning, the more pronounced becomes the prophetic warning of the coming judgment, and the more unwaveringly the eyes of the prophets are turned to the future from which they expect the deliverer; their prophecies concerning the Redeemer become increasingly definite. The judgment takes place; for God is not mocked. The Kingdom of the Ten Tribes, for its exceptional hardness of heart, is led captive to Assyria; Judah, better than the sister kingdom of Israel, but also exceedingly wicked, is exiled to Babylonia after it had shown its unwillingness to heed the call to repentance issued by Isaiah and Jeremiah, and to profit by the example of Israel. Because it was not to perish utterly—for God is faithful to His promise and gracious above all that we ask or think—there are given to it the prophets Ezekiel and Daniel, and Isaiah’s book of comfort becomes its stay. Nor has it been thrown into the crucible in vain: it learns to pray and cry to God. Accordingly, God, in His faithfulness and grace, restores a remnant, with a member of David’s house at the head. After the return, it makes a covenant with God, becomes the people of the Law, and clings, at least outwardly, to the covenant of Sinai. Idolatry is definitely done away with. The people was to regard its deliverance from Babylon as the last great type of the great promised redemption which also Malachi, after the exile, expects in the future. In the subsequent period three classes gradually develop among the people: the Sadducees who compromise with the world and pay little attention to the prophets; the Pharisees who observe the Law, but do not understand its spirit and expect to realize their own glory in the expected Messianic kingdom; but also a throng of devout believers by no means negligible in number, who look for the hope of Israel, such as Zacharias and Elizabeth, Simeon, Hannah, Joseph and Mary, the
\quellseite{326}
shepherds in the field. These pious souls moved, and had their being, in the prophetic books and were filled with \emph{hope and faith}. Zacharias utters this faith in his “Benedictus,” Mary makes it the theme of her song in Elizabeth’s house. This group constituted the fruit of Israel’s training at the hands of God. Now the time was fulfilled, and \emph{God kept His promise}.

As of old, angels appeared—a sign that something great was impending. The forerunner is promised; the birth of the Redeemer Himself is announced. The birth of John warrants the truth of the words spoken to Mary. We come once more upon the old principle: \emph{God promises; man must appropriate the promise in faith}. In Bethlehem the Saviour is born; angel lips proclaim the fact. He is the promised Christ, or Messiah; and He can be Saviour because He is, at the same time, Lord. Here is the goal of the whole movement: \emph{God Himself becomes man} in order to bring about a true deliverance. The Son of God and the “Seed of Woman”—this is the dominant fact in all subsequent events. A human child is born; yet the angels are stirred by His birth. Like any other child in Israel, He is circumcised and presented in the temple; yet extraordinary words are spoken of Him by Simeon. The Son of God, He is worshipped by the wise men; yet as a poor human child He must take flight from Herod. Wonderful words He speaks in the temple concerning His fellowship with God; yet, as a true son He deports Himself in the parental home in Nazareth. He matures in solitude; yet John the Baptist salutes Him as his superior, and God Himself proclaims Him as His dear Son. By stepping into the waters of Jordan he declares that He shall fulfill all righteousness and become the deliverer even though it will be necessary to step into the waters of suffering.

As the future deliverer he frustrates the designs of the tempter. He gathers about Him His first disciples who, without exception, recognize Him as the promised Messiah. He enters upon His redemptive office with the proclamation that He is come to establish the promised \emph{Kingdom of God} in which deliverance can be found from each and every ill, from sin, death, and the devil, and through His miracles He proves that He really can work such deliverance; but as \emph{indispensable condition of entering His kingdom} he requires regeneration, faith, i.e., the firm confidence that He alone is the true Saviour and Redeemer. That is the reason why “Thy faith hath saved thee” occurs so frequently. Faith is the condition, indeed; but He Himself fulfills the condition, for His message of the kingdom of God is the means
\quellseite{327}
whereby He would produce faith. By His miracles He would invite men to enter into personal fellowship with Him. Who thus has entered His kingdom through faith is \emph{blessed} even in the midst of tribulation, and may well be expected to live a new life (Sermon on the Mount). But the rulers and the larger part of the people will not permit this faith to be wrought in their hearts; in self-righteousness they believe that the only deliverance needed is that from Roman dominion. Their enmity aroused, they decide to kill Him. At Caesarea Philippi the fact is revealed that but few persons recognize Him as \emph{the true Redeemer} and Son of God. For that reason Jesus begins to proclaim His impending passion for which He as well as the disciples receive strength through the transfiguration. In order to carry out His mission, i.e., to give His life as a ransom for many, He goes to Judea; the resurrection of Lazarus fans the flame of the Pharisees’ hatred: they decide when and how to put Him out of the way. Permitting Himself to be anointed, He interprets this act as a symbol of His impending burial. He presents Himself at Jerusalem as a lamb that is led to the slaughter; but not without demonstrating once more by the manner of His entry that He is the promised King. He weeps over Jerusalem which does not know the things that belong to her peace; He cleanses the temple: for the last time He appeals to the Pharisees, earnestly impressing upon their hearts the Father’s love and their contemptuous response to it (parable of the laborers in the vineyard), and silences them by the question, Whose son is Christ? He prophesies the destruction of Jerusalem, the end of the world, and final judgment, always inviting the hearers to faith. He celebrates His last Passover with the disciples: He unmasks the traitor; and, certain of the redemptive character of His impending death, He distributes His body which is to be given into death, and His blood which is to be shed in death, and enters upon His great suffering which shall fully reveal that He is the \emph{true Deliverer}.

As our substitute, He takes upon Himself the burden of the whole world’s sin together with the wrath of God. He surrenders to His enemies in order to fulfill the Scriptures, but first reveals Himself once more as the Almighty One: voluntarily he goes into death. His heart is set on Calvary, the consummation of His redemptive work: He permits Himself to be sentenced to death and to be led to Calvary, from the cross He asks for forgiveness in behalf of His murderers, is forsaken of God, dies for the human race with the cry of victory, “It is finished”. Yes, redemption was finished: performed was the sacrifice; fulfilled were the Scriptures; bruised was the head of the Serpent, although
\quellseite{328}
His own heel was bruised in the process! The goal of the whole history of mankind down from the Protevangel, the antitype of Moses, Joshua, and David, and of all the sacrificial lambs of the Old Covenant—we find in the Crucified One. Here all threads converge. How great and world-embracing the work of our Redeemer must be; how great He Himself, since in Him we find consummated a time of preparation ranging over four thousand years; since in Him is comprehended the salvation and hope of all the world!

But is He really the Redeemer; is the Father satisfied with the sacrifice; is the guilt of humanity really atoned for, covered in His sight? Yes, for on the morning of Sunday He awakes Him from the dead: in vain are seal and guards whereby the rulers endeavor to keep life in the grave. While the rulers defraud the people of this glad message, the Lord appears to His disciples, even to doubting Thomas. He assures them of the reality of His resurrection; He shows them \emph{the necessity} of His passing through suffering and death in order to become the Redeemer of whom Moses and the prophets have spoken; He delivers to them as fruitage of his work the message and the gift of \emph{peace} between God and man; He accepts the confession that He is God, worthy to be worshipped; He restores the apostle to his apostolate; He shows the disciples by the manner of His coming and going that the manner of His presence with them will be changed thereafter; He commands them to take His gospel into all the world and to disciple \emph{all nations through Baptism and teaching} for which purpose He pledges to them His presence, the gift of the Holy Spirit, and continued support at His hands through miracles. He appoints them His witnesses—\emph{witnesses of His death and resurrection}, and then ascends to heaven in order to enter upon the government of the world and to gather on earth, by the Spirit, a congregation of the faithful which is blessed through Him and acknowledges Him as its Redeemer and Lord.

Such is the conclusion of the first chief part of the New Testament story, which deals with the life of the Lord. The Lord keeps His promise and on Pentecost sends His Spirit who gathers the disciples into a congregation, or Church, and inspires them with courage and joy of true discipleship. Peter proclaims the two fundamental facts which continue resounding in all the sermons of the apostles that \emph{Jesus, through suffering, death, and resurrection, has become both Lord and Christ}, and that on but one condition men may receive remission of sins and the Holy Spirit, namely repentance (i.e. sorrow for sin and faith) and Baptism in His name. Through such preaching
\quellseite{329}
three thousand people are added to the Church with the result that the holy nation typified by Israel is now established. But no one shall lose sight of the fact that the Church is to be holy: Ananias and Sapphira suffer sudden death; God himself sets an example of most thoroughgoing church discipline in order to bring about a permanent impression. Like her Lord and Master, the congregation must pass through \emph{cross and suffering}: there is no kingdom of glory on earth, but only one that bears the cross. Stephen exemplifies especially this truth; but also the other truth that true believers can be full of cheer even in the midst of death because their path leads to Jesus who sitting at God’s right hand shall receive their souls. The stoning of Stephen is lightning foreboding the storm; but even persecution serves to promote the spread of the Christian Church. The most relentless persecutor is converted by Christ and appointed as an apostle to the Gentiles. However, before the work among the Gentiles can be begun, Peter and the other disciples must be convinced that also the Gentiles have been called to membership in the Church of God and that they need not pass through Judaism or submit to circumcision. This is effected through a special revelation in Joppa and an object lesson in the house of Cornelius. Paul on his missionary journeys takes the gospel to the Gentiles, observing meanwhile how the Jews spurn the Saviour. Such rejection spurs him on to greater missionary efforts among the Gentiles. He goes as far as Rome in order to sound in the capital of the world the message of Him through whom fellowship has been established between God and the whole human race, and in whom \emph{everyone that lays hold of Him in faith, shall have life and shall have it abundantly}. The history of salvation finds its consummation in a congregation of believers, composed of Jews and Gentiles, \emph{in which the gospel of the free grace of God in Christ is preached}, i.e., in a congregation bearing all the marks found in the Evangelical-Lutheran Church of today; for the greatness of our Church consists in this that the gospel of Christ \emph{the only Saviour is preached within her in its purity}.

\emph{Augustine} deserves credit for being the first to recognize clearly the necessity for instruction in the history of salvation (cf. p. 34). More recently, one of its warmest advocates, both in theory and practice, was Buchrucker who was influenced by the ‘historical’ theology of Hofmann. Compare with the above sketch Buchrucker, Katechetik, pp. 142 ff.; 222 ff.

If treated in the manner here described, it is manifest that instruction in Biblical History must cease to be merely
\quellseite{330}
ancillary to the teaching of the Catechism, and to be used merely as a welcome collection of stories to illustrate its truths. It rather occupies the same level with the Catechism as a distinct factor. Instruction in Biblical History is to be given before there is specific instruction in the Catechism; from the seventh to the fourteenth year it should be given in every grade. Treated in connection with “Introduction to the Bible”, and from the viewpoint of a history of salvation it is to be recommended even for the confirmed; compare the Senior Department of the \emph{Wartburg Lesson Helps}. \emph{Biblical History, and again Biblical History, must be our motto}; but Biblical History taught in the right manner—according to principles as here elaborated. With the right sort of treatment, especially when the truths drawn from the several stories, or groups of stories, are clothed in the words of the Catechism (cf. Wartburg Lesson Helps, Intermediate and Junior Grades), instruction in the Catechism will not lose, but rather gain, it will become easier and more effective.

\quellseite{331}
\chapter{§ 26. THE CATECHISM}

*G. v. Zezschwitz, ii 1, \emph{Der Katechismus}, 1864. K. Buchrucker, pp. 119-135. E. Sachsse, pp. 338-365. *F. Zange, pp. 73-142. A. Eckert, pp. 36-47; 53-55; 116-163. O. Baumgarten, pp. 68-72; 82-84. *J. Gottschick, pp. 142-158. J. Berndt, pp. 59-68; 99-105. R. Kabisch, pp. 137-139; 240-246. E. Chr. Achelis, pp. 420-443. A. Rude, pp. 100-117. Chr. Burkstuemmer, pp. 157-169. J. Steinbeck, pp. 201-218. K. Knoke, \emph{Ueber Katechismusunterricht}, 1886. G. Heimerdinger, \emph{Zur Reform des Katechismusunterrichts}, 1890. G. v. Rohden, \emph{Ueber christocentrische Behandlung des luth. Katechismus}, 1891. *L. Schultze, \emph{Katechetische Bausteine z. Religionsunterricht in Schule u. Kirche}, ⁵1891. A. Braasch, \emph{Reform d. Religionsunterrichts in d. Volksschule}, 1891. Malo, \emph{Zur Katechismusfrage gegen von Rohden}, 1892. W. Bornemann, \emph{Der zweite Artikel im Lutherschen kleinen Katechismus}, 1893. G. Voigt, \emph{Die Bedeutung des christl. Religionsunterrichts}, 1895. J. D. von der Heydt, \emph{Der Religionsunterricht in Schule u. Kirche}, 1896. *S. Bang, \emph{Katechetische Bausteine z. christocentrischen Behandlung d. I. Hauptstuecks}, ²1897. K. Just, \emph{Der abschliessende Katechismusunterricht}, 1896-07. A. Eckert, \emph{Die Formalstufen im Katechismusunterricht in Ev. Schulblatt}, 1898. F. H. Kähle, Der Kleine Katechismus L's., ¹⁹1902. G. v. Rohden, \emph{Ein Wort zur Katechismusfrage}, ³1902. *H. Brammer, \emph{Neue Bahnen f. d. Religionsunterricht}, 1900, 1903. R. Staude, \emph{Katechismusunterricht}, 1903. S. Bang, \emph{Zur Reform d. Katechismusunterrichts}, 1904. M. Reu, Explanation of Luther's Small Catechism, 1905. *A. Auffahrt, \emph{Die religiöse Frage und die Schule}, 1905. A. Reukauf, \emph{Didaktik d. ev. Religionsunterrichts in d. Volksschule}, 1906. *R. Staude, \emph{Das 8. Schuljahr}, ³1907. R. Hempel, \emph{Zum Katechismusunterricht}, ⁵1907. *H. Huebner, \emph{Wie laesst sich der Katechismusunterricht einfach, interessant und fruchtbar gestalten}, 1907. *H. Huebner, \emph{Was der kleine Catechismus fuer ein grosser Schatz ist}, 1907. A. Richter, \emph{Religionsunterricht oder nicht}, 1908. Arzt, \emph{Welche Maengel zeigt der gegenwaertige Religionsunterricht und auf welche Weise ist ihnen zu begegnen}, 1908. *W. Rein, \emph{Stimmen zur Reform des Religionsunterrichts}, 1904-1908. \emph{Die Umgestaltung des Religionsunterrichts i. d. saechsischen Volksschulen}, ed. by the Vorstand d. saechs. Lehrervereins, 1908. Pfeiffer, \emph{Der Sturmlauf gegen die Zwickauer Thesen}, 1909. Th. Franke, \emph{Der Kampf um den}
\quellseite{332}
\emph{Religionsunterricht}, 1909. *S. Bang, \emph{Luthers Kleiner Katechismus ein Kleinod der Volksschule}, 1909. *J. W. Horine, \emph{The Catechist’s Handbook}, 1909. *H. Matthes, \emph{Neue Bahnen f. d. Unterricht in Luth. Katechismus}, 1909. *F. Rendorff, \emph{Das Problem der Konfirmation und des Religionsunterrichts in d. Volksschule}, 1910. G. Rietschel, \emph{Zur Reform d. Religionsunterrichts in der Volksschule}, 1909. E. Sachsse, \emph{Zur Reform des Religionsunterrichts in der Volksschule}, 1909. F. Niebergall, \emph{Biblische Geschichte, Katechismus. Gesangbuch}, 1910. H. Matthes, \emph{Der Religionsunterricht im Dienste d. Erziehung innerhalb d. relig. Gemeinschaft}, 1910. *F. Eberhard, \emph{Der Katechismus als paedagogisches Problem}, 1912. *B. Doerries, \emph{Erklaerung d. kl. Katechismus. Ein Beitrag zur Reform d. Religionsunterrichts: Zehn Gebote}, 1912; \emph{Der Glaube}, 1912. *K. Eger, \emph{Evang. Jugendlehre}, 1912. *Th. Kaftan, \emph{Auslegung d. luth. Katechismus}, 1913. M. Reu, \emph{Die grosse Bedeutung des kleinen Katechismus Luthers}, 1913. *L. Himmels, \emph{Der Katechismus als Lebensbuch}, 1922. G. Haccius, \emph{Der Kleine Katechismus in der evang. Heidenmission}, 1923. M. Reu, \emph{A New English Translation of the Small Catechism}, 1926.

If Biblical History is properly taught it will not only introduce the child into the past, but also exhibit to him the norms of a true Christian life in the present, that is to say, it will touch upon most of the subjects expressed in the Catechism. Notwithstanding, it would be a grievous mistake to regard specific instruction in the Catechism as superfluous. On the contrary, such instruction is indispensable for the reason that the truths diffused through Biblical History and gained only by means of inference, require to be combined into groups and to be turned to account from one and the same point of view. By removing the Catechism as an independent factor, we would pour contempt upon the historical development as it has taken place under the guidance of the Holy Spirit; we would disparage the Church’s catechetical labors of almost two thousand years; we would undervalue the revelation through the Word and those parts of Scripture, particularly of the New Testament, which are not strictly historical; last but not least, we should commit a pedagogic blunder of farreaching consequences because the youth of the Church as well as the
\quellseite{333}
common people require brief, compact sentences in which the most important elements of what they may and should believe are summarized, and which may serve at any time as a medium of orientation. This much is certain: if there were no Catechism, one would have to be created.

But that does not dispose of the question whether Luther’s Small Catechism still satisfies all the requirements of today. In our time this is largely denied with great vehemence. It is gladly conceded that, in the past, it was entitled to great consideration, and that it has left many blessings in its wake. For this reason there is no objection to recognizing it as a significant historic monument of the Evangelical Church. But as medium of instruction for the young it is either decisively repudiated or largely emptied of its contents. The misgivings felt as to the employment of Luther’s Catechism in the instruction of the young are both dogmatical and pedagogic. So far as the dogmatical objections pertain to the Evangelical faith as expressed in the Catechism, we pass them by; for the opponents have gradually acquired sufficient honesty to admit that the expression of the Evangelical faith as found in the Catechism tallies perfectly with that of the New Testament. Now, since that is the case and since this faith is laid down in Luther’s Catechism in simple form, unburdened with theological problems and terms (cf. pp. 110-113), Luther’s Catechism is fully adequate; we desire, neither for our own person nor for the instruction of the adolescent youth, any other expression of Christianity than that laid down in the New Testament. But it remains to examine as to their correctness and bearings those objections against the Small Catechism which are raised from the pedagogical point of view and, likewise, those dogmatical objections which are not based on opposition against Holy Writ.

We concede from the outset that the Small Catechism has often been handled, and is still handled, in such a way as to strike all healthy pedagogic requirements in the face. When children are made to commit
\quellseite{334}
it to memory at a stage of life when they are not yet equal to the task; when the teacher assigns it for memorizing without in the very least paving the way for its appropriation; when its material is not properly outlined and the teacher fails to show his pupils the inner structure of the component parts; when the explanation fails to shed the needed light upon it from the characters of Sacred History and to connect it with the fresh, throbbing life of the present; when the catechist becomes a lecturer and possibly even lectures beyond the scope of the catechism; when he never lets the children feel the fervor of his own devotion; when by adducing features from the history of salvation he attempts to make it something that it was never designed to be; when he enlarges it into a textbook of dogmatics and chokes its own fresh life—in all such cases the hour for catechization may become one of positive torment for the children, and the bond between the Church and her young may, in consequence, be loosed rather than strengthened. But is Luther’s Catechism, this “\emph{golden gem}”, to be made responsible for these disastrous results when it is maltreated and abused? Is the Catechism to be blamed when it is accorded a treatment radically in conflict with its unique character?

Two objections are directed specifically against the \emph{First Chief Part}. It is said (a) that it is not suited for the evangelical instruction of youth because it incorporates the Mosaic Decalogue, that it leads the children back to a sub-Christian level, and that it teaches Old Testament, instead of New Testament, ethics. But the objection is not valid: the Catechism is made to shoulder a reproach deserved by its interpreters. Luther’s explanation has so thoroughly infused \emph{New Testament life} into the Old Testament commandments as to constitute this one of the greatest advantages of his Catechism. This matter has been treated at length on previous pages (101-103); only one fact is to be added here owing to the present form of the \emph{First Chief Part}. The introductory words, “I the Lord thy God,” were omitted in the editions of the Catechism published in Luther’s life time; but through the influence of the Nuremberg Sermons for Children they have been incorporated in the Catechism and retained. Thereby the Evangelical character of this part has become even more pronounced. This quality, indeed, is not brought out when one departs so far from a correct understanding of the text in question as to endeavor to educe theoretically from the words “I am” the existence and personality of God, and from the words “the Lord thy God”, his essence and character. For such a teacher, these words are merely an opportunity for an extremely unfruitful,
\quellseite{335}
because altogether unchildlike, and, at least in this connection, perfectly superfluous and fatuous treatment which, while it may possibly enable the catechist to subject his dogmatical knowledge to a wholesome review, is positively discouraging to the souls of the poor children. In truth, the words of introduction invite and attract like the open door of home. For when God gives His commandments the superscription, “\emph{I the Lord},” or “\emph{I, Yahweh, am thy God}”, He opens therewith His heart and home and offers Himself to us with all His saving power and mercy. He does not begin with the cold, severe, difficult demand: “You shall fear and love Me”; but He begins with the gift of Himself to us, with a revelation of Himself as \emph{Saviour and Father}. Not until He has given Himself to us as God and Father and opened wide the gate of home, does He open His lips to pronounce the rules and regulations of His house in the observance of which His children are to exemplify their love. With His love He intends to arouse our love so that He may reap where He first has sown. Is it not the foundation of all true godliness that God gives Himself to us? And, again, could there be a more perilous and fatal error than the opinion that one must make his own way to God by means of legal obedience? We have therefore reason to be grateful that through these words the gospel has been given precedence in the \emph{First Chief Part}; the introduction reminds us of Baptism: Yahweh has become our God and Father, in the commandments He gives His children the regulations of His house, and true children of His will gladly observe them. The authority divine is not at all weakened by this treatment, but rather made more effective. From what has been said, it is evident that the children, in the \emph{First Chief Part}, are not led into a world altogether foreign to them, as has been averred, a world for which they lack every connecting link in their own life. On the contrary, just as in their home life, the father, author of their life and provider of their bread, enacts certain regulations to which they are required to conform, so the Father in heaven with the words of the introduction opens the gate of His home, whereupon He lays down the several rules of His house. These are things quite suited for the concept sphere of the child. Nor is (b) the other objection to the \emph{First Chief Part} well grounded that the explanation of the several commandments dovetails inadequately with the life of the present. When enumerating the various virtues and faults, Luther merely means to give examples of the transgression or fulfillment of the commandments, as is clearly indicated by the “etc.” added in the original editions in the explanation of the Ninth Commandment.
\quellseite{336}
Moreover, in the Large Catechism, Luther expressly said that those features which were most commonly disregarded should receive the most careful attention. According to the Reformer’s own intention, therefore, there is no reason whatever not to supplement the examples of true obedience to the divine will given by Luther through others taken from the life of today. Personally, we are of the opinion that virtually all that is necessary is already found in the original outline. What is needed is not so much new additions as an exhaustive amplification of what is given.

When it is asserted that the \emph{Second Chief Part} (a) postulates so profound an insight into sin and a faith so strong as altogether to transcend the experience of the child, the fact is overlooked that according to Luther’s intention, it is primarily the father or the catechist who confesses himself to be a lost and condemned sinner and expresses his faith. This is irrefutably proved in the explanation of the \emph{First Article} where “wife and child, house and home” are enumerated. This is further corroborated by the practice, in Luther’s day, of the father or the catechist pronouncing, and the children and servants repeating, the words. Moreover, it must be borne in mind that the initial stages of contrition and faith can well be wrought by the Holy Spirit in the period of childhood and that, during adolescence, they sometimes reach a surprising depth (pp. 262, 268). (b) When the assertion is made that the notion of creation ex nihilo, as contained in the \emph{First Article}, is not an essential part of a saving evangelical faith, a grave error is committed. For it is the indispensable foundation for the article that God is the absolute Lord of the world and that nothing can harm us save by His permission. That this conviction is essential to the Christian faith, no one can deny. It is sometimes said (c) that the explanation of the \emph{First Article} does not correspond to the conditions in which most of our children are reared because we find enumerated “house and home, wife and children, cattle and all goods”. Now, it is quite true that Luther when he wrote this explanation had in mind the rural family; but the fact must not be overlooked that this specification is merely intended as an illustration. It is not likely that serious difficulties will arise for an experienced catechist in this connection. If (d) the fact is emphasized that the socalled intermediate causes have been completely ignored, everything being traced directly back to God, it should be remembered that the Bible deals with the subject in precisely the same way without thereby denying the operation of the intermediate causes. The judgments expressed at this place are
\quellseite{337}
judgments of faith; and faith does not deal with powers and objects belonging to the economy of nature. It is rather the very nature of faith to look beyond the material world for the invisible one lying behind it; nor is it satisfied until in everything pertaining to our natural life, the concealed hand of God has been seized which through those selfsame intermediate causes supplies food and everything else that is needed for our body and life, protects us against all danger, and guards and keeps us from all evil. It is just such emphasis upon God as the \emph{causa prima} which is bitterly needed in our day where the highly cultivated natural sciences tend to blind men against the governing hands of the invisible God. Again, (e) if one objects to the fullness of assurance with which Luther confesses that “He richly and daily provides me with all that I need for this body and life, protects me against all danger, and guards and keeps me from all evil”, we reply (1) that it is a judgment of faith which is here expressed, viz., the confidence that one is safe in the arms of God for all time to come; (2) that the accent lies on “God” in contrast to one’s own strength and wisdom; and (3) that the terms “danger” and “evil” are not to be interpreted from the narrow standpoint of our limited human reason, but, in the light of Rom. 8:28, from that of the all-seeing divine ruler of the world and of every individual life. Offense has also been taken (f) because the Apostolic Symbol silently passes over all intermediate steps between the birth of Christ and His passion and also over the prophetic ministry of Christ, and because the same omission occurs in Luther’s explanation. But let it be remembered that this is precisely the method of Paul in all his letters, and that the passion and death of Christ is in reality the climax and crowning consummation of His work.

The objection against the \emph{Third Chief Part} is that the tenor of the explanation is too heavy. We readily concede that due to the lengthy sentences and massive thoughts there is an occasional lack of lucidity and that the catechist faces peculiar difficulties on this account. If, however, he has faithfully explained the first two Chief Parts, the children are able to comprehend what Luther offers in his explanation of the Third. The last trace of difficulty will disappear if he will direct the pupils to repeat in brief sentences the respective burden of the individual petitions (cf. our Explanation of the Catechism). It is also an easy matter to add the thought of missions in the Second Petition.

When it is said that the Fourth Chief Part is too largely oriented
\quellseite{338}
upon adult baptism and that infant baptism is virtually passed by, we answer that whatever pertains to the essence of Baptism is of necessity true of infant baptism as well if the latter really is Baptism.

While it is true, finally, that Luther erred in connection with the \emph{Fifth Chief Part} in that he applied the words of institution, “Given and shed for you for the remission of sins”, to the act of distribution, and not, as is proper, to the act of Christ’s death, and then, acting on such erroneous interpretation, fixed these words as central\footnote{Toward the end of his life Luther conceded their bearing upon the death of Christ without, however, casting aside their bearing upon the Sacrament; cf. Mathesius’ xv. sermon on the life of Luther.}, we do not endorse his error; but the catechist who is convinced that the sacramental blessing, the remission of sins, is bestowed not through the body and blood of Christ, but through the Word used in connection with the Sacrament, need not consider this formula—thrice recurring—as anything but recapitulation of the words of institution whereupon he will find his difficulty vanished. While the matter is not so simple for one to whom the body and blood of Christ is intended not only to seal but also to transmit the sacramental blessing of the remission of sins, the difficulty can be surmounted nevertheless. It would be altogether wrong for such reasons to take offense at Luther’s Catechism and to doubt its availability as a medium of instruction in the present. The few occasional difficulties are more than outweighed by the numerous merits of this peerless little book.

Accordingly, if Luther’s Small Catechism is the best summary of that faith which is the life of the mature congregation, pains must be taken in catechetical instruction not to obliterate its great merits, but rather to give them full scope so that the uniqueness of the Catechism is reflected in the instruction based upon it as its text. This will be accomplished most readily when instead of going elsewhere for leading points and grafting them as well or ill as possible upon the several parts of the Catechism, all of the instruction is nothing more or less than an \emph{unfolding of the fullness of the evangelical life of faith expressed in the words of the Catechism}. It was Loehe particularly who reaffirmed this principle; Theo. Kaftan in his \emph{Auslegung des Lutherschen Katechismus} (⁶1913), and likewise
\quellseite{339}
the \emph{Synodical Catechism of the Iowa Synod} have consistently carried it out.

This principle recognized, it follows naturally that instruction in the Catechism should be restricted to the truths that constitute the center of the Christian life. That involves the exclusion of special complementary parts which, it is thought, form a necessary addition to Luther’s Catechism, as for instance, a special section relating to the Word of God or the new life of the Christian (Beatitudes, Twin command of love). It involves also the exclusion of supplementary sections taken from dogmatics and sacred history whereby the verbal limits of the Catechism are transcended. There are catechists who believe a discussion of such concepts as revelation, religion, conscience, to be necessary to an introduction to the Catechism. They believe that no adequate treatment of the Second Chief Part is possible without first having laid down a general doctrine of God, His essence and attributes, and the Holy Trinity. They cannot refrain from interweaving the whole history and doctrine of creation although now separate courses in Biblical History are found everywhere; from making an “excursus” to the subject of good and bad angels; from foisting upon their treatment of the First Article the dogmatical division “creation, preservation, and government”. They consider as a weakness of Luther’s explanation of the creed his failure to interlink between the First and the Second Article the doctrine of the fall of man, and try to improve this condition by adding an elaborate statement concerning sin and the need of redemption. In connection with the Second Article they demand a supplement treating of the sinlessness of Christ or of His whole ministry. They think that the explanation of the Second Article must be treated according to the conventional scheme of dogmatics, “Christ’s Person and Work”. It is simply incomprehensible to them how one can dispense with so important(!) a distinction as that of the prophetic, priestly,
\quellseite{340}
and royal offices of Christ (a distinction which, as a matter of fact, is to be made only with due caution, and which was not fully developed until the seventeenth century). And then, to cap the climax, they introduce an extended discussion on the Two States of Christ (although they have treated the same material before, even if in different connection!). In spite of the fact that the expression “redemption from death” is found midway between redemption from sin and redemption from the power of the devil, they cannot refrain from adding the distinction—quite appropriate, to be sure, at the right place—in regard to natural, spiritual, and eternal death. In connection with the Third Article, they endeavor to give a detailed explanation of the person of the Holy Spirit; and they believe themselves guilty of a serious omission when they fail to enumerate all the steps of the order of salvation ever constructed in dogmatics or to force everything that can be said on the subject into the few words found in the catechetical text proper. Instead of demonstrating in the light of the Lord’s Prayer itself what true prayer is, they premise an elaborate doctrine of prayer. The conception of the Sacrament with which they operate has a cast altogether dogmatical; and then they strain under obstacles self-imposed when they attempt correctly to state the heavenly blessing of Baptism, etc. That these elucidations are correct in themselves, is not the point in question. They are not in place in any kind of catechetical instruction which would preserve the uniqueness of Luther’s Catechism. That book excludes the technical ballast of dogmatical terms as much as is practicable; and instead of serving as an introduction to theology, it has no other purpose than to serve as a succinct and lucid expression of the substance of the evangelical life of faith.

If instruction in the Catechism is confined to what is central in the Christian life, it will, naturally, refrain from fusing the Five Chief Parts into a system. Such attempts
\quellseite{341}
have been made again and again. The several parts have been joined into a connected whole by such an arrangement as this: The First Part is to show us the disease, the Second the physician, the Third the means whereby the physician is to be summoned, the Fourth and the Fifth the remedies employed by the physician to effect a cure. Or by the following arrangement: The First Chief Part shows us the goal to be reached, the Second shows us faith as the way to this goal, the Third, Fourth, and Fifth show how the power to prosecute the way is to be found, namely in prayer and the Sacraments. According to the first systematic effort here described, the Ten Commandments are of moment exclusively as a mirror for sin, according to the second as the rule and norm of the new life. The attempt is made to justify both these conceptions by appealing to Luther, but without warrant (cf. pp. 111 f.); for, according to the Small Catechism at least, the First Chief Part is, indeed, primarily the norm of the true Christian life, but also the mirror for sin (see the part treating of Confession). In a type of instruction in the Catechism that takes its cue from the Small Catechism itself, the Decalogue will be used to impress the ideal Christian conduct upon the vision of the pupil with as much vividness and animation as possible. Through such a treatment, the Decalogue, apprehended in an evangelical sense, will of itself become a mirror for sin. Thus the danger of formulating a system has been obviated; for that can exist only when the Decalogue is treated either as a mirror for sin or as the norm of the new life, but not when it is treated as both. At the same time a service is rendered to the young who often have been tormented with questions relating to the connection between the one Chief Part and the other. The young and the common people require brief and definite statements as a treasure secure and rich, but not a system. — Moreover, the treatment of the
\quellseite{342}
Catechism is rather one-sided when it is based upon the notion that each Chief Part contains but one side of Christianity and excludes the others. The fact is that each Chief Part contains all of Christianity although the light and the relation in which it is viewed do not remain the same. In the First Chief Part Christianity is viewed as a divine requirement; in the Second, as the expression of personal conviction; in the Third as a divine gift to be received through prayer; in the Sacraments as salvation appropriated by visible acts; in the additions, as exemplified in the various relations of life.

Where the principle stated above is adhered to, a third wrong method stands excluded. We refer to the method which is not satisfied with merely unfolding the explanation of Luther but insists upon adding another explanation of the Catechism text proper (here and there, even one that is not at all in concord with that given by Luther). Such method has a most disturbing effect when, e.g., each step of the humiliation and exaltation of Christ is examined separately, and its meaning and importance, as distinguished from that of the preceding and following steps, is brought out; or when in connection with the Third Article its five component parts are treated separately (e.g., in \emph{Caspari’s explanation}). Such method, considered from the standpoint of pedagogics, is thoroughly vicious, connoting as it does a misapprehension of one of the greatest merits of the Catechism (cf. pp. 103 ff.). No, the text of the Catechism has been explained by Luther in accordance with Scripture, and therefore in an authoritative manner. For this reason the text confronts us with no other requirement than that a brief definition be given of difficult words; that when occasion suggests it, a few historical remarks be interspersed; and that wherever necessary, the outline of the text be disclosed, as in the Second and Third Articles, the catechist always
\quellseite{343}
taking his cue from Luther’s appended explanation (cf. Reu, \emph{Explanation of Luther’s Small Catechism}, pp. 63 f., 73-76, 88 f.). Instead of introducing the pupil to theology and dogmatics, the catechist should make every effort to bring out the very meaning of the truly central religious elements, to impress it upon the hearts of the children, and to make the application to their life. He will try to eradicate the erroneous notion as though the ‘lifeless, useless dogmas’ of the Catechism must simply be hoarded up in the memory; wherever possible he will rather exhibit the great significance, for the present life of the pupil, of the several truths of the Catechism. It will be necessary for him to draw upon Biblical History, upon secular history, upon the history of the Church and of missions; if need be, also upon the kingdom of nature, in order to gain an appropriate starting-point or illustrative material. He will find much useful material in H. Caspari, \emph{Geistliches und Weltliches}, ²1915; J. H. A. Fricke, \emph{Handbuch des Katechismusunterrichts, zugleich Buch der Beispiele}, ²1892; L. Pestalozzi, \emph{Die christliche Lehre in Beispielen}, ³1901; P. von Zychlinski, \emph{Illustrirende Aussprueche, Sentenzen und Geschichten zu Gottes Wort}, 1900; J. Besch, \emph{Aus der Lernstube des Lebens}, ²1913; and \emph{Narratives on the Catechism}, Columbus, O.; valuable sketches and narratives from the field of missions are found in Warneck, \emph{Die Mission in der Schule}, ¹⁴1912; Schaefer, \emph{Die Innere Mission in der Schule}, ⁷1912; and Koenig, \emph{Die Mission im Katechismusunterricht}, 1913. But whatever material the catechist may have selected for the purpose of bringing home to the pupil the truth in a quickening and touching manner, of impressing upon him its value, and of producing permanent impressions and promptings of the will—always the statements of the Catechism must constitute the cord upon which everything is strung, or the ultimate issue of the catechesis; nor should any element of knowledge receive consideration
\quellseite{344}
which is not an organic outgrowth of the Catechism. Only in this way will the youth of the Church be well grounded in the Catechism, katechismusfest; only in this way will it become an instrument for the orientation of their life in days to come. The teacher who is not content with letting his explanation turn upon Luther’s explanation, might well be asked why he insists at all upon an exact memorization and repeated reviews of the same.

Finally, all instruction in the Catechism should be characterized by the tendency — manifested also in Luther’s Enchiridion — to individualize instead of dealing in impersonal generalities (compare “we should” in the First Chief Part; “I” and “me” in the Second Chief Part; “we,” “us,” “our” in the Third Chief Part). The teacher should be a \emph{Seelsorger} rather than a dogmatician: he should try to arouse not only the intellectual faculties, but also the feelings and the will, in short the whole man. To summarize: instruction in the Catechism should possess all those marks that constitute the superlative excellence of the Small Catechism. The more it preserves and reflects its unique character, the more likely it will become truly evangelical and pedagogically correct. And it will reflect its uniqueness best when the instruction from the beginning has no other aim than to be an unfolding of the fullness of evangelical faith, as it is expressed in the words of Luther’s Catechism.

We take the liberty of adding a few hints as to a correct understanding of the Small Catechism. While we must confine ourselves to essentials at this place, we refer for details to M. Reu, \emph{Explanation of Luther’s Small Catechism} (larger edition, Chicago, ¹¹1926; smaller edition, ibid., ³1926); M. Reu, \emph{A New English Translation of the Small Catechism}, 1926, to Vol. xxx of the \emph{Weimar Edition of Luther’s Works} Weimar. Boehlau, 1910, upon which O. Albrecht has bestowed great care; A. Ebeling, \emph{Historisch-kritische Ausgabe von Luthers Kl. Katechismus}, Hannover, 1901; A. Hardeland, \emph{Luthers Katechismusgedanken in ihrer Entwickelung bis zum Jahre 1529}, Guetersloh, 1913; J. Meyer, \emph{Luthers Grosser Katechismus, mit Kennzeichnung seiner Predigtgrundlagen},
\quellseite{345}
Leipzig, 1914; and esp. to Th. Kaftan, \emph{Auslegung des Lutherischen Katechismus}, Schleswig, 1913, the best commentary to Luther’s “golden gem”. In the last named work all important books on the subject are noted. Parts of Kaftan’s book are extant also in the English language, in the form of \emph{The Catechist’s Handbook}, by J. W. Horine, Philadelphia, 1909. While my own views in the premises are largely in agreement with those of Kaftan, justice to our subject requires the statement that, as my notes show, my views, on the whole, were formed years before I became acquainted with the book of Kaftan. The impulse for my conception of the Catechism in general came from the foreword of W. Loehe’s Catechism of 1845.—Abundant material relative to the history of the exposition of the Catechism is found in M. Reu, \emph{Quellen zur Geschichte des kirchlichen Unterrichts im evangelischen Deutschland zwischen 1530 und 1600}, Guetersloh, 1904 ff.

\emph{The introduction should be brief.} Following Brenz’s example (a similar procedure is found already in \emph{Fundamentum aeternum felicitatis seu summa religionis Christianae} which appeared in print in 1498 after some years of manuscript circulation) and that of many expositors of Luther’s Catechism, notably Bischoff (1599; cf. Reu i, 2, pp. 222 f.), the teacher will use as a starting point the sacred act by which the catechumens have become children of God, i.e., Baptism; touch upon Holy Scripture which bears witness to that fact; and then pass on to the Catechism which, in five phases, is designed to teach them in what the faith and the life of God’s true children consist.

\section{The First Chief Part}

“The Ten Commandments,” Luther says at one place, “reveal to us what we have been, what we are, and what we should be.” His mode of presenting the subject indicates that he wishes the last point to be stressed: always he begins his explanation with the words “we should”. His aim is to hold before the young the ideal Christian life and to show them, in clear outline, what they should do and what they should avoid. But are the commandments of Sinai adapted to that purpose? Are they not an integral part of the Law which has found its end in Jesus Christ, in fact, are they not that very Law summarized? Quite true, in their Old Testament form they are obsolete as is clearly seen when we compare the Fourth and the Third Commandments with Col. 2: 16 f. and Gal. 4: 10 f. Nor did Luther incorporate them in the Catechism in their Old Testament form; still less did he give them an Old Testament interpretation (cf. pp. 101-103). But he gathered the divine thoughts of eternal value that were contained in the Old Testament form and these he explained in the sense
\quellseite{346}
and spirit of the New Testament. This, and this alone, is the acceptation of the Ten Commandments in the present: so understood, they are gratefully welcomed by the Christian, eager to do the will of God of which he is often ignorant, and gladly accepted as the rule and norm of his life. The children will undoubtedly observe the difference between the biblical wording of the Decalogue (as found also in their Biblical History) and that in the Catechism; this gives the catechist a fine opportunity to spread healthy evangelical conceptions of Holy Scripture. In fact, it is his duty to make this matter plain at the outset because the Reformed never cease accusing Luther of mutilating the text of the Bible in this respect. Afterward, in connection with the Third Commandment, a similar opportunity presents itself.

The Commandments are divided into three parts: (1) The Introduction; (2) The Commandments proper; and (3) The Conclusion. \emph{The Introduction}—omitted by Luther who followed the traditional text—was rightly restored almost everywhere (cf. p. 312 f.). Luther himself took note of it not only in his sermons on Exodus (1525) and Deuteronomy (1529), but also in his catechetical sermons of 1528. The exodus from Egypt referred to in Ex. 20 and Deut. 5 was to Luther the “sign by which the Jews were to lay hold of God”; what the exodus was for the Jew, redemption by Jesus Christ is for the Christian; at one place there is even a reference to Baptism (Weimar Ed. xvi, p. 425). “For this reason it is not necessary that they [the Christians] should fear and love God under the name of deliverer from Egypt: He has received another name now, namely Christ who, true God, has died for us” (Weimar Ed. xxviii, p. 605). In his annotations on the Decalogue of 1530, he therefore is able to say in regard to the words of introduction: \emph{Promissio omnium promissionum, fons et omnis religionis et sapientiae caput, evangelium Christum promissum complectens} (Weimar Ed. xxx, 2, p. 358). Accordingly, the catechist is following in Luther’s steps when he uses the Introduction as an occasion to remind his pupils of their redemption by Jesus Christ or of the baptismal sacrament by which the former is appropriated, and says to his class of Christian children: “These commandments are given you by Him, who, through Baptism, has become your God and Father. But if He has granted you this wonderful blessing that He has become your God and Father, it is His right to give you rules and regulations for your conduct, and He can justly expect you to live up to them in grateful love.” Sometimes a passage from Luther’s Table Talk (Erlangen Ed. lviii, p. 206 f.) is adduced as disproving the statements made above;
\quellseite{347}
but the objection is not well taken: Luther does not argue at all against the retention of the Introduction, but against its inclusion among the commandments (cf. Reu i 1, p. 451, 758), because to him such a procedure would, naturally, be a perversion of a “promise” into a commandment, a suppression of the comforting gospel which stands at the head of the commandments (cf. p. 313).

We are ignorant of the arrangement of the Commandments upon the tables of Moses. This should prompt the catechist to pay no attention to this point when he discusses the question as to the division of the commandments; still less will he exaggerate the division 1-3 and 4-10 into a shibboleth of the Lutheran Church. While there is no adequate reason to depart from this division, no attempt should be made to maintain, with an unwarranted reference to Matt. 22:36-40, that the first three commandments treat of that which we owe to God while the last seven treat of that which we owe to our neighbor; for the child will not number his parents among the “neighbors” (note that Luther does not introduce the term “neighbor” until he comes to the Fifth Commandment). For this reason Melanchthon and many interpreters of Luther’s Catechism made a distinction between what we owe to God and what we owe to men. While this distinction has a better warrant, it is still misleading. A mere glance at Luther’s explanation is sufficient to assure us that our duty toward God is the subject under consideration also in the last seven commandments; the difference is this that in the first group the obligation imposed is direct, while in the second group it is indirect. The first table shows what we owe to God himself, and the second table what we owe to our fellow-man for the sake of God. A reference to Matt. 25:40 will at once make the subject clear to the child: he will begin to comprehend the important truth that all he does bears not only on men, but primarily on God (Psalm 51:4).

He who has come to a proper understanding of the Introduction will find the First Commandment the necessary logical sequel. If God through Baptism has become our God, it is His right to require that we should have no other gods before Him, but have Him alone as our God. The question when that is the case that we have God as our God, is usually answered by Luther, quite in accordance with the New Testament, in that he simply points to faith. E.g., he writes in the explanation of the Decalogue as found in his Sermon on Good Works (1520): “Since I alone am God, thou shalt place all thy confidence, trust and faith in Me alone, and on no one else. For that is not to have a god, if you call him God only with your lips, or worship him with
\quellseite{348}
the knees or bodily gestures; but if you trust Him with the heart, and look to Him for all good, grace and favor, whether in works or suffering, in life or death, in joy or sorrow . . . And this faith, faithfulness, confidence deep in the heart, is the true fulfilling of the First Commandment; without this there is no other work that is able to satisfy this Commandment. And as this Commandment is the very first, highest, and best, from which all the others proceed, in which they exist, and by which they are directed and measured, so also its work, that is, the faith or confidence in God’s favor at all times, is the very first, highest, and best, from which all others must proceed, exist, remain, be directed, and measured. Compared with this, other works are just as if the other Commandments were without the First, and there were no God” (Weim. Ed. vi, p. 209). Similarly he says in his sermons on Exodus (1525): “All the First Commandment requires is genuine faith and trust in God; it does not enjoin anything external. But no one can fulfill it unless the Holy Spirit work it in his heart” (Weim. Ed. xvi, p. 445; cf. p. 464 where he repeats the exposition of 1520 almost literally). In the same strain he expresses himself once more in his Catechetical Sermons of 1528: “\emph{Ergo intentio huius praecepti est dass es will gebieten einen rechten Glauben}” (W. E. xxx 1, p. 28), and in the Large Catechism: “Therefore it is the intent of this commandment to require true faith and trust of the heart which settles upon the only true God, and clings to Him alone (W. E. xxx 1, p. 133; Concordia Triglotta, p. 581, § 4). Just because it requires faith, he aptly calls the First Commandment “the chief part of all our Christianity” (in his sermons on Deuteronomy of 1529; W. E. xxviii, p. 601), and properly describes the Church as “the number of those who trust in nothing save the mercy of God alone and understand the First Commandment” (W. E. xxviii, p. 580). In the sketch De Loco Iustificationis he even writes: “pueri et infantes confirmant suo catechismo solam fidem absque operibus iustificare . . . Primum praeceptum est promissio quod velit esse Deus. Et fidem exigit ante omnia opera quae sequentibus praeceptis exiguntur. At fide secundum primum praeceptum habita filii Dei sumus, remissis iam peccatis, ipsa fide iusti” (W. E. xx p. 663). Now, when Luther in his Small Catechism in answering the question when we regard God as our God, points to something threefold, it is not likely that we are to look for anything else in his answer than a development of the conception of faith from the pedagogical standpoint. And indeed, he who believes in God as his God and Father is bound to fear, love, and trust in Him above all things. To trust in
\quellseite{349}
him; for, to trust in God as One who in Jesus Christ has become our Father and who now so governs all things that nothing can harm us against His will,—that is the very central part of our faith. To love Him; for, to love means to surrender to God; and this, too, is involved in faith: at first (logically) faith can, indeed, do no more than receive God and His precious gift, but on the strength of these gifts divine it subsequently rises to personal self-surrender. To fear Him; for how should faith maintain itself without feeling awe of that God who has shown Himself so great and sublime in all His deeds? The assertion has indeed been made, in view of the conclusion of the Ten Commandments, that Luther did not have in mind the awe of God natural to a believer, but the fear of God's wrath and punishment. Now it is quite true that Luther when he has occasion to explain the expression “jealous God” does stress this aspect; and that the fear of God's wrath is for the Christian as long as he is garbed in the flesh, and just because of that fact, an additional motive for obedience, has been conceded also by Luther. We read, e.g., in a letter dated Oct. 27, 1527: “\emph{Timeant poenam et infernum omnes impii, Deus aderit suis, ut simul timeant Deum cum poena. Neque fieri potest, ut sine timore poenae sit timor Dei in hac vita, sicut nec spiritus sine carne, etiamsi timor poenae sit inutilis sine timore Dei}”. (Enders vi, 109). But this does not at all justify the view that Luther, in his explanation of the First Commandment and in repeating the words, “We should fear and love God”, in his explanation of the other commandments, had this idea of fear in mind. This is impossible because his explanation of the First Commandment is always under the influence of the superscription which is to him sheer gospel and promise, and because he here unfolds the conception of faith. In the third series of Catechetical Sermons, of 1528, in which the First Commandment, even more emphatically than usually, is summarized as timor and fides, timere, according to one manuscript, is also found in juxtaposition with colere and venerari (Weim. Ed. xxx 1, p. 59, lines 32 and 34). In other words, what Luther desires to bring out forcefully is this: God is so great and august, so sublime and wonderful that, filled with holy awe, the Christian can only marvel and adore in His presence and that he must tremble at the very thought of ignoring His will or offending His love by ingratitude or sinful conduct. No one can doubt that this idea of God's sublimity, majesty, omniscience, omnipresence is a most powerful motive for a righteous life, and at the same time a motive most appropriate for him who, in faith, has received a new vision of
\quellseite{350}
God’s greatness and wondrousness. It is the conclusion of the commandments that supplies the thought of divine wrath and punishment; it reminds the Christian: If you, in your carnal mind, are so thoughtless and foolish that the idea of God’s wondrous greatness does not move you to sanctify your life, remember that this great and august God can and will prove His greatness also by punishing, in terrible wrath, all those who despise His commandments and frivolously or maliciously ignore His will!” According to Luther, then, the fear of God’s wrath and punishment is added—but only added—to the more important factor (awe of God’s greatness as inherent in faith and inseparable from it). Ten times he speaks of the latter, only once of the former; just in proportion as we can dispense with the former in our life and the latter becomes the dominant factor in it, we shall become more and more what we should be. The catechist who has recognized this fact will, therefore, with a reference to a passage like Gen. 17:1 or Tob. 4:6, explain this conception to his children in the following manner: “To fear God, means always to have before our eyes the great and majestic God, the almighty and omniscient One, so that we fear to sin against Him.” He had better not introduce the dogmatic distinction between filial and servile fear since the fear taught in the First Commandment, thus defined, does not quite tally with filial fear.—Likewise, since he expounds the First Commandment to Christian children, he will refrain from any reference to pagan idolatry; still less will he accept the more than questionable distinction between ‘coarse’ and ‘refined’ idolatry, but, with Luther (cf. the Large Catechism), confine himself to that idolatry which fears, loves, and trusts something else more than it does God, and which, for a Christian, to whom God has revealed Himself as his God and Father, is coarser and more wicked than any species of pagan idolatry. Considered pedagogically, a discussion of heathen idolatry at this point is a deflection involving the danger that the children regard Christian idolatry as the lesser offense. The vigorous underscoring of the “above all things” will so completely engage the teacher’s attention that no time is left for anything else. The best exposition of the First Commandment of the Lutheran Catechism is found in the Nuremberg Sermons for Children (in Reu i 1, p. 465 ff.), except that at one place a reference to God’s wrath has erroneously crept into the explanation of the fear of God.

In explaining the remaining commandments, the catechist will take pains regularly to stress the “we should fear and love God” of Luther lest he inculcate Old Testament instead of New Testament,
\quellseite{351}
medieval instead of evangelical and reformatory, ethics. Eventually he will accept no other answer to the question, What does God require in this or that commandment? than this: We should fear and love God. In this way he will effectively drive home both the unity of Christian ethics and the fundamental truth that no work pleases God unless it comes from a heart that fears and loves Him (p. 102 f.). Says Luther in the Large Catechism: “Thus the First Commandment is to shine and impart its splendor to all the others. Therefore must this declaration run through all the commandments, like a hoop in a wreath, joining the end to the beginning and holding them all together; that it be continually repeated and not forgotten, etc.” (Weim. Ed. xxx 1, p. 181; ‘Triglotta, p. 675, § 326-329). In his Sermon on Good Works (1520) Luther calls the faith required in the First Commandment the Christian’s “\emph{health}” and says: “. . . a man must first have \emph{health}, which will work all the works of all the members of the body. So faith also must be in all works the master-workman and captain, or they are nothing at all” (Weim. Ed. vi, p. 213). The other commandments, from the second to the tenth, merely bring out in what manner the fear and love of God is to show itself in thought, word, and deed: ‘whereunto it is to “\emph{urge and impel us}”’ (Weim. Ed. xxx 1, p. 181, line 6; cf. Reu, Explanation of Luther’s Catechism, p. 33). That Luther, in constantly repeating the First Commandment in his explanation of the others, does not likewise repeat the word “trust” but confines himself to the words “fear and love”, is probably explained by the fact that both the positive and negative elements in the several commandments correspond better to the words “fear and love” than to the word “trust”; especially since Luther used it in a unique sense—trusting God for leading us in the best way. This view is borne out by the Large Catechism; for there we read: “Whatever you lack of good things, expect it of Me, and look to Me for it, and whenever you suffer misfortune and distress, creep and cling to Me. I, yes I, will give you enough and help you out of every need; only let not your heart cleave to or rest in any other . . . . Look to Me for everything, and regard Me as the one who will help you and pour out upon you richly all good things . . . We are to trust in God alone, and look to Him and expect from Him naught but good, as from one who gives us body, life, food, drink, nourishment, \emph{health}, protection, peace, and all necessaries of both temporal and eternal things. He also preserves us from misfortune, and if any evil befall us, delivers and rescues us, so that it is God alone (as has been sufficiently said) from whom we receive all good, and by whom we are delivered
\quellseite{352}
from all evil” (Weim. Ed. xxx 1, p. 133 f.; Tr. p. 581 ff.). So understood, trust can not at once be recognized by the child as a motive for the fulfilment of the Law. The child can grasp this idea much better in the context in which Luther has it in the conclusion: If now we keep the commandments, and still perceive but little of the promised ‘grace and every blessing’, we should not be dismayed at this unexpected experience, but cling to our conviction and trust to God that, in due time, he shall acknowledge our actions as done for Him and crown our efforts with a happy issue.—That Luther meant to correlate the fear of God rather with the negative features of the commandments, and the love of God with the positive ones, is a view doubtless warranted by the facts; but that, after all, not much weight is to be attached to this distinction, is evident from the Sixth Commandment which, like the others begins with “We should fear and love God” and yet contains only positive statements.

The conjunction “dass” in the explanation of Commandments Two to Ten has not received the same acceptation everywhere. The first Latin translation (in the Enchiridion piorum precationum, 1529) and, in part, also the third (by Justus Jonas, 1539), change the whole construction and offer instead, e.g.: \emph{Debemus Deum timere et amare, divinum eius verbum non contemnere aut negligere, sed illud potius magnificare et venerari, libenter tum audire ab aliis, tum alios docere}. The second Latin edition (by Sauromannus, 1529) and the fourth (by Hiob Magdeburg, 1560) construe the conjunction as final, rendering it by ut, resp. ne (e.g. \emph{Debemus Deum timere et diligere ne divinos sermones, eius verbum contemnamus, sed ut sanctum repetemus} etc.; or, \emph{Timendus et amandus Deus est ne sacras conciones et ipsius verbum negligamus, sed ut illud sancte veneremur} etc.). The Greek translations have δεῖ with the infinitive (Hiob Magdeburg: Δεῖ ἐν φόβῳ καὶ ἀγάπῃ τοῦ θεοῦ ἡμᾶς μὴ καταφρονεῖν κηρύγματος τοῦ θείου λόγου κτλ.) or they have ὥστε or ἵνα. Menius (1532) and others think that it introduces the manner in which we are to fear and love God (to make the matter plainer they add also; e.g., \emph{Wir sollen Gott also fuerchten und lieben, dass wir die Predigt und sein Wort} etc.). With this corresponds the former English translation: We should so fear and love God as not to despise His word, etc. It appears that Luther looked upon the conjunction in question as introducing the contemplated result, for he writes in the Large Catechism: “\emph{In like manner shall such fear, love, and trust urge and impel us not to despise his Word, but to gladly hear, learn and honor it, and esteem it holy}.” This quotation warrants the conclusion
\quellseite{353}
drawn by Albrecht that the clauses beginning with \emph{sondern} are to be construed as governed by \emph{dass}, and not as independent clauses. The explanation of the Sixth Commandment does not contradict this view; and the explanation of the Eighth Commandment which repeats \emph{sollen} after \emph{sondern} represents an exception.

The history of Luther’s Catechism in the English language evidences the difficulty to reconcile the meaning couched in Luther’s original with the genius of the English language. The following quotations may serve as illustration: We should fear and love God that we may not despise preaching and His Word; We should so fear and love God as not to despise preaching and His Word; We should fear and love God and for His sake not despise. . .; We should fear and love God and not despise preaching and His Word. From the standpoint of diction the last named translation deserves preference. Although syntactically the two clauses of the sentence are coordinate, a brief and simple explanation is sufficient to bring out the original meaning of Luther’s.

As regards the text of the Second Commandment, the threat connected with it in the Old Testament was not accepted by Luther nor is it found in the editions of the Catechism which appeared during his life-time. The catechist has so much less reason for wishing to correct Luther as the demand for the addition in this country is largely based upon the erroneous conception, that we are tied to the words of the Old Testament (cf. p. 323 f.). At all events, the catechist will take care not to interpret Luther’s “fear” in the light of the divinely threatened penalty. By that construction he would merely vitiate what he has taught concerning the meaning of this word in the First Commandment and what he will teach again in Commandments Three to Ten. Moreover, he would anticipate the comments to be made, according to Luther, upon the Conclusion of the Commandments. We may note also that Luther did not use the word \emph{missbrauchen}, misuse, as in his German Bible, but \emph{unnuetzlich fuehren}, take in vain. — The term “swear” has been misunderstood for so long a time that the correct view of it cannot gain ground without encountering considerable obstacles. As commonly understood, it means swearing in court. If this meaning is accepted, the word must be presumed either to forbid the taking of an oath altogether or the taking of a false oath. Luther, however, knows nothing of the first interpretation, declaring in his Large Catechism that taking an oath is “a very good work whereby God is glorified.” The second interpretation would be permissible; for the confessional
\quellseite{354}
booklets of the Middle Ages often contain the phrase “not to swear” as an equivalent of “not to swear falsely.” But the term “swear” occurs in these booklets also in the sense of thoughtless swearing in every-day life. Which of the two meanings Luther had in mind, will best be decided by his catechetical sermons of 1528 and the Large Catechism, drawn from these. Now, there is no doubt whatever that, in the Large Catechism particularly, all emphasis is placed upon “lying and deceiving by the name of God” and that false swearing, or perjury, is expressly mentioned as one form of lying by the name of God. It is also quite significant that the Latin translation of that passage in the Large Catechism which suggests most strongly the concise form of the Small Catechism (cf. Mueller, p. 397, § 62; Tr. p. 597), does not mention \emph{iurare} at all, but reads: \emph{illud blasphemando, execrando, maledicendo, incantando, ignominiose usurpamus}. The obvious conclusion is that Luther, in this connection, had in mind the thoughtless swearing in every-day life, a conclusion drawn also by Kantz, 1542; Meckhart, 1553; Trotzendorf, about 1556; the Joachimstal Catechism, 1574; Opitius, 1588; and others. Or, “swearing” might possibly be linked still more closely with “cursing” as a virtual synonym of that word so that the whole phrase would be practically equivalent to “using profane language.” The catechist who is guided by Luther’s intentions will treat of perjury under the head of lying and deceiving by the name of God; of the just, divinely warranted oath under the head of calling upon the name of God (cf. Large Catechism: Trotzendorf; Dressler, 1561; Hubernus, 1544; and others). He will thus avoid the pedagogic error of treating something that is commanded and classed as a good work, under the head of things forbidden (cf. Reu, Explanation etc., p. 36 f.)

In explaining the \emph{Third Commandment} the catechist will take special pains to set forth clearly the evangelical conception of the Sunday and to make plain that the Sunday is an institution arranged by the Church (cf. pp. 101 f.; 323 f.). For this reason he will not feel inclined to surrender the form prescribed by Luther, “Thou shalt sanctify the holy day,” and possibly to exchange it for “Remember the Sabbath day, to keep it holy.” He is aware, of course, that Luther did not originate the version which is found in his Catechism; and he also knows that even this version does not entirely correspond to the evangelical conception of the Sunday: on basis of Col. 2:16 we Evangelicals recognize no divine distinction between one day of the week and another. He also realizes that the adoption of the
\quellseite{355}
conventional English form, “Remember the Sabbath day, to keep it holy,” does not necessarily imply a denial of evangelical liberty; it is found, e.g., in all the Latin translations of the sixteenth century (also in the Book of Concord of 1580), and, through the influence of the Nuremberg Sermons for Children, also in several German reprints and enlarged editions (e.g., even in the \emph{Nuremberg Textbook for Children}, composed in 1628 and in use throughout the seventeenth century). Still, because it departs from the wording of the Old Testament, Luther’s version is better adapted to impress upon the minds of Lutherans their freedom from the literal phraseology of the Old Testament Decalogue. Formerly, perhaps, little depended upon the wording as such because Calvin as well as Luther rightly interpreted the words in question; but at this day it is well to pay more attention to a correct formulation. In the first place, the sections that went over to the Reformed Church (e.g., Anhalt, Lower Hesse, the Palatinate), felt called upon to “purge” or “reform” Luther’s explanation of the Decalogue by adding the words of the Old Testament, totally ignoring the New Testament freedom from the Jewish Decalogue (Reu i 1, p. 217; i 2, 2, p. 377; i 2, p. 278 ff.). In the second place, it is a fundamentally false attitude toward the Old Testament and the Law that is evidenced by the continued demands of the Reformed for a literal rendering of the Mosaic Decalogue. The catechist who would interpret the \emph{Third Commandment} in a really evangelical manner, will best equip himself for the task by a repeated perusal of the pertinent section in the Large Catechism. After emphasizing the correct attitude toward the Word and its proclamation and enlarging upon the blessing and importance of these, the catechist will find that his pupils do not despise the Church institution of Sunday; they will rather gladly meet together Sunday after Sunday, pray together, hear God’s Word, partake of the Holy Sacrament, and by that practise, manifest and edify themselves as a Christian congregation, and they become willing to lay aside all temporal labor that might prevent them from appearing at the common service, and quietly to meditate upon what has been heard and learned. Thus is made possible a Sunday rest that, quite apart from a purely human need of rest and the social aspect of the question, bears the impress of the gospel.—When the origin of Sunday as an institution of the Church is under discussion, it is customary to mention, alongside the fact of Sunday being the day of Christ’s resurrection, that Pentecost, too, fell upon that day (so also in our \emph{Explanation of the Catechism}, larger edition).
\quellseite{356}
But since this circumstance was no factor at all in the choice of just Sunday as the day of rest, that argument is irrelevant.

Although the several statements of the Fourth Commandment, grammatically, have reference to superiors as well as parents, they were chosen by Luther with a view to the latter rather than the former, a fact that the catechist dare not forget. Nevertheless, while enjoining upon his pupils the duties toward their parents, it will be incumbent upon him to flash the light of Scripture also upon the relation toward superiors, which is done best by referring to the Table of duties (\emph{Haustafel}). The duty of the catechist becomes so much more imperative as it is largely thought that, in this respect, the injunctions of the Bible have become more or less antiquated. Yet, however radical the change of social conditions, the duties of deference to authority and of faithful service stand unaffected. The idea—so classically expressed in the Large Catechism—that the parents are God’s own representatives and that all authority of superiors in school and state is grounded in the parental office thus defined, is of abiding moment. The better the catechist succeeds in arousing the conscience in this respect, the more his teaching will contribute toward preserving unimpaired the foundations of a healthy public life.

Rightly Kaftan points out that the \emph{Fifth Commandment} refers to the exhibition of the fear and love of God merely in the sphere of the neighbor’s bodily life; a teacher influenced by Luther’s example should not treat here of the commission of spiritual harm and murder. It is true that \emph{Leid tun} is not equivalent to German \emph{schaden} or our “harm” and that it actually denotes suffering in the sphere of the soul; nevertheless Luther refers in this connection to the neighbor’s body (we should fear and love God, and not hurt nor harm our neighbor in his body). Accordingly the subject under consideration is an “injury dealt to the body through harm done to the soul.” The right explanation would therefore be: “To embitter and shorten the life of the neighbor by envious looks, by spiteful words, or by hostile acts, as Esau and the sons of Jacob did.” When, however, Kaftan proposes to treat also of suicide under this head, he goes beyond Luther; quite likely also beyond the original intent of the Mosaic Law since in all the commandments, from the Fourth to the Tenth, the relation to our fellowmen is the dominant factor even though the term “neighbor” does not appear until the Eighth Commandment. Suicide had best be treated in connection with the Sixth Petition.
\quellseite{357}
In the \emph{Sixth Commandment}, the words, \emph{\emph{ein jeglicher} sein \emph{Gemahl} lieben und ehren}, are hardly to be construed as an additional main clause. The words are governed by the conjunction \emph{dass}, and \emph{ein jeglicher} refers back to the subject \emph{wir}: “and that each of us love and honor his spouse.” \emph{Gemahl}, “spouse”, as usual with Luther, is neuter although he probably had woman in mind (cf. \emph{ein jeglicher}). Present conditions and the duty to explain the text compel the catechist to discuss briefly the nature of marriage and to place strong emphasis upon its divine origin and its indissolubility. While the duties of husband and wife need merely be touched upon, the phrase “live chaste and pure etc.” requires a thorough as well as tactful treatment. Here Luther’s restriction to the positive element, of transcendent value pedagogically, should be the pattern for every teacher, especially when boys and girls are taught together. That which is forbidden “should be allowed to show itself only as the shadow, as it were, of that enjoined.” If there should be an occasion to enter into details, it should be done in private conversation with the individual. For the practical treatment of the subject much can be learned from Ziethe, v. Rohden, and Heydt, \emph{Die unterrichtliche Behandlung des 6. Gebots in der Schule. Drei gekrönte Preisschriften}, Berlin, 131894.

The \emph{Seventh Commandment} is to show us how the fear and love of God is to reveal itself in relation to the neighbor’s property. It is impossible to unfold Luther’s text without demonstrating the divine sanction of property and the unequal distribution of wealth as divinely intended. The teacher, likewise, must cast the light of truth upon the various forms of unfaithfulness or commercial dishonesty which his pupils may observe. The salient words serving this latter purpose are “nor get it by false wares or false dealing” (falsis mercibus aut impostura, in Sauermann’s translation; corruptis mercibus aut fraude aliqua, in the first; falsis mercibus, ullis malis artibus proximum decipiamus, in the third; fucosis mercibus aliisque actionibus fraudulentis nobis ea vendicemus, in the fourth Latin translation). “This applies to the indirect ways of stealing as it is done in all sorts of transactions, esp. in buying and selling, when the wares bought and sold are false, i.e., adulterated or counterfeit; or when the dealing is false i.e., when deception is practiced through misrepresentation, short weights and measures, etc., or when debts are not paid, either to the individual (loans, credit, keeping back wages), or to the company or state (fares, taxes, customs duties); or when trustees are faithless (embezzlement, breach of trust), or when advantage is taken of need (usury), or when
\quellseite{358}
no equivalent is given (gambling), or when false pretense is made (begging), and the like’ (Kaftan-Horine, p. 48); also practices of labor unions and trusts, insofar as they result in material loss to the neighbor, are entitled to attention should the occasion for a discussion of them arise (cf. our Explanation, p. 50). The catechist, however, should always bear in mind that he is to judge of all these things neither from the legal nor from the economical, but from the religious standpoint so that, also here, not natural, but Christian, Evangelical ethics are inculcated. Here, too, Luther has placed in the foreground the fear and love of God as the source of power and the motive of action. He who fears and loves God will recognize the divine order, also in regard to property, and beware of displacing the divinely imposed bounds and barriers; for he knows that he has in his heavenly Father the treasure of treasures and a helper in poverty and distress, and that earthly property has never permanently satisfied the heart. Only the lover of God can help his neighbor to improve and protect his property and business: he has learned from the heavenly Father that it is more blessed to give than to receive. The term \emph{Nahrung} here does not signify food, or provisions, as in the First Article, but is used in a wider sense, signifying the means of livelihood, the source of income, trade, business (cf. Gen. 46:33 f. in the German version), and, in general, income, earnings, property (\emph{conditio}, second Latin translation). The positive part of the explanation does not contain three statements: help, improve, protect; but only one: we should help to improve and protect our neighbor’s property and business (Sauromannus: \emph{sed demus operam ut illius opes conserventur et eius conditio melior reddatur}; H. Magdeburg: \emph{sed ad augendas facultates eius et conservandam rem familiarem prompti simus}).

According to the \emph{Eighth Commandment} we should exhibit our fear and love of God by exemplifying both truth and love in all that we say about our neighbor. Well it is said by Th. Harnack in his explanation of the Catechism (in \emph{Katechetik} ii, Erlangen, 1882): “Truth and love must always be found together. Truth protects love from the fault of good-natured indulgence and weakness which arises from indifference to the truth of the matter. Love, conversely, protects truth from the severity of injustice arising from indifference to the person. For this reason truth is the light and salt of love, and love is the fire and warmth of truth”; compare in connection with this and the other Commandments, in addition to Luther’s Large Catechism, the Nuremberg Sermons for Children (Reu, i 1, pp. 462-564), a work always
\quellseite{359}
instructive, though occasionally going beyond Luther. \emph{Beluegen}, “belie”, according to Luther’s usage, does not by any means signify to lie to anyone, but to lie about him, to spread lies about him; “falsely” is added not only to add force to the conception of falsehood conveyed by the verb “belie”, but also to disclose the motive back of such falsehood,—to speak as a deceiver, purposely and craftily to utter a falsehood in order to harm the neighbor. “To betray” means to imperil some one through disclosure of what he desires to keep secret; it is equivalent to Latin \emph{prodere}. This being its meaning, the example of Judas, frequently adduced, is hardly suitable. The motive of Judas was quite different; nor did Jesus desire the place where He kept Himself to remain concealed. \emph{Afterreden oder boesen Leumund machen}, Luther writes and thus connects these two verbs as closely \emph{luegen and truegen} in the Second, or \emph{abdringen and abwendig machen} in the Tenth Commandment. The slandering (\emph{boesen Leumund machen}) is done by means of backbiting (afterreden); by speaking evil about the neighbor behind his back, the coward brings him into bad repute. Hiob Magdeburg combines both notions into one in his splendid translation: “\emph{Timendus et amandus Deus est ne de proximo dicamus mendacium aut invite illum deferamus neve existimationem eius violemus lingua futili, sed ut illum excusemus, bene de illo loquamur et in meliorem partem interpretemur omnia}”. Sauermann retains also the adverb “falsely”: “\emph{ne proximum falsis mendaciis involvamus}.” The German word \emph{Leumund} is by no means an abbreviation for der Leute Mund; its Old High German equivalent is \emph{hliumunt}, from the same root as the Gothic \emph{hliuma}—ear, hearing; so \emph{Leumund} is—\emph{fama}, Engl. fame, and \emph{boesen Leumund machen}—defame. The Joachimstal Catechism (1574; probably based on Mathesius) says: “ ‘Falsely belie’ means to raise false accusations against one before the authorities and other people, to impute to others the abuse of which oneself is guilty, like Gehazi. ‘To betray’ means to divulge everything, to peddle around what we have heard, and to be mongers of scandal; likewise, without necessity to reveal secrets, as the Ziphites did in the case of David. ‘To Backbite’, to make known behind their backs what people do and say; to censure and put the worst construction upon it, as did Doeg. ‘To slander’ means by cutting remarks to hack his honor to pieces as meat is hacked to pieces upon the butcher’s table with the result that thus evil suspicion fastens to him and he is robbed of his good name as was done by Absalom in the case of his father David.” Tetelbach, 1568: “ ‘To excuse’ means to defend the neighbor against evil mouths and false tongues; ‘to speak well
\quellseite{360}
of him’, to evince a kindly disposition toward him wherever he is the subject of conversation; ‘to put the best construction on everything’, to cover up the frailty and fall of the neighbor and to excuse and explain them as best we can to those who speak ill of them.” Siber, 1575: “ut ad calumniatores excusemus, obtrectantes refutemus, et commode sentientes de illo in melius omnia, interpretando accipiamus.” Tetelbach, who understood Luther, takes pain, in connection with every commandment, to point to the root of true fulfilment—the fear and love of God; here he asks the question, Who transgresses this commandment?, and his answer is: Those who do not fear God, but follow the devil, the father of all liars and slanderers. Afterward he asks the other question, Who fulfills this commandment? Answer: Those who love God; for the love of God works in us truthfulness, sincerity, and a helpful tongue (Reu i, 1, p. 681).

As to the Ninth and the Tenth Commandments, Luther had previously been in the habit of treating them as one, even in the Large Catechism. If now, in the Small Catechism, he treats them separately, the reason is, of course, not that in the meantime it had dawned upon him that the verb used in the Ninth Commandment forbids \emph{concupiscentia inheridata}, and that used in the Tenth, \emph{concupiscentia actualis} (cf. Deut. 5:18, \emph{chamad} and \emph{awah}). In previous expositions he had brought both Commandments in connection with original sin, but here he says nothing whatever on that subject. Nor did he keep them apart because he realized that the object in the Ninth Commandment is in the singular while that of the Tenth Commandment is in the plural number; there is no evidence whatever that he had this difference in mind, and his contraction of two of the objects named in the Bible—ox and ass—into one—cattle—is an argument against this view. Luther’s motive was probably concern for the young that had to learn these Commandments: they would have been bewildered at finding no separate explanation for each Commandment as long as they were held to number them separately, in accordance with the traditional practise. The catechist who does not aim to supplement Luther’s Catechism, but merely to unfold it, will, for this reason, hardly take the pains to find a difference between the Ninth and the Tenth Commandment, especially since, until today, no real difference has irrefutably been shown to exist. But if he feels the need to explain why the Tenth Commandment is numbered separately, the method of the Nuremberg Sermons for Children commends itself as most practicable, provided we shall not bring in the difference between inherited
\quellseite{361}
and actual lust. There we read: “Having learned in the previous commandment that we should not desire to dispossess the neighbor of his house, position, or calling, we might come to the conclusion that while we should not desire his house as a whole, to desire a part of it, such as a good servant or ox, were no sin and that we might well try to accomplish our desire. Lest we should entertain such a thought, God forbids that, too, and says: ‘Thou shalt not covet thy neighbor’s wife, nor his man-servant, nor his maid-servant, nor his ox, nor his ass, or anything that is thy neighbor’s.’ Otherwise the thought might come to us: Oh, even though I take this or that from my neighbor, it will not harm him, for he has enough of such or similar possessions left. But the Lord does not want this; hence he summarily forbids us to covet anything whatever possessed by the neighbor.” In this connection color may be given to the explanation of the Ninth Commandment by the reflection that the house or piece of ground inherited by the Israelite had an added value in that it was to him a guarantee of his share in the promised land (cf. the laws bearing upon the year of restoration; see also Kirch. Zeitschrift, 1917, p. 372 ff.).

More important than the question as to the difference between the Ninth and the Tenth Commandment is that as to the progress made by these two upon the preceding ones, especially the Sixth and the Seventh. According to the words of Scripture, the element of progress is indicated by the verb. In the previous Commandments, the wicked word and deed are forbidden; here God forbids the wicked desire; but desire, whether we take it as inherited lust or as actual lust, is still a matter of the inner world of the heart. That these two Commandments were, in consequence, of special importance for the Jews who fondly thought that they had kept the Law if only they did not transgress it in word and work, is self-evident. More difficult it was for Luther to find something new in the last two Commandments since he was of the opinion that every Commandment is to be applied to the heart. The element of progress, accordingly, could be found only in this that the Ninth and Tenth Commandments contain an injunction against inherited lust. He writes, e.g., in Decem Præcepta Wittenbergensi prædicata populo (1518): “Videtur . . . duobus præceptis prohiberi ipse fomes et invincibilis cupiditas, ipsa, inquam, radix malarum cogitationum, ut sc. 6. et 7. intelligantur prohibiti consensus cordis et signum membrorum, sc. verbum oris et opus corporis mali, hic vero etiam ipsi primi motus una cum fomite, qui est origo illorum”; he refers to Rom. 7: 23 f. Taking
\quellseite{362}
this view, it is indeed possible for him to say that the fulfillment of the other Commandments is possible while these two are impossible of fulfillment for anyone, no matter how holy he is, and that in their light one becomes aware of his sinfulness. Through them man learns to despair of himself and his merits, and to trust in the mercy of God. So Rom. 7:7 and Luke 11:39 would be the key to the matter. While, in the Large Catechism, he holds fast to the conclusion that the element of progress is signalized by the verb, the word “covet” is to him now no longer an expression of original sin, but of actual lust, of lust to which man’s will has yielded (although he admits incidentally “that the scope reaches somewhat farther and higher”). He writes, “God has added these two Commandments in order that it be esteemed as sin and forbidden to desire or in any way to aim at getting our neighbor’s wife or possession.” It is required that no one “should think or purpose to obtain for himself what belongs to another, such as his wife, servants, house and estate . . . even with a show of right or by a subterfuge, yet with injury to his neighbor” (\emph{Weim. Ed. xxx 1, p. 175; Tr. p. 665}). Instead of original sin, he now recognizes avarice and envy as the root of all sinful conduct against the neighbor. “Therefore we allow these commandments to remain in their ordinary meaning, that it is commanded, first, that we do not desire our neighbor’s damage, nor even assist, nor give occasion for it, but gladly wish and leave him what he has, and besides, advance and preserve for him what may be for his profit and service, as we should wish to be treated. Thus these Commandments are especially directed against envy and miserable avarice, God wishing to remove all causes and sources whence arises everything by which we do injury to our neighbor, and therefore He expresses it in plain words: Thou shalt not covet, etc.” (\emph{Weim. Ed. xxx 1, p. 178; Tr. p. 669}). The impression is received that Luther, in consideration for the common man and children, has confined himself to such “ordinary meaning.” Nor is it likely that he intends to turn aside from this meaning in the concluding words: “God would especially have the heart pure, although we shall never attain to that as long as we live here; so that this commandment will remain, like all the rest, one that will constantly accuse us and show how godly we are in the sight of God!” (ibid.) What he means to say, is in all probability this that we shall not even succeed in attaining to such a degree of godliness that we never yield to the evil desire and never begrudge our neighbor his possessions. In the Small
\quellseite{363}
Catechism he takes his cue altogether from his statements in the Larger Catechism. In those of the preceding Commandments which have reference to our conduct toward the neighbor, he has as a matter of fact (Kaftan to the contrary notwithstanding) merely spoken of the words and works in which our fear and love of God is to find expression, but not, as yet, of the control that such fear and love is to exert over our thoughts in regard to the possessions of the neighbor. This, however, he now supplements in the Ninth and Tenth Commandments, where he explains “covet” as denoting “to seek to gain.” To be sure, he does not stop there, but at once proceeds to enumerate all the evil practices to which one is but too likely to resort when he seeks to gain the neighbor’s possessions; but by using in his explanation the term “to seek to gain,” the verb, and herewith the element of progress in these commandments, has not remained entirely unexpressed. The catechist that desires to remain in Luther’s footsteps will make the concept “covet” his starting-point, and explain: “Covet means to envy the neighbor his possession and to be discontented until it has come into one’s possession.” Then he shows how such coveting leads to the acts mentioned by Luther in his explanation. While he thus would not find original sin forbidden here directly he would come to speak of it as he develops the subject; for he would surely have occasion to make clear to the class how it comes that covetous thoughts will arise in us again and again (cf. our \emph{Explanation}, p. 55).

Just in view of the persistence with which covetous thoughts arise in us, fear and love are to be manifested in not only opposing them, but in doing one’s best to help the neighbor preserve his possessions. We read in the \emph{Nuremberg Sermons for Children}: “Since we are aware of the fact, my dear children, that coveting is a sin, let us guard against it, by the grace and help of God, as much as we can. But we should take special pains not to yield and conform to any evil desire . . . God has created everything and is its true Lord; therefore He gives it to whom He pleases. Has He given it to your neighbor, let him be welcome to it, and remember: if God had wanted me to have it, or if it had been of any benefit to me, He might have given it to me, too” (Reu i 1, p. 503). Tetelbach: “He who loves God does not begrudge his neighbor any blessings but is rather of help to him in keeping them” (Reu i 1, p. 683).

In regard to “house”, the Nuremberg Sermons say: “You must not merely understand the building by the little word house—the structure that the people live in, but the whole household and everything
\quellseite{364}
that pertains to it. For, in the Old Testament, the cities and villages were distributed among the Jews according to tribes, and if any one did not belong to the particular tribe occupying a certain place, it was impossible for him to obtain a house there. Therefore we may understand by the word house as much as tribe. This, therefore, is the meaning: If our neighbor belongs to a superior tribe or family; if he has inherited much property; if he is the owner of a fine household, possessing citizenship, honors and dignities, an honorable office, and everything pertaining to a householder, we should not covet his house, i.e., we should not have the least desire that he should lose it, so that we might take his place“. (Reu i, 1 p. 501 f.). Tetelbach: “By inheritance any kind of real estate is meant, such as fields, meadows, gardens, vineyards, timber, ponds, country-seats, and possessions of that sort. House means the dwelling of the neighbor and everything outside and inside that pertains to it, for instance, furniture. To seek craftily implies every kind of intrigue and wicked practise, bringing the neighbor to grief, depriving him of his income, damaging his property, or compelling him to sell out and to assign his possessions to another. Show of right means to lay claim to property by fraudulent titles and false pretenses; to outrage the law and to precipitate legal action; to bribe unscrupulous people; to advocate the side known to be wrong; to pervert the law through financial maneuvers in order to secure a favorable decision and thus to obtain the neighbor’s property”.—The second half of the Ninth Commandment should be construed in this manner: “but help and serve him in keeping it”. Compare the first Latin translation, Sed potius \emph{iuuare eum ut suas fortunas retineat integras}; the second, Sed \emph{detur opera sedulo ut ista proximo diligenter custodiantur}; the fourth, Sed \emph{ad conservanda ea ut prompti et parati simus}. “Estrange” means to alienate by allurements or promises; the German word \emph{abspannen} is by no means identical with \emph{abspannen}=to unhitch; its root is either \emph{spannen}=persuade, lure, entice, deceive (cf. High German \emph{Gespenst}), or \emph{spenen}=to wean (cf. \emph{Spanferkel}—a little pig recently weaned). Luther preferably connects the word with wife (cf. Weim. Ed. xxx 1, pp. 9. 42. 84). Tetelbach, while mistaken in regard to the derivation, gives the correct explanation: “With smooth words and promises to persuade them to neglect their calling and duty and to cease doing right, as when a horse is unhitched or a crossbow unstrung.” “To force away” Tetelbach explains as follows: “By threats, force, and undue pressure to urge them until they leave the neighbor”. “To entice away” means, according to the same author, “to steal one’s
\quellseite{365}
heart and favor by gifts, matchmaking, and other unscrupulous measures so that he turns against his environment and wishes for another.” The first Latin translation has “\emph{ad nos arte aliqua transferre, non vel vi vel dolo aliquo eos ab aliis abalienare}”; the second, “\emph{abalienemus aut abstrahamus}”; the fourth, “\emph{ne ulla ratione abalienemus et ad nos alliciamus uxorem proximi, servos, aut quicquam ex ipsius familia}.” Kaftan is right when he says: “All three expressions have substantially the same significance; ‘to entice away’ is the general term, ‘to estrange’ refers to the person enticed away and denotes that this is accomplished by allurements; ‘to force away’ refers to the owner and denotes the use of persistent urging.” That Achelis is mistaken when he applies the term ‘estrange’ to animals, ‘force away’ to servants, and ‘entice away’ to the wife (in Der Dekalog als katechetisches Lehrstueck, p. 53), requires no argument.

In regard to the function assigned by Luther to the Conclusion, Albrecht (in Weim. Ed. xxx 1, p. 361) recapitulates the matter bearing on this subject as follows: “In deference to tradition, Luther left out of his text the prohibition of images together with the threat and the promise. That he accounted these words as parts of the First Commandment, is evident from his comments on Exodus (1525, cf. Weim. Ed. xvi, pp. 436 ff.; 445, 23) where we find the significant word: ‘The Promise and Threat Contained in the First Commandment.’ Aside, however, from this significant statement, Luther, in his earlier catechetical expositions, failed to take note of this addition, although he maintained from the outset the fundamental, controlling importance of the First Commandment. Not until the second and third series of Catechetical Sermons of 1528 (W. E. xxx 1, pp. 43. 85), did he take the addition into consideration. Here he emphasized the fact that the addition was virtually part of each commandment (ibid., p. 43), and that he had appended it at the end of the commandments for the purpose of ‘joining them into a wreath, the last to the first.’ The threat with its power of arousing fear and the promise with its power of arousing love are to him the ‘cord with which the wreath is tied together’ (ibid., p. 85). Quite in harmony with these statements is the exposition in the Large Catechism where the conclusion receives attention both in connection with the First Commandment and after the last” (ibid., pp. 136 ff., 179 ff., Tr. pp. 589 ff., 669 ff.). That this does not justify the inference that the “fear” of the First and the succeeding commandments is to be understood in the same sense as in the conclusion, has already been dwelt
\quellseite{366}
upon (p. 327 f.; compare also A. Hardeland, \emph{Der Begriff der Gottesfurcht in Luthers Katechismus}, Guetersloh, 1914, against J. Meyer, \emph{Fuerchten, lieben und vertrauen in Luthers Kleinem Katechismus}, in \emph{Neue Kirchl. Zeitschrift}, 1913, pp. 793 ff.). In regard to the position of the conclusion see also p. 101. The catechist who favors Luther’s arrangement of appending the threat and the promise to the Commandments as a whole will certainly let his treatment of the Commandments culminate in an exposition of this threat and this promise as two powerful incentives to obedience. To raise the question at this point in regard to “the use of the Law”, is to such a catechist an incomprehensible procedure. Has he not in all preceding catecheses tried to direct his pupils’ eyes to the Law as the norm of life, which certainly is a use of the Law? No one would ever have thought of introducing such a discussion at this point but for the serious misconception that dogmatical topics must be interpolated in the Catechism. Again, the catechist following in Luther’s steps will not think that he must pause to show the children their sin in the light of the Law, true as it is that the Ten Commandments can serve as the mirror for sin, and should serve for it at the proper time, namely in connection with confession. Still less will he now take the opportunity to convince them of their inability to keep the commandments. Nobody would have been led to these two pedagogic follies had it not been for the desire to systematize the Catechism by linking the two Chief Parts together. On the contrary, just as it was the catechist’s aim to show his pupils in connection with each Commandment how child-like fear and love of the Father should impel them to the fulfillment of the divine will, even so he reminds them, for the same reason, of the wrath and jealousy divine which is to be visited upon the transgressors of the Law, and of all the blessings in which the obedient may rejoice, in order to help them by all means to learn obedience to the divine will. A treatment of the Commandments which conforms to Luther’s pattern, dare not sound as final note the confession: “I am a sinner, groaning under the curse”, nor the conviction: “I am too feeble to fulfill the Commandments of God”. Rather should it lead to the resolution: “As the child of my heavenly Father I mean to fulfill the will of my God and Father, and to fulfill it cheerfully” (it is just this voluntariness and cheerfulness of an obedience prompted by fear and love which Luther emphasized again and again), and to the prayer: “Help me, my God and Father, better and more perfectly to fulfill Thy will!”

The text of the Bible contrasts those that love God with those
\quellseite{367}
that hate Him. It speaks of hating because, according to its very nature, sin is rebellion and hatred against God (Rom. 8:7), a fact quite in evidence in the case of those who wilfully and defiantly disregard the commandments. “Unto the third and fourth generation of them that hate Me” and “showing mercy unto thousands of them that love Me” are phrases bringing out the energy, the earnestness, and holy zeal of God, both when He inflicts penalties and when He scatters blessings. The thought of them is to add strength to the child-like fear and love in our hearts toward the heavenly Father as an incentive to obedience. Both punishment and blessing bear primarily upon the bodily life; the fundamental principle is brought out that sin is followed by evil, obedience by blessing. In connection with transgression we are to think of every kind of misfortune, failure, trouble, sickness, etc.; and in connection with obedience, of blessing in body in that God preserves our health and strength, of blessings in soul in that He rewards obedience with peace and joy, of blessings in property in that He prospers our work, of blessings in honor in that He gives us a good reputation among all godly men. But because these results are not invariably the instantaneous fruit of obedience, Luther, once more, rightly lays stress upon trust. If God permits us to languish in misery in spite of our obedience, it behooves us to trust Him for finally changing suffering into bliss, even though we should have to wait till eternity’s day. The words employed are broad enough to embrace this last thought. Well it is said in the Nuremberg Sermons for Children where the Conclusion is brought in connection with every one of the Commandments, and each sermon ends with the following words:

\begin{quote}
“Thus, my dear little children, you have the correct, simple, everyday meaning of this Commandment. You should lay it diligently to heart and fear God lest you transgress this commandment; for what He commands is good and holy, and what He forbids is wrong, sinful, and shameful. It is the determined will of God that we should keep His commandments and not despise them; for He says: ‘I, the Lord, thy God, am a jealous God’. . . . Therefore, my dear little children, you should fear God and diligently keep His commandments, imploring His grace and help so that you may be enabled to do this. For the fear of the Lord is the beginning of all wisdom and makes us pious and good, pleasing in the sight of God and of service to other people; and we shall enjoy peace and quiet, honor, prosperity, and happy days. If we thus continue in obedience to God and in the right faith to the end, God will add eternal life as the last gift. May He grant
\quellseite{368}
this to us all. Amen.”
\end{quote}

This is the final note to be sounded in an exposition of the Commandments. Should anyone think that we should thus turn into the Roman channel, he merely reveals his failure to understand either the Scripture or Luther. Albrecht says aptly: “There is no trace here of pleading of merit; trust clings to a promise of grace, not of a reward” (W. E. xxx 1, p. 362). The believer, justified by faith, who, if all is well with him, endeavors to fulfill the will of God by inner necessity, finds an additional motive and incentive in the “grace and every blessing” promised of God upon the sole condition that we obey.

The catechist will much better succeed in impressing upon the heart of his pupil the thought that the regulations in the Father’s house, as found in the First Chief Part, are to be faithfully and scrupulously kept, if he sums up each Commandment by reaching into the life of Jesus, the perfect fulfillment of the Law. The fulfillment of this or that single requirement is typified by one person of sacred history or another; but of the fulfillment of every requirement we have in Jesus the most winsome and attractive type: all rays of obedience meet in Him; not one feature is lacking. Such treatment as this will arouse in the catechumen the conviction that the object to be realized is to follow Christ; that it is His image that is to shine forth from our life. Thus the New Testament character of the First Chief Part stands duly revealed. Leopold Schulze has laid great stress upon this thought in \emph{Katechetische Bausteine} (Magdeburg, ⁵1891); in our own \emph{Explanation of the Catechism} it has been carried out consistently.

The needlesness of a transition to the \emph{Second Chief Part}, and Luther’s attitude in the matter, have been discussed above (pp. 111 f.; 319). Here we only add what Albrecht says on the subject (Weim. Ed. xxx 1, p. 362): “In the Small Catechism the \emph{Second Chief Part} is adjoined to the First without any connecting thought. This is the case also in the third series of Catechetical Sermons of 1528 where we read in the sermon of December tenth: ‘You have heard the first part of Christian doctrine, namely, the Ten Commandments. I have admonished you diligently that you should exhort the family to learn that part literally by heart so that they might obey God accordingly, and you as their masters, and yourselves likewise might obey God. For if you will instruct and urge the family, there will be progress; and there has never been a doctor who has failed to become more learned by study. These Ten Commandments having been dealt with, we will now proceed to the next part’ (xxx 1, p. 86). At other places
\quellseite{369}
connection is made by harking back to the First Commandment, or by asking the question, \emph{Qualis est deus tuus qui tulit tibi decem praecepta?} (xxx 1, pp. 9. 10. 183), or by reviewing the connection given in the short form [above pp. 85, 111 f., 319]: those incapable of fulfilling the commandments [for the disease of sin does not here come into consideration as guilt, as is still maintained by many, but rather as moral impotence] will receive help and strength from God through faith (xxx 1, p. 43 f.; xi, p. 48). This thought though expressed less forcibly, returns also in the Large Catechism (xxx 1, p. 184).” One who explains the Introduction, the First Commandment, and the Conclusion according to Luther’s plan, has learned that the First Chief Part presents all of Christianity, insofar as it is a divine requirement for the baptized Christian, and he will not only feel no need for the interpolation of transition material but will find it absolutely incongruous.

In the \emph{First Article} Luther places all possible emphasis upon the word Creator or “Maker.” Not only does he give this article the heading “Of Creation”, but, in the edition of 1531, revised and corrected by him (cf. p. 95), he had the printer capitalize every letter of the word Schoepffer (just as the word Herr in the Second Article). He made it the keyword in his explanation. Previously his mode of operation had been different. In the Short Form he did not emphasize the objects of faith, but the verb—the concept of faith, although mention is made of the former (cf. p. 90). In the sermons of 1523 he lays more stress upon the objects of faith, especially on God as the “almighty Father” (Weim. Ed. xi, p. 49); similarly in the rather brief allusions in the first two series of sermons of 1528. In the third series, however, Luther confines himself altogether to the word “Maker”: “The children and the common people are to learn the matter in short, plain, and simple terms as follows: There are three articles in the faith, (1) of the Father, (2) of the Son, (3) of the Holy Spirit. What do you believe of the Father? Answer: He is the Creator. What of the Son? He is the Redeemer. What of the Holy Spirit? He is the Sanctifier. For the benefit of the learned the article might be divided into as many parts as there are words; but now I would teach the young and the common people that God the Father is the creator of heaven and earth. Now, what is the meaning of that? Just this, that I am to believe that I am God’s creature; that He has given me body, soul, sound eyes, reason, property, wife, children, fields, meadows, hogs and cattle. Then that He has given the four elements, etc. This article, accordingly, teaches that you do not have your life of yourself, not
\quellseite{370}
even a hair. All things in existence are comprehended in the little word Creator” (Weim. Ed. xxx 1, p. 87). Such is the tenor also in both the Large and Small Catechisms although here, in the phrase “fatherly goodness and mercy,” the word “almighty Father” comes into its own. Whether this thought served as basis for the following statements “clothing and shoes . . . . . from all evil” is doubtful.—That Luther used the word “almighty” as an adjective, connecting it with “Father,” not with “creator,” which would be a tautology, hardly requires any proof. In the edition of 1531, so carefully revised by him, he puts a comma after “almighty”; also in the first edition of the Large Catechism (W. E. xxx 1, p. 130); then, after repeatedly omitting all punctuation, again in the reprint of the third edition (1540) and all subsequent prints; with unmistakable clearness it is indicated in that part of the exposition where we read: “But to speak more of this belongs into the other two parts of this article where we say ‘Father almighty’” (W. E. xxx 1, p. 184.: Tr. p. 681, § 17, 18). Capito (1528) and Brenz (1529) write without hesitation “almighty Father”. See M. Reu, \emph{Die Konstruktion des ersten Artikels} in Kirchliche Zeitschrift, 1904; Cohrs, \emph{Katechismusversuche vor Luthers Enchiridion} iv, pp. 295 f. That the right connection was retained for a long time after Luther, can be abundantly demonstrated from our sourcebook, \emph{Quellen zur Geschichte} etc.; why Luther did not simply write “almighty Father”, transposing the adjective as Capito and Brenz did, we have seen (p. 100).

Another problem is closely connected with, and to some extent solved by, the problem just considered. We refer to the question whether God comes into consideration primarily as “Father of Jesus Christ” or as “our Father.” Now, God is unquestionably considered here, in some way, also as the Father of Jesus Christ; this appears from the origin of the Creed as a baptismal confession, and from any number of New Testament passages which show that God is our Father only through Jesus Christ, His incarnate Son, coequal and coeternal with God the Father. But that this relation is to be chiefly emphasized at this place, as is conventionally done, or that it is the sole fact to be considered, is totally erroneous and most certainly not in harmony with Luther’s views. He writes in the Short Form: “Because He is God, He can do the thing that is best for me, and knows what that thing is. Because He is Father, He wills to do what is best for me, and to do it with all His heart. Because I do not doubt, but put my trust in Him, I am assuredly His child, His servant and His heir forever, and as
\quellseite{371}
I believe, so will it be done unto me.” In the second series of his Catechetical Sermons of 1528 we read: “\emph{Pater hat sich ausgeschuettet cum omnibus creaturis quas creavit}.” In the Large Catechism the same thought occurs in this form: “And He does all this out of pure love and goodness without our merit, as a benevolent father who cares for us that no evil might befall us” (Tr. p. 681). It is found once more in the hymn of faith: “\emph{Wir glauben all an einen Gott, Schoepfer Himmels und der Erden, Der sich zum Vater geben hat, Dass wir seine Kinder werden}.” In his \emph{Hauspredigt von den Artikeln des Glaubens, in Schmalkalden gehalten} (1537) he says: “Moreover, He not only teaches us what and whence we are, but also where we belong. This is indicated by the word “Father”, which means that God lays claim to fatherhood as well as to omnipotence. The animals cannot call Him Father; but we can call Him Father and ourselves His children. With that word He expresses His purpose in regard to us. Previously we had learned whence we are and what honor and glory are ours; now the question arises as to our destiny. It is this that you shall be children, and I Father. I have not only created you and seen to your preservation: I want you for my children and appoint you as my heirs so that you shall not be thrust out of the house like other creatures—oxen, cattle, sheep, etc., which, without exception, die or are devoured; but in addition to being My creatures, you shall be also My children for eternity. This is our prayer and confession when we recite the creed: I believe in God the Father” (W. E. xlv, p. 16). This was the sense in which, before the appearance of the Small Catechism, Brenz, Althamer, Lachmann, and others interpreted the word “Father.” This was its interpretation also after its appearance. Matthaeus Zell (1536) writes: “What does he mean when he says: Father? This is the most comforting thought that we may entertain concerning God. Why? In this one word the truth is expressed that the high divine Majesty loves us poor worms as heartily as a father, and not as a stepfather; that He desires us for children and for heirs of an eternal blessed life—a thought than which there is no higher and surer comfort in the extremity of death and in any other ills for anyone who, in true faith, speaks and confesses that word . . . . What do we mean when we confess that God is almighty? First of all, the man that really lays hold of that truth finds it a comfort exceedingly great. How is that? Not until we realize that He is almighty, have we the sure comfort that no one can prevent Him from doing as He pleases, unlike human parents who are often prevented from carrying out their will. But
\quellseite{372}
God, being almighty, can be prevented by no one: no matter when His help is required or by whom it is required, no one shall be too strong and powerful to be in His way” (Reu i 1, p. 109). The \emph{Nuremberg Sermons for Children} are less clear on the subject; but they give a splendid explanation of the term “almighty”: “Nobody is so sick but He can heal him; nobody so poor but He can make him rich; nobody so despised but He can bring him to honor; nobody so great a sinner but He can make him godly; nobody so bereft of faith but He can make him a believer; nothing so incredible but He can perform it if He will” (Reu i 1, p. 69). Butzer is clearer (1537): “I believe that God, the Father of our Lord Jesus Christ and the first person of the Godhead, desires to be my Father also. How is He your Father? In that He imparts to me His nature and heritage. What is God’s nature? All righteousness, godliness, and love. What is His heritage? Eternal, blessed life” (Reu i 1, p. 69). Huber (1544): “Therefore I believe and confess that Almighty God, the only, eternal God, is the almighty Father of myself and all the faithful” (Reu i 1, p. 792). Meckhart (1553): “Why do you call God your Father? Because He, through grace, has elected me as His child, promised me fatherly faithfulness and every blessing, and made me an heir of all His possessions” (Reu i 1, p. 824). But even those who take the word in the trinitarian sense, do not overlook the other aspect. A case in point we find in Kantz (1542): “Why do you call God Father? Because He has a Son of the same essence as His own, born from eternity, who is God equal with the Father . . . And through the same beloved Son He has elected as children all those who believe in Him so that they may inherit His kingdom with Him; John 1:12; Rom. 8:17. Such is the meaning of the word Father which is an exceedingly gracious, comforting word, assuring us that God is our faithful Father, and will be such forever; Isa. 49:15” (Reu i 1, p. 608). Or Tetelbach (1568): “Why do you call Him Father? He is such; for He has brought forth from His divine nature and essence a Son—our Lord Jesus Christ through whom He has also become my dear Father and the Father of all believers in the Lord Jesus Christ; and He is willing to deal with us as does a father with his children” (Reu i 1, p. 658).

Understanding the First Article thus we can comprehend why Luther, in 1523, could assign first rank to the First Article: “Hic primus articulus, quod credo Deum esse patrem meum, qui articulus summus est . . . Quicquid est in tota Scriptura, huc referri potest. Si deus pater meus est, liber sum ab omnibus, inferis, morte, mundo etc.” (Weim.
\quellseite{373}
Ed. xi, pp. 50 f.); and in that very connection he makes the words “almighty Father” the nucleus of his whole explanation. It is quite plain that Luther did not intend to cast aside Christ’s death and resurrection and the atonement conditioned by them; Luther says often enough that God is our Father only through Christ. And it is only when this is borne in mind that Harnack’s statement is acceptable: “The First Article is for Luther an epitome of Christianity” (\emph{Dogmengeschichte} iii, p. 711, note 3). From this standpoint another matter becomes clear, namely, that it is a thorough misunderstanding of the First Article when the view is still entertained as though it were nothing else but general faith in providence, made so much of in the era of rationalism, which is here confessed, and as though the First Article represented a kind of preparatory stage, a sub-Christian creed which rationalist and Lutheran Christian could very well join in confessing. Nothing could be more erroneous. That a faith in providence is here spoken of, is true; but it is the faith in providence as entertained by the Christian—a faith based upon his faith in redemption, yes, a faith which, far from being inferior, is nothing less than an expression of it. It is God, the almighty Father and Creator, of whom I am here speaking, who has become my Father through no one save Christ, and only through His redemption. Likewise, it is the baptized Christian that here confesses his faith, already a child of God through baptism, and now looking up to Him as Father and finding comfort in His care and protection. Instead of being something sub-Christian, the assurance of faith here manifest is specifically Christian and radically different in principle from the natural man’s faith in providence. It finds its true parallel in that stanza of Paul Gerhardt’s: “With gladness I confess it, What long my heart has known: That God, most high and blessed, To me His grace has shown; And that, whate’er betide me, He by my side shall stand, Through storm and wave to guide me, And every ill at hand.” Only thus can we justify the fullness of assurance of Luther’s words: “He richly and daily provides me with all that I need for this body and life, protects me against all danger and guards and keeps me from all evil.” The more the uniqueness of the Catechism is preserved also in this respect, the more clearly the gospel will stand disclosed.

In regard to the explanation of the First Article, a difficulty arises from a mooted point of construction. So much is clear that Luther’s explanation is readily divided into three parts: (1) What God my Maker has done, and still does, for me; (2) Why God has done and
\quellseite{374}
still does this; (3) What I owe him therefore. But how to construe, in the first part, the words \emph{dazu Kleider und Schuh . . . . . reichlich und taeglich versorget} is the question. Frequently it is advocated to subordinate the whole section \emph{dazu Kleider und Schuh . . . . . bekleidet und bewahrt} to the words \emph{und noch erhaelt}, and still preserves. Divine preservation takes place when God supplies everything which I need for the preservation of life and when he safeguards me against everything that might cause me injury. Such, it seems, was the construction preferred by H. Opitz (1583) who answers the question, “Whereby does God preserve our human body and life?” in this manner: “By clothing, shoes, . . . cattle and other goods, together with all that is needed to support this body and life” (Reu i 1, p. 532). With him agrees H. Magdeburg, so far as his Greek translation is concerned (the Latin one is different: v.i.): “\emph{καὶ νῦν ἔτι τρέφει, ἐπιμελούμενός μοι τῶν ἐσθήτων, τῶν ὑποδημάτων . . . . . οἰκίας, πρὸς δὲ τούτοις τῆς γυναικός . . . . . πρόνοιαν ἔχων, ὅπως ἂν ταῦτα ἱκανῶς καθʼ ἡμέραν ἑκάστην ὑπάρχῃ}.” The result of this construction was that the word \emph{dazu} (= also) was no longer taken in the sense of “further”, but connected with the preceding “still preserves to me”, and thus understood to mean “for this purpose”, namely for the purpose of preservation. After the words “and all my goods” a semicolon was placed, the missing verb being supplied from what preceded, i.e., “gives”. That this cannot be the meaning of \emph{dazu}, becomes clear from the Large Catechism and the Latin translations according to which \emph{dazu} is equivalent to “further”. This is the very circumstance that occasions the question as to the predicate of this part of the explanation whether it must be found in one of the preceding verbs or in one that follows, that is, in \emph{gegeben hat und noch erhaelt} or, perhaps, in \emph{versorgt}.

The last named construction enjoyed great vogue during the sixteenth century. We find it in the first Latin translation: “\emph{Credo, quod Deus me una cum omnibus aliis creaturis creaverit, quod mihi corpus et animam, hos oculos, aures, omnia alia membra, quod rationem et illos sensus donaverit. Neque id solum, sed credo quoque, quod omnia illa alias peritura sustentet, quod vestes, calceos, cibum et potum, certas sedes, uxorem, liberos, agros, iumenta et quidquid bonorum est, abundè ad sustentandam hanc vitam mihi impertiat praebeat. Credo etc.}” Again in the second translation: “\emph{Credo, quod Deus creavit me una cum omnibus creaturis: Quod corpus et animam, oculos, aures, et omnia membra, rationem et omnes sensus mihi dedit et adhuc sustentat. Ad haec quod vestes et calceos, cibum ac potum, domum, uxorem, liberos, agros,}
\quellseite{375}
\emph{iumenta cum omnibus vitae necessariis copiose et quotidie largitur: „Me contra omnia pericula protegit etc.”} Likewise in the third Latin translation: “\emph{Credo, quod Deus Pater me creavit . . . omnesque sensus dedit et conservat. Victum, vestitum, cibum, potum, domum, fundos, uxorem, liberos, iumenta et omnia bona, omnia necessaria ad vitam quotidie suppeditat, quod contra omnia pericula nos defendit, etc.}” Nor does the fourth Latin translation, by Hiob Magdeburg, show a different construction: “\emph{Credo, me cum tota rerum universitate a Deo creatum esse, qui corpus mihi et animam . . . sensusque omnes et largitus est et conservat. Vestibus mihi insuper et calciamentis, cibo et potu, aedibus, uxore et liberis, agris, pecoribus et ceteris bonis ad conservandam vitam necessariis abunde quotidie prospicit.}” Similarly in the Greek rendering in the Catechesis Graeco-Latina: “\emph{ἔτι τὰ ἐνδύματα καὶ τὰ φράγματα . . . καὶ πάντα τὰ ἄλλα ὑπάρχοντα . . . δίδωσιν}.” Accordingly, the words, \emph{dazu Kleider und Schuh . . . versorgt}, were accepted as a coherent statement, and the word \emph{versorgt} (=to provide someone with something) was taken in the sense of \emph{besorgt, schafft herbei} (=to procure something for someone), with the accusatives \emph{Kleider und Schuh} etc. as its object. \emph{Mit aller Leibsnahrung und Notdurft} would, accordingly, be detached from versorgt and become: “and, in addition, all that I need for this body and life”. The English translation corresponding most closely to this construction is found in Cranmer’s Catechism of 1548: “I believe that God the Father hath made me and all creatures in heaven and hearth, that he hath gyven to me and conserveth me bodye and soule, reason, senses, eyes, eares, and all my other members. And I believe that the same almightye Lord and God doth dayly gyve to me and to us all meat, dryncke, cloth, wife, children, house, lande, riches, cattell, and all thynges necessarye to the mayntenaunce of our lyves etc.” Joachim Moerlin (1547 or 1554) does not hesitate to paraphrase: “That God the Father daily and richly gives me these (namely body and life) and all things necessary to this life and temporal food”. The same construction is found again, e.g., in Marbach, 1557; in Aquila, 1538; in Bischoff, 1599; in Hercko, 1554; in Siber, 1575; in Dresser, 1581; in Hadamarius, 1552; in Maurer, 1575; in Joseph Opitz, 1538. The usage of versorgt as equivalent to besorgt was evidently not unknown to any of these; and because an example is found also in Luther (in a letter to Hans Schotten, May 1524, he says concerning the duties of the father: “\emph{er ist schuldig, dem Kind Essen und Trinken, Kleider . . . zu versorgen, für des Kindes Not und zu seinem Besten}”; cf. Erl. Ed. liii, p. 243),
\quellseite{376}
it might appear that this construction is the right one. It would, in that case, have the additional advantage of giving in the words: “I believe that God has made me my reason and all my senses” . . . a definition of the word “Creator,” and in the words “also clothing and shoes . . . guards and keeps me from all evil” an exposition of “Father almighty”. What an earthly father does for his children, Luther would raise to the highest degree and apply to the almighty Father.

Notwithstanding, this is hardly the correct construction. Taking the Small Catechism alone into consideration, there are three points which arouse our suspicion. (1) This construction would express that God richly and daily (!) provides me with wife and child; (2) it is incongruous, after “all goods”, an expression indicating that the enumeration is finished, to continue immediately with the words “and in addition, all that I need for this body and life”; and (3) even the rhythm, by no means unimportant in Luther’s explanations, requires a pause after “goods”. More significant is another factor: Luther’s preliminary labors militate against this construction. In the sermon of December 10, 1528, we read: “The meaning is that this I am to believe that I am God’s creature; that He has given me body, soul, sound eyes, reason, property, wife, children, fields, meadows, hogs, and cattle. Then that He has given the four elements etc. This article, accordingly, teaches that you do not have your life of yourself, not even a hair. All things in existence are comprehended in the little word Creator” (W. E. xxx 1, p. 87). In the Large Catechism we read: “What is the force of this, or what do you mean by these words: I believe in God, the Father almighty, Maker, etc.? Answer: This is what I mean and believe, that I am a creature of God; that is, that He has given and constantly preserves to me my body, soul, and life, members great and small, all my senses, reason, and understanding, and so on, food and drink, clothing and support, wife and children, domestics, house and home, etc. Besides, He causes all that exists to serve for the uses and necessities of life, sun, moon, and stars in the firmament, etc.” (W. E. xxx 1, p. 183; Tr. p. 681). These two quotations inevitably lead to the following conclusions: (1) That the section, “also clothing and shoes . . . all my goods” belongs to “he has given and still preserves to me” and that the objects here enumerated (though following the verb, in the German original) are to be taken as coordinate with the first named objects (body, soul, etc.). (2) That there is by no means to be supplied, as, for instance, Luehrs and Zezschwitz require, a \emph{mich} in connection with \emph{erhaelt} which would produce the
\quellseite{377}
following text: “that He has given me my body and soul, eyes, ears, and all my members, my reason, and all my senses, and still preserves me”; but that the \emph{mir} belongs both to \emph{gegeben hat} and to \emph{erhaelt}, thus: “that He has given and still preserves to me my body and soul etc.” This, then, is the thought expressed: I believe that God has given and still preserves to me body and soul, eyes, ears . . . . all senses, further, clothing, shoes, meat, drink, . . . and all goods; and that He also preserves to me all these objects. In preserving them He, indeed, preserves also me; but by interpolating mich as the object of erhalt we would destroy the whole sentence structure that Luther intended to rear. (3) That the words mit aller Leibesnahrung und Notdurft . . . versorgt constitute a new, complete thought, in the wording of which the terms \emph{versorgt} and \emph{mit} are to be taken in the usual sense, instead of being exchanged for besorget and nebst. The absence of a direct mich does not disprove this statement, for Luther often omits the personal pronoun in the second part of a sentence, and does this even if the cases do not agree (cf. the sentence in the Large Catechism which expresses the very same thought: “\emph{Da sehen wir, wie sich der Vater uns gegeben hat samt allen Kreaturen und aufs allerreichlichste in diesem Leben versorgt}, Tr. p. 682, § 24). (4) That the conventional interpretation of \emph{Notdurft und Nahrung dieses Leibes und Lebens} according to which these words are a summary of the blessings previously enumerated, is an error (even Kaftan still adheres to this error although in other respects he has consistently rejected the construction of the Latin translations). Luther, far from virtually repeating his statement, thinks of every necessity of this life not previously named, as air and light, fire and water, sun, moon, stars, etc. Careful reading of the Large Catechism shows that Luther in this phrase includes everything enumerated by him in the explanation of the Fourth Petition alongside of “clothing, shoes . . . all goods.” That thus, especially through his reference to the elements, a vital and timely thought is introduced, is evident to everyone who knows how often in our day nature is placed alongside of God, if not above Him, as something independent (cf. p. 315). Compare the lines of Paul Gerhardt: “Heaven, earth, and all their host, As my servants He has given, etc.” When we remember that Luther had previously made mention of the “creatures”, and when we have come to realize that the words “has given my body and sould” down to “richly and daily provides me with all that I need for this body and life,” is nothing but an unfolding of the preceding main clause: “I believe that God has made me and all
\quellseite{378}
creatures”, there will be still less desire to eliminate his reference to air, light, etc. (5) From the Catechetical Sermon quoted above (p. 354) and the Large Catechism it is clear that the giving, preserving, providing, here mentioned are by Luther intended to serve as a definition of divine creation. For this reason it is totally wrong to add to creation “preservation” which is said to express itself in giving (clothing, shoes, etc.), providing and protecting (so Opitz, H. Magdeburg); and even greater is the error when also “government” is intercalated. The benefit (?) resulting from this procedure is that we have in happy union the dogmatical trio—creation, preservation, government. “All is comprehended in the word creator. Emphasize the word creator!” says Luther. In the Small Catechism he even goes beyond his preliminary labors in that he so extends the statement: “I believe that God has made me etc.”, as to find in it the notion of protecting, guarding, and keeping. This has been recognized, e.g., in the Joachimstal Catechism (1574) where we read: “Creation does not merely mean that God once upon a time made everything but that He is present with His creature and active for its sake, as the words of the explanation (down to “from all evil”) and the hymn (Luther’s hymn on the Creed) show” (Reu i 2, p. 696). All this—giving, providing, protecting, etc., can, according to the Large Catechism, be classified under a double head: (1) God has given us everything we have and see before us; (2) He also wards off whatever could hurt me and mine.

The Catechetical Sermons of 1528 and the Large Catechism justify the inference that all the statements in the Small Catechism are designed to lay stress upon God’s activity as contrasted with human wisdom and strength. In the Sermon of December 10 we read: “This article teaches, accordingly, that you do not have life of yourself, not even a hair . . . Whatever you may possess, no matter how little it be, remember this when you say ‘Maker.’ Even though you wear a wreath on your head—let us not think that we have made ourselves, as do the proud princes . . . Nothing of all this I have of myself. The Maker, that is, God, has given everything . . . If all things are God’s gifts, you owe it to Him to praise and thank Him for the same . . . How many, I ask, are in the world who understand this word ‘Creator’? For no one serves Him. Accordingly the First Article may well humble and dismay us because we do not believe . . . He, however, who believes this article, when he looks at his cow, will say, ‘The Lord has given it to me.’ The same in regard to his wife, his children . . . I have nothing of
\quellseite{379}
myself lest I should become proud. I can neither give it to me nor preserve it” (W. E. xxx 1, p. 87 f.). Compare also the Nuremberg Sermons for Children (Reu i 1, p. 510). This, then, is the construction of the first part of the explanation: I believe that it is God who made me together with all creatures; that is (1) I believe that whatever I am and have is from Him alone; (a) I believe that it is God who has given me my body and soul; (b) I believe that it is God alone who also preserves to me my body and soul; (c) I believe that whatever else I have besides body and soul is given me and preserved for me by Him alone; (d) I believe that it is God alone who also provides me with all I need for this body and life (for, being Lord of the world, He gives me rain and sunshine, light and air, fire and water; yes, He has appointed heaven and earth, air and sea for my service); (2) I believe that it is God who wards off whatever could hurt me and mine; (a) I believe that it is He alone who protects me and mine from all danger, (b) guards me from all evil, (c) keeps me from all evil. Should the argument be urged against this construction that the “me” it “provides me with all that I need for this body and life” is an intercalation (\emph{mit aller Notdurf und Nahrung dieses Leibes und Lebens reichlich und taeglich versorgt}), let it be remembered that the other construction has likewise an intercalation, in that the “me” in “keeps me from all evil” is not in the original (\emph{wider alle Faehrlichkeit beschirmet}); and that Luther has by no means always been exact in the use of the personal pronoun so far as the Small Catechism is concerned. A case in point is the second half of the explanation of the Fifth Commandment or the third part of the explanation of the First Article. It is, accordingly, quite likely that Luther when he used the verbs “protect, guard, keep”, had in mind as object not only “me” but every object enumerated in the preceding explanation. For this reason we have explained in the epitome above given: “Me and mine” (see Large Catechism, Jacobs, Book of Concord, p. 440, § 17; Tr. p. 680 adds uns). On the whole question of construction compare Reu, \emph{Die Konstruktion des ersten Artikels} in \emph{Kirchl. Ztschr.}, 1904, pp. 28–33; 83–89.

Finally we must note how well Luther succeeds in filling the article of Creation with contents of vital significance for the believer. He attained this end principally by relating God’s creative activity to the individual (cf. p. 105), but also by viewing all creatures as existing for our benefit (e.g. in “provides me with all that I need
\quellseite{380}
etc.”); and the fullness of assurance with which Luther confesses that God protects, guards, and keeps us from all evil, has its basis in just this faith that God is the maker of all creatures: He is the Lord of all, no creature can thwart His plan, all of them are bound to be His instruments in protecting, guarding, keeping His children. Magnificent is Luther’s presentation of the subject elsewhere (Walch viii, p. 318): “If we should behold the downfall and collapse of the world and all its elements, and ourselves floundering in the wreck, we still should say: You will not fall though you do fall unless it be the will of God. And though the world should lie, a huge weight, upon our very head, we still should say: World, you can not hurt or crush me; but if it should please God to have me felled and crushed by your weight, let there be done in the name of the Lord whatever pleases Him: my time is in His hands. But should the contrary be His design, I will defy you, heaven and earth, and all power of pope and Turk and the whole world’s wrath with it” (cf. p. 296). The same thought he expresses in the Short Form: “If He is almighty, what can I lack that He cannot give me and do for me? If He is Creator of heaven and earth and Lord of all things, who will take anything from me, or harm me?” (above p. 85). Not all of the catechists of the sixteenth century have been unmindful of the fact here brought out. Compare the Nuremberg Sermons for Children (Reu i 1, p. 405), and what Meder (1590) says on basis of these sermons: “If God has made heaven and earth, He is also Lord of them; and everything in heaven and earth must move according to His will; and if He has made everything for His sake, everything is bound to serve us” (Reu i 1, p. 495).

For lack of space, the remainder has to be compressed into a few brief notes; for a more detailed discussion see M. Reu, \emph{A New English Translation of Luther’s Small Catechism}, 1926. In the text of the Second Article, the edition of 1531 emphasizes the word “Lord” as the one most important by capitalizing every letter: HERR. In keeping with this conception, clearly expressed in the sermon of December 10, 1528, and in the Large Catechism where that sermon is largely reproduced, Luther chooses the heading “Of Redemption”, and groups the whole explanation around the statement, “I believe that Jesus is my Lord”. Already at an early day, this simple and lucid construction has unfortunately been misunderstood by some (for instance, H. Magdeburg: “\emph{Credo Jesum Christum verum}
\quellseite{381}
\emph{esse Deum et a patre ab æterno esse genitum: veritum item ex Maria virgine hominem natum, eundemque dominum meum.}” Still worse is W. Han’s product, of 1553; Reu i 2, 2, p. 583; i. 2, 1, p. 459); later it was ignored almost universally in the interest of dogmatics. If there is any point in regard to which the Large Catechism, here a virtual reproduction of that sermon of December 10, 1528, can serve as commentary upon the Small Catechism, it is the one under consideration; and for this reason everyone should study the Large Catechism who would grasp Luther’s intentions and escape the danger of mixing foreign ingredients with his Catechism.—“Redeemed, purchased and won” are words expressing one and the same thought; but each time under a different image. The first of these words reminds one of the great price paid for us; the second of the arduous labor necessary; the third of the fierce battle that was fought.—The section stating the ultimate purpose of redemption (“that I might be His own, etc.”) treats solely of God’s, of Christ’s, purpose, and should not be deprived of its objective character. It remains for the Third Article to picture the execution of this purpose in detail; there the very question is answered, how the individual is brought to the enjoyment of all the blessings which Christ has won for him by His redemptive work. Instead of anticipating the Third Article, the object must be to let the pupil envisage the blessedness of the change purposed by Christ (formerly the devil’s property, now Christ’s; formerly a lost and condemned sinner, now righteous, innocent, and blessed). This blessedness is to begin here below (“in His kingdom” describes the life of the faithful in the kingdom of Christ as it is here on earth, i.e., the Church; cf. W. E. xxx 1, pp. 80 f.; 186); it includes, however, the life beyond (cf. “in everlasting righteousness”, a phrase equivalent to “state of justification”, cf. W. E. xxx 1, p. 187). It is to be hoped that no catechist will henceforth dissect this magnificent period and endeavor to correlate its constituent parts with this or that part of the text of the Apostolic Creed. He would deprive “the golden gem” of the fairest part of its beauty and dim its power (cf. pp. 103–105).

He who has come to the Third Article by way of the First and Second, will expect Luther to select some one word as the nucleus for the whole explanation. The heading, “Of Sanctification”, strengthens this expectation. And Luther actually begins this article in his Large Catechism: “This article I cannot relate better than to Sanctification, that through the same the Holy Ghost, with His
\quellseite{382}
office is declared and depicted, namely, that He makes holy. Therefore we must take our stand upon the word Holy Ghost, because it is so precise and comprehensive that we cannot find another” (Tr. p. 687 § 35). But how does Luther now correlate the other parts of this article with the main one? The Large Catechism does not leave us without an answer to this question. We read presently: “But how is such sanctifying done? Answer. Just as the Son obtains dominion, whereby He wins us, through His birth, death, resurrection, etc., so also the Holy Ghost effects our sanctification by the following parts, namely by the communion of saints or the Christian Church, the forgiveness of sins, the resurrection of the body, and the life everlasting; that is, He first leads us into His holy congregation, and places us in the bosom of the Church, whereby He preaches to us and brings us to Christ” (Tr. p. 687 § 37). Subsequently he repeats this outline of the article in practically identical words and says of the exposition given on basis of this outline: “Let this suffice concerning the sum of this article.” But he himself does not let this “sum” suffice. “Rather, because the parts which are here enumerated are not quite clear to the simple, we shall run over them also”, he says (Tr. p. 689 § 46), and then devotes a special paragraph each to Church, Communion of saints, Forgiveness of sins, Resurrection of the body, and Life everlasting. Not so in the Small Catechism. Here he consistently carries out his principle that “everything must take its stand upon the word Holy Ghost.” He touches, indeed, upon the other parts, but he neither coordinates them with the main part nor explains them as \emph{disiecta membra}; he subordinates them all to the Holy Spirit, and confesses his faith in Him and His work alone. At a previous place we have emphasized the fact that these attributes of unity and compactness constitute the principal excellence of his explanation of the Third Article (p. 105).

It is regrettable that this advantage has been recognized and turned to account by but few commentators. Quite a number have adopted the mediæval twelvefold division, while others, who succeeded in preserving the unity of the first two articles, have dissected the Third into five, more or less discrete parts. This has been the case down to modern times when the assertion was made, or rather renewed, that Luther had accurately followed in his explanation the several parts of the text. There is some truth in this; but when the subjects, “Church”, “Forgiveness” etc., are treated independently, the fact is overlooked
\quellseite{383}
that Luther does not see them at all in that light, but only as subordinate to the one subject—the Holy Spirit. Even Kaftan, although more intent than most catechists upon maintaining the connection between the Holy Spirit as subject and the other parts, nor unmindful of the tense employed by Luther (“has called me . . . kept me in the true faith”), does not do justice to the matter. He says in regard to the division: “Besides our faith in the Holy Ghost we also confess in the Third Article that there is a holy Christian Church, which is His workmanship and in which He Himself dwells and works” (Kaftan-Horine, p. 131). Now let us apply to this statement the touch-stone of Luther’s exposition: according to that, we confess the Church not as something “alongside the Holy Spirit”, but we confess of the same Holy Spirit, in regard to whom we have confessed something in the first section, something additional, namely, that He calls the whole Christian Church on earth, etc. We have been of the conviction for more than thirty years that Luther’s explanation best comes into its own when outlined as follows: (1) What the Holy Spirit has done in the past; (2) what He is doing in the present; (3) what He still will do in the future. When the article is thus outlined, everything, precisely as in Luther’s explanation, is \emph{dominated by one subject}—the Holy Spirit; the other parts of the article receive the same treatment that Luther accorded them; the perfect tense in the first section which is so largely but none the less erroneously exchanged for a present tense, retains its distinct tense value; and the whole article, as we can assert from experience, becomes clear and intelligible to the children. Nor is it necessary to repeat here that, where the Third Article is thus taught, the fundamental fact that the salvation of the Christian, from beginning to end, rests not on himself but on Christ and His Holy Spirit, will effectively be impressed upon the pupil’s consciousness (pp. 105 f.). Kaftan, briefly, repudiates this outline (p. 229). But the Large Catechism shows that it was not at all remote from Luther’s mind; for he there connects the subjects of Church and forgiveness of sins with the present, and those of the resurrection and eternal life with the future, exactly as he does in the Small Catechism. Here are his words: “Therefore we believe in Him who through the Word daily brings us into the fellowship of this Christian Church, and through the same Word and the forgiveness of sins bestows, increases, and strengthens faith, in order that when He has accomplished it all, and we abide therein, and die to the world and to all evil, He may finally make us perfectly and forever holy; which now we expect
\quellseite{384}
in faith through the Word” (W. E. xxx 1, p. 191; Tr. p. 695 § 62). With even greater clearness, these thoughts are found in the Catechetical Sermon of December 10, 1528: “He is here called the Holy Spirit. Why? Because He sanctifies. And I believe in the Holy Spirit for the reason that He has sanctified me and sanctifies me to this day . . . And the Holy Spirit will sanctify you through the resurrection of the body . . . then the Holy Spirit will complete His work and my sanctification will be complete . . . The Third Article means that I believe in the Holy Spirit, that is to say, that the Holy Spirit of God will and does sanctify me . . . He has already begun to sanctify me; when I shall have died, He will complete this sanctification through these two parts, the resurrection of the body and eternal life . . . The other parts signify the means, or method, through which He sanctifies me, for the Holy Spirit does not sanctify you outside the Church. That is the reason that the Christian Church in which all His gifts are found, is placed immediately after the Holy Spirit. Through this Church He proclaims His Message, calls you, makes Christ known to you, and kindles faith in you that, through the Word and the Sacraments, you may become free from sins; and thus you are truly free on earth. If you die and remain in the Church, He will raise you from the dead and will sanctify you wholly . . . And the Holy Spirit sanctifies me through the Word and Sacraments, which are in the Church, and will sanctify us perfectly on the last day” (W. E. xxx 1, pp. 91. 93 f.). We refer to our Explanation as proof that when the Third Article is thus skeletonized, Luther’s explanation comes into its own better than with any other outline so far suggested.

The term sanctification occurs in the Third Article in both a narrower and wider sense: in the wider sense in the heading, in the narrower in the sections “sanctified in the true faith” and “sanctifies the whole Christian Church”. Sanctify in the wider sense means to Luther “to lead to Christ for time and eternity”; compare the above quotations. In the narrower sense, therefore, no other meaning is possible than “to bring one to faith and thus to lead him into a state of righteousness.” The latter view has been impugned by those who understand by “sanctified” subjective sanctification—the victorious life-long conflict with sin, or at least, the union of both—justification and the daily subjective sanctification. That sanctification can be taken in a subjective sense, and has been taken in that sense also by Luther, is perfectly true; but that
\quellseite{385}
does not prove that he used it in that sense in this connection. That Luther has rather used it at this place in an objective sense, is proved in the first place by the fact that the objective sense, in his writings as well as in the New Testament, is the one which predominates (cf. Koestlin, \emph{Luther’s Theologie} ii, pp. 178 f. “sanctification = justification”); in the second place, by the fact that it is so used in the Large Catechism, e.g., “Therefore sanctifying is nothing else than bringing us to Christ to receive this good [Christ’s redemption] to which we could not attain of ourselves” (W. E. xxx 1, p. 188; Tr. p. 689 § 39).—With this question another one is intimately connected, namely, how the four verbs “\\emph{call, enlighten, sanctify, and preserve}” are to be understood. The view has found favor that they denote four stages of the \emph{ordo salutis} and that they must be interpreted in strictest conformity to our dogmatical terminology (some have found even a chronological order in these words!). Thus e.g., Tetelbach: “How many acts of the Holy Spirit are here enumerated? Four. First, He calls us to a diligent hearing of the Word of God, and makes and effects in us the beginning of conversion. Which is the other? He enlightens by the gospel our dark and blind hearts through the light of faith so that we, erstwhile unbelievers, become believers. Which is the third? He sanctifies us in the true faith so that a new life proceeds from faith, a life of new obedience, of love, joy, peace, patience, kindness, charity, meekness, and chastity. Which is the fourth? He preserves us constantly in the truth which we have recognized, in the true faith, and a blessed confession to the end. Accordingly, beginning, middle, and end of our salvation is the work of God the Holy Spirit” (Reu i 1, p. 689 f.). Formerly we had the same view; but by paying more attention to biblical usage regarding this point (see Reu, \emph{Heilsordnung}, 1914), and by renewed study of Luther’s preliminary labors, we have come to the conclusion that the first three of the verbs enumerated are not at all intended to describe three distinct, successive stages of the Holy Spirit’s work, but rather the one saving work from three different aspects, in order effectively to reveal its magnitude. When we were far from Christ, the Holy Spirit led us to Christ (call—effectual, successful call, not invitation); when darkness was in our hearts, He enlightened us with His gifts, i.e., He kindled a bright light in our hearts by the knowledge which He imparted; when we were servants of sin, He transplanted us into a state of holiness and righteousness, through faith in the forgiveness obtained in Christ, a faith wrought
\quellseite{386}
by Himself. In the first clause the emphasis is on the Word through which the act of sanctification takes place; in the second, on the knowledge conveyed by the Holy Spirit; in the third, on saving faith as wrought by the Holy Spirit. Looking back, the Christian cannot but confess: “If the Holy Spirit had not led me to Christ through the gospel, I should still be far from Him; if He had not enlightened my heart with His gift of knowledge, I should still sit in darkness; if He had not kindled faith in Christ in my heart and thereby enabled me to obtain forgiveness of sins and to enter into a state of holiness and righteousness, I should still be in my sin”. But in making this confession, the Christian does not enumerate three successive acts; he merely means to give expression, from every possible point of view, to the one blessed experience of salvation. This was understood by A. Siber already who paraphrased this section of the Enchiridion in his \emph{Sabbatum puerile} (1575) as follows: “\emph{Qui quidem, lege conclusum et occisum, sermone reconciliationis, quae est vox Evangelii, in vitam me vocavit} (note well: effectual call!), \emph{sua gratia collustravit, fidei sanctitate insigniverit}” (Reu i 2, 2, p. 570). When discussing this section with its past tense, the catechist will have to impress upon the children the fact that the Holy Spirit has worked upon them through Baptism and instruction in the Word; but he dare not forget what was said by us on p. 314: let him speak from the standpoint of his own life or that of well-known men of God lest guile or lip-service insinuate themselves into the teaching, especially in view of the fact that he also confesses: “and preserved me in the true faith.”

Finally, a brief remark relative to the much mooted section, “enlightened me with His gifts”. He who thinks that the fundamental features, or the various successive steps, of the order of salvation are here indicated, and accordingly mistakes the “call” as a mere invitation to come to Christ, will surely find himself tempted to understand by the “gifts” with which he has been enlightened, the Law and the Gospel—a view that, for a long time, constituted the exclusive catechetical tradition. While already the sequence thus arising (Gospel, Law, Gospel) should make one suspicious, that it in itself would not be an invincible obstacle, at least not from the standpoint of dogmatics. Even more significant would be the change of prepositions: for we should expect durch, by, a second time instead of mit, with. What is more important is the fact that nothing in Luther justifies such an interpretation, at least nothing in his preliminary labors on the Catechism. Albrecht
\quellseite{387}
has well summed up what is to be said in this respect: “In the sermon of December 10 (W. E. xxx 1, p. 94), the \emph{dona spiritus sancti} evidently mean the objective spiritual gifts or official endowments in the sense of 1 Cor. 12:4 ff. and Eph. 4:11 f. The phrase “all His gifts” in the Large Catechism will have to be understood in the same sense (W. E. xxx 1, p. 192), probably also that found on page 190: “with various gifts.” But the clause on page 192, line 15 f.: “and besides they are not enlightened and favored by any gifts of the Holy Spirit”, when studied in the light of its context, leads to the conclusion that here the subjective operations of the Holy Spirit are meant. Luther’s point of view may have been the widely spread medieval doctrine of the seven gifts of the Holy Spirit, according to Isa. 11:2 (that Luther retained this view is proved from his pentecostal hymn, “Come, O Creator God . . . . Thou art with powers sevenfold the finger on the hand divine); or it may have been certain biblical passages, such as Rom. 8:1 ff.; Gal. 5:22; or possibly a suggestion such as that contained in the German Patrem: \emph{Der aller Bloeden Troester heisst und mit Gaben zieret schoene}” (W. E. xxx 1, p. 367 f.). R. Hadamarius writes (1552): “\emph{Quae dona sunt Spiritus sancti? Vera sapientia, intellectus, consilium, fortitudo, scientia, pietas, timor Dei}” (Reu i 2, 2, p. 1030). Bischoff (1599): “Thereupon He enlightens our hearts with the true knowledge of God and true faith in Jesus Christ, whereby we obtain the forgiveness of sin” (Reu i 2, 2, p. 231). The Joachimstal Catechism (1574): “Who enlightens our blind hearts? Alone the Holy Spirit, the very love and fire in God; He is poured out into our hearts and regenerates us, enkindling a new light in us and leading us to the knowledge of God the Father and His Son, Matt. 16:17; 1 Cor. 12:3” (Reu i 2, 2, p. 609). The interpretation of the gifts of the Spirit in an objective sense did not appear in the expositions of the sixteenth century. In itself it would be quite correct to say that the Holy Spirit enlightens us by arousing apostles, prophets, evangelists, etc. A warrant for this view might be found in the December sermons and the stress laid upon the ministerial office in the Large Catechism; but this interpretation, too, would call for the preposition durch, by, instead of mit, with. The subjective gifts must be meant. The Spirit has enlightened and favored, graced and adorned me with knowledge, wisdom, comfort, power, and every variety of virtue (cf. Large Catechism, W. E. xxx 1, pp. 190, 192; Tr. p. 691 § 51, 697 § 66); or, as more in harmony with the context in the Small Catechism (consider, e.g., the contrast to
\quellseite{388}
“not by my own reason”): He has enlightened me within by giving me the true knowledge of Christ and my sin: of my sin, by teaching me that it poisons my life and accuses me before God; of Christ, by showing me that He is my Saviour and my Lord, who has redeemed me from sin, and whose own I am to be henceforth and forever.

Whatever needs to be said concerning the \emph{Holy Trinity}, will readily suggest itself when the teacher at the conclusion of his discussion sums up the contents of the three Articles. “We could never attain to the knowledge of the grace and favor of the Father except through the Lord Christ, who is a mirror of the paternal heart, outside of whom we see nothing but an angry and terrible Judge. But of Christ we could know nothing either, unless it had been revealed by the Holy Ghost”, says Luther in the Large Catechism (Tr. p. 695 § 65).—In regard to the other Chief Parts we refer to Kaftan and our \emph{Explanation}, cf. also M. Reu, \emph{A New Translation} etc.

Whether it be advisable to supply the child with an elaborate explanation of the Small Catechism as an aid in instruction is a moot point. Our local educational conditions would strongly speak in its favor. In any event, the printed explanation is to be considered only as \emph{the result or thesis of the previous catechesis}; it is given to the child in order that he may have a fixed lesson for home study. For the catechesis proper pupil as well as teacher knows only the text of Luther’s Catechism. Whether the form of such catechetical explanation is erotematic or thetic, is in itself immaterial. This much, however, can be said in favor of the thetic form: (1) the catechist will find himself less \emph{shackled} by the employment of the thetic form than the other; (2) the child is more easily trained for independent thinking where the thetic method obtains for the reason that, in reviewing, he is compelled to seek the answer for himself; (3) the thetic form supplies the thought as a whole whereas the other method usually supplies only half the thought, since the child is very much inclined to memorize only the answer and to put the question aside; and (4) the thetic method leads to better mastery of the Catechism, for here the text of the Catechism can be plainly traced by the eye, as it winds through all the explanation like a string of pearls.


\quellseite{389}
\chapter[§ 27. Holy Scripture]{§ 27. HOLY SCRIPTURE}

*G. v. Zezschwitz ii 2: \emph{Der bibl. Unterricht in der Volksschule}, pp. 188–220. K. Buchrucker, pp. 154–158. E. Sachsse, pp. 334–336. *F. Zange, pp. 73–142. A. Eckert, pp. 33–35; 55–61. O. Baumgarten, pp. 74–80; 87–88. J. Gottschick, pp. 126–141. *J. Berndt, pp. 36–55. *R. Kabisch, pp. 129–156. E. Chr. Achelis, pp. 266–283. A. Rude, pp. 90–99; 126–129. J. Steinbeck, pp. 193–200. E. Kautzsch, \emph{Bibelwissenschaft und Religionsunterricht}, 1900. T. Raymont, \emph{The Use of the Bible in the Education of the Young}, 1911. H. E. Weber, \emph{Historisch-kritische Schriftforschung und Bibelglaube}, ²1914. *Geo. Hodges, \emph{The Training of Children in Religion}, 1917, pp. 167–217. *M. Reu, \emph{Quellen zur Geschichte des kirchlichen Unterrichts, ii: Der Bibelunterricht}, 1906. On Bible Passages: G. W. W. Thilo, \emph{Der Bibelspruch im Dienst des Religionsunterrichts}, 1846. M. Reu, \emph{Quellen etc.}, 1906. Eckhart, \emph{Kurze Erklaerung der wichtigsten Bibelsprueche}, 1887. Peschel, \emph{Ausfuehrliche Erklaerung der wichtigsten Bibelsprueche fuer den Katechismusunterricht}, 1887. L. Schultze, \emph{Katechetische Bausteine}, ⁵1891. *K. Buchrucker, \emph{Der Schriftbeweis im Katechismus}, 1893. Hauffe, \emph{120 Sprueche schulmaessig erklaert}, 1896. *H. Toegel, \emph{Der konkrete Hintergrund zu den 150 Kernspruechen d. relig. Memorierstoffes}, ²1904. H. Toegel, \emph{Die Behandlung der Sprueche} in \emph{Paedag. Studien}, 1904. Th. Meinhold, \emph{Die biblische Grundlage f. d. Katechismusunterricht}. 1907. Zeissig and Fritzsche, \emph{Prakt. Volksschulpaedagogik} (containing seven catecheses on Bible Passages by Toegel), 1907. W. W. Stellhorn, \emph{Schriftbeweis im luth. Katechismus}, 1912. *L. Wessel, \emph{The Proof Texts of the Catechism with a Practical Commentary} in \emph{Theological Quarterly}, 1917 f., book edition 1923. On Bible Reading, Introduction to the Bible, School Bibles and Bible Readers: Boettcher, \emph{Handreichung f. d. Bibellesestunde}, 1894. O. Zuck, \emph{Das Bibellesen im Anschluss an Bibl. Geschichte und Katechismus}, 1895. Melinat, \emph{Das Bibellesen im Volksschulunterricht}, 1898. Habermas, \emph{Handbuch d. Bibellesens u. der Bibelkunde}, 1898. H. and F. Falcke, \emph{Bibellesen}, 1899. J. H. A. Fricke, \emph{Bibelkunde}, i ²1904; ii 1898. J. A. W. Haas, \emph{Bible Literature}, 1903. General Council Graded System, \emph{Bible Outlines, Scripture Quarterly, Lesson Commentary}. H. C. Alleman, \emph{The Book and the Message}, 1914. *M. Reu, \emph{Wartburg Lesson Helps}, course iii: \emph{The Book of Books}, 1926; \emph{The Book of Life}, 2 vols., 1917 ff. John Schaller, \emph{The Book}
\quellseite{390}
\emph{of Books}, 1918. F. Sheatsley, \emph{A Guide to the Study of the Bible}, 1918. R. Hofmann, \emph{Die Schulbibel}, 5¹896. Martin, \emph{Die Schulbibelfrage}, 1888. Baehlnisch, \emph{Ist eine Schulbibel notwendig und wie muss sie beschaffen sein?} 1892. J. Schlier, \emph{Bibel, Schulbibel, Biblisches Lesebuch} in \emph{Neue kirchl. Zeitschr.}, 1904. Behr, \emph{Die Schulbibelfrage}, 1895. Evers, \emph{Die Schulbibelfrage}, 1895. Mueller, \emph{Bibel oder Schulbibel}, 1896. Dix, \emph{Neuere Geschichte der Schulbibel}, 1898. R. Locke, \emph{Vollbibel, Schulbibel, Bibl. Lesebuch}, 1898. Bremische Bibelgesellschaft, \emph{Biblisches Lesebuch oder Schulbibel}, 1894. Schaefer and Krebs, \emph{Bibl. Lesebuch f. d. Schulgebrauch}, 1896. Voelker and Strack, \emph{Biblisches Lesebuch}, 1890. 1898. \emph{Biblisches Lesebuch f. evang. Schulen}, Stuttgart, 1901. Frischi, Schremmer and Holzinger, \emph{Bibl. Lesebuch}, 1909. *Ostermai, Toegel and Neuberg, \emph{Bibl. Lesebuch mit Bildschmuck}, 1911. *Thisselton Mark, \emph{The Young People's Bible}, 2 vols., 1910. *Ella Broadus Robertson, \emph{The Heart of the Bible}, 1911. Kent and Foster, \emph{The Shorter Bible}, 1923. H. Spanuth, \emph{Die Propheten d. A. Bundes. Lebensbilder u. Entwurfe z. unterrichtl. Behandlg.}, 1903. P. Staude, \emph{Der Prophet Jeremiah}; \emph{Praeparationen in darstellender Form}, 1906. Witzmann, \emph{Die unterrichtl. Behandlung d. Psalmen}, 1906. Werkmeister, \emph{20 ausgew. Psalmen; Praeparationen}, 1908. E. Greenwald, \emph{Questions on the Gospels for the Church Year}, 1873. Id., \emph{Questions on the Epistles for the Church Year}, 1874. Sperber, \emph{Erklaerung der Sonn- und Festtagsevangelien}, 31888. R. Kabisch, \emph{Die Episteln u. Evangelien f. Volksschullehrer}, 1897. Lohof, \emph{Kurze Handreichung z. Behandlg. d. Perikopen d. Kirchenjahrs}, 1902.

Biblical History and the Catechism are based exclusively on the Scriptures and in part literally reproduce the Bible text. This might seem to render superfluous further, specific instruction in Holy Scripture. The omission of the Bible from the religious curriculum, however, would not correspond to the supreme importance which Evangelicals attach to that book. As the record and the permanent re-presentation of the divine revelation; as the fountain and norm of all saving knowledge; as the voice of true authority and the book of life for the mature congregation, from which comfort, counsel, incentive, and power are to be drawn every day, the Bible requires a special method of treatment: and instruction in Biblical History and Catechism would be deficient
\quellseite{391}
at a vital point, did they fail to lead the pupil back to Scripture at every step, and thereby prepare him for an independent use of it. If the members of the mature congregation of the present day do not use the Bible with greater diligence, the cause is, in great part, that in their youth they have not been properly introduced to this precious book.

One feature of the instruction of the youth in Holy Scripture is the literal \emph{memorizing of some of the principal parts of Scripture}. Of late, opposition against this practise has arisen, on the plea that it degrades the Bible to the level of a book of oracles; but the advocates of such a view have failed to comprehend one of the peculiar features of Scripture, viz., that fundamental truths have so crystallized in some of its individual passages that the whole Christ, with His holy love and divine grace, shines forth from them. Nor has due account been taken of the need of the ordinary Christian, who lacks the faculty of surveying large complexes, and who is unable so to retain them that they are at his disposal as soon as needed. It is in harmony with his manner of learning a subject and retaining it for future direction and solace to entrust to his heart and memory a treasure of passages in which this or that truth is concentrated as light is in the sun. The history of the Christian Church bears abundant testimony to the fact that just such individual passages have been the stay and comfort of many.

Bible passages are most effectively memorized when associated with instruction in Biblical History or the Catechism; for the Lutheran catechist feels the need of summarizing the fundamental thoughts of the individual Bible narratives in words of Holy Scripture and, likewise, of proving the doctrinal statements of the Catechism by Scripture. The last-named province suggests the necessity for a fund of \emph{dicta probantia}, or proof passages. In the selection of these, the
\quellseite{392}
main requisite is that they are clear and readily intelligible, brief and pithy, trenchant and apt, and also easily memorized and retained. Above all, the connection in which they are found in Scripture must tally with that in which they are employed in catechetical instruction: only thus can they possess the force of proof passages. In connection with some of the chief passages, the catechist will do well to enter more fully upon the biblical context. If a number of passages is used in support of a certain doctrine, they should be so arranged as to correspond to the progress of revelation, e.g., when the Bible proof for the divine sonship of Christ is furnished. Particular attention should be devoted to such passages as embody the consummation and fruitage of whole periods of development. But the criterion for the selection of Bible passages is not merely their evidential value, but also their vital content in the form of comfort and admonition. Words that awaken life and quicken life, words of majesty that engrave themselves, as with style of diamond, upon the heart and conscience of the child, to become a real treasure for all after-life, a stand-by in breath and in death—such words require particular attention. Especially the Psalms and the Prophets, with their abundance of pithy maxims, should freely be drawn upon. It is of the utmost importance that the teacher should guard against the mistake of having too many Bible passages memorized. Where there is no regular parish school, two hundred and fifty passages should be the limit; these should be reviewed with such frequency that they become the children’s never-failing treasure. Only thus time is found to explain them—a duty that must never be set aside, and, also, to have the more gifted children in the course of time memorize some larger sections of Scripture, e.g., Ps. 23; 51; 90; 91; Isa. 9:1–7; 40:26–31; 53:1–12; portions of the Sermon on the Mount and of the last sayings of Christ (such as Matt. 25:31–46);
\quellseite{393}
Phil. 2: 5–11; Rom. 3: 21–28; 5: 1–11; 8: 31–39; 1 Cor. 13; Rev. 3: 17–20.

Another element of instruction in the Holy Scriptures has been, since the sixteenth century, the reading and explanation of the pericopes (usually on Saturdays, to facilitate the understanding of the sermon on the day following). Since instruction in Biblical History has become general, the need of this feature is not so urgent as in the past. Where such instruction is given, the children should look up the pericopes in their Bible, not in their Church Book; the portions hard to understand should be carefully explained; in all cases the main thoughts should be made to stand out; and the whole pericope should be explained to the children in an edifying manner, and impressed upon their minds, in keeping with the respective level of knowledge to which they have attained. But still other features must be observed in the endeavor to introduce the young to Scripture: \emph{Regular reading of the Bible} must take place. This will be most effective where special periods are appointed for it, so that larger sections can be read and studied connectedly. Where that proves impracticable, Bible reading should at least regularly accompany instruction in Biblical History and in the Catechism. It is recommended that, when the creation narrative or the First Article is under consideration, such portions be read as Ps. 104; 121; 91; Job 38:1–40:5; Isa. 40:12 ff.; Ps. 8; 1 Kings 17:1–16; Acts 12:1–24; Gen. 32:10–13; Matt. 6:25–34. In connection with the Conclusion of the Commandments the class might read the narrative of the flood, of the destruction of Sodom and Gomorrah, of the destruction of Jerusalem, of the exile of the children of Israel, but also the narrative of Joseph. In connection with the Second Article the following portions would be appropriate: Luke 2:1–14; John 1:1–18; Isa. 53; Ps. 110; with the Third Article Joel 2:28–32; Acts 2:1–47; John 16:1–15. The Third Chief
\quellseite{394}
Part would receive light from John 17:1–26; the Fourth from 2 Kings 5:1–19; John 3; Acts 8:26–40; Rom. 6:3–11. The Fifth Chief Part suggests the reading of John 6; Rev. 3:14–22; 1 Cor. 11:23–32; and the part dealing with Confession Ps. 51; 32. Where there is a good parish school, time will be afforded in connection with instruction in the Catechism to introduce the pupils to Holy Scripture in the manner here shown; where that opportunity does not exist, larger, connected sections of the Bible ought to be read at least during the last weeks of catechetical instruction. Regular reading and study of the Bible should take place in the Bible class of Sunday Schools and in the devotional meetings of the Young People’s Society.

With the reading of the Bible the more important elements of \emph{Introduction to the Bible} as a whole are readily connected. In Biblical History the sequence of the historical books and their main content are learned; the knowledge of the poetic and prophetic books especially must be imparted by means of Bible reading. The content of the Gospels is imparted in a similar way while the order of the Pauline Epistles can be discussed best in connection with the missionary journeys. Thus the most important phases of introduction to the Bible can easily be taught in connection with the rest of religious instruction. The order of biblical books should not be memorized disconnectedly, but in the connection above suggested. The use of rhymed verses for this purpose can be recommended only upon the condition that better ones than those hitherto used shall have been formed. The children must also necessarily be drilled in finding any one passage at a moment’s notice, they must know something about the technical make-up of the Bible, and they must learn the conventional abbreviations of the biblical books. The proper place for exercises of this sort will be catechetical instruction. As the young begin to mature, they should be
\quellseite{395}
informed concerning the more important sections of the chief books. For the initial stage of Bible introduction the appendix to our \emph{Explanation of Luther’s Catechism} might be acceptable (pp. 147–151). For the confirmed the third course of the \emph{Wartburg Lesson Helps} (\emph{Book of Books}; for the more advanced: \emph{Book of Life}) is intended. The whole process of introduction should glow with respect for the Book and joy in it; and the teacher should at all times aim at equipping the young for personal use of the Bible. Otherwise we need not entertain any hopes of training a generation that delights and takes root in the Bible. It is also recommended to link all truths gleaned from the Scriptures directly with the corresponding portions of the Catechism, because it will help to bring clearness into the child’s religious concept world, cf. § 35.

But is it at all proper to put the Bible into the hands of the children? In view of the frank and unvarnished presentation of sexual matters and sins, this has been looked upon as morally doubtful; for this reason so-called \emph{School-Bibles and Biblical Readers} have been produced. For various reasons, a Biblical Reader is quite desirable (Th. Mark and E. B. Robertson have given us such readers); but even when using the complete Bible, we can, in a measure at least, obviate the difficulty if (1) there is no continuous, but only select, reading, and (2) if we tell our children as soon as the Bible has been put in their hands, that this book pictures man as he is: that even the most heinous sins are there called by name; but that this is not done for our entertainment, but for our humiliation and warning, which may be seen from the fact that, generally, the penalty is pictured immediately after the sin; and that, when we come upon such passages we should read them with holy grief and the sincere prayer that the Father in heaven preserve us from similar transgressions. A determined fight against the carnal appetite of the young should be the chief aim.

\quellseite{396}
\chapter{§ 28. THE CHURCH BOOK OR HYMNAL}

K. Buchrucker, pp. 158–162. E. Sachsse, pp. 366–368. J. H. Schueren, pp. 105–125. J. Berndt, pp. 57–59. R. Kabisch, pp. 152–157; 247–249. A. Rude, pp. 117–122. J. Steinbeck, pp. 218–222. *W. Thilo, \emph{Das geistl. Lied i. d. ev. Volksschule Deutschlands}, ²1865. *E. Sperber, \emph{Gesch. d. Behandlung des Kirchenlieds} in Schumann u. Sperber, \emph{Gesch. d. Religionsunterrichts}, 1890. *K. Buchrucker, \emph{Der Gesangbuchunterricht}, ²1887. Schumacher, \emph{Lehrbeispiele z. Behandlung d. Kirchenliedes i. d. Volksschule}, 1885. Gattermann, \emph{50 evang. Kirchenlieder}, ³1900. Koeller, \emph{Das ev. Kirchenlied i. d. Volksschule}, 1904. O. Zuck, \emph{Das Kirchenlied}, ³1906. Schultz and Triebel, \emph{Die gebräuchlichsten Lieder d. ev. Kirche}, ¹⁷1907. Lehmensick, \emph{Kernlieder der Kirche in Stimmungsbildern}, 1907. Schlegel, \emph{Praeparationen fuer Kirchenlieder u. Psalmen}, 1908. *Achenbach, \emph{Behandlung d. Kirchenlieds auf historischer Grundlage. Lehrbeispiele nach psychologischer Methode}, ⁴1910. *F. Niebergall, \emph{Bibl. Geschichte, Katechismus, Gesangbuch}, 1910. Nic. Smith, \emph{Hymns Historically Famous}, 1901. Th. Brown and H. Butterworth, \emph{The Story of the Hymns and Tunes}, 1906. *W. L. Hunton, \emph{Favorite Hymns. Stories of the Origin, Authorship, and Use of Hymns We Love}, 1917.

Since training for participation in congregational worship is one of the Church’s duties toward the rising generation, and since the Church Book or Hymnal is the indispensable manual for worship, it will also be necessary to make the young acquainted and familiar with this book. Already in the sixteenth century when, as a rule, Church Books were not placed in the hands either of the children or adults, the catechism, for this very reason, contained as an appendix some hymns, prayers, and other parts of the liturgy, even the forms for Baptism and marriage.

Public worship adjusts itself to the seasons of the church year; and this idea of the church year is worked out in the Hymnal both in the liturgical section and in that containing the hymns. Our youth must therefore be introduced to the
\quellseite{397}
meaning of the church year. This must be done very carefully in our country where the Lutheran Church is surrounded by so many churches either altogether ignorant of the idea of the church year or acquainted with mere fragments of it. As the outstanding redemptive acts of God pass in review before our eyes in church year after church year, they invite us time and again to devout contemplation and to earnest application; neglecting the church year we would deprive ourselves of a most valuable element of religious instruction. Special lessons on the church year are, perhaps, not necessary; the information can well be imparted incidentally (in connection with the reading of the pericopes where that is still done, or in Biblical History if the stories have been selected with a view to the church year); in any event, where it has not been done previously, it should be given in a concise and coherent form in connection with the instruction of confirmands. Of course, care must be taken to avoid unnecessary historic details and popular but unwarranted theories concerning the selection of the pericopes. What has been said in our Explanation of Luther's Catechism (pp. 151–155) might be sufficient for the beginning, later to be supplemented in the Young People's Society.

The next subject requiring the attention of the catechist are the \emph{hymns}. They are the sublimate of the life of the Church; they reveal the experiences of sin and grace through which other Christians before us have passed; they witness to the abiding power of faith that has strengthened and comforted the Christian saints; they are the fairest flowers in the garden of evangelical piety. This is the very factor constituting their educational value. When the young perceive how these men of God have been sustained in affliction by their faith and by grace divine, how they have found shelter and comfort and cheer in God's omnipotent Word, their own respect for a life of faith will increase and they
\quellseite{398}
will gain greater confidence in the fountains of Christian life from which these hymns have streamed; and they will tremble at the awful reality of sin when they see how it has pressed sighs of anguish from the lips of the Christian poets. Another reason why hymns should be employed in religious instruction is that they sink more deeply into the soul when gladly and expressively sung than the spoken word; not rarely it occurs that when all else is forgotten, a hymn awakens in the soul and unfolds saving force. If there is to be fresh, vigorous congregational singing—and it is the predominance of congregational instead of choir singing that corresponds best to the Reformation principle—the young must be drilled with diligence in the melodies of those hymns; and they should be taught to sing them in their \emph{rhythmical setting} (which is finding favor owing to the efforts of Tucher, Layritz, and especially of Zahn), not in their perverted forms dating from the time of rationalism. Compare, e.g., the musical setting of the chorals in the \emph{English Hymnal} of the Missouri Synod, in the \emph{Hymnal of the Common Service Book}, in the \emph{Norwegian Lutheran Hymnary}, and in the \emph{Wartburg Hymnal for Church, School, and Home} (ed. by O. Hardwig).

In \emph{Gottes Wort im Feld und daheim}, O. Eberhard says: “The church hymn is not, like its brother, the folk song, the property of merely a certain period of life, that of youth; it belongs to the whole congregation, to all periods of life. It is learned by heart by the children, criticized by those in a state of inner ferment, explained by the lessons of life, learned anew and comprehended by the adults, fathomed more deeply through progressive experience, transfigured by age, and tested by death. The hymn is your companion from the cradle to the bier, the expression of your every need. Its deep, central notes not only thunder down from the organ loft; they also rise from the pew. And when the time for singing is past, it trembles upon the lip as a prayer; when the word loses its power in depths of woe, the old comforters once more begin to gleam; and in the darkest night of suffering they sparkle as inextinguishable stars” (\emph{Kirch. Zeitschr.} 1915. pp. 278 ff.). And Sneath, Hodges, and
\quellseite{399}
Tweedy, express themselves as follows: “One of the best gifts which parents can give to their children is constant familiarity, intelligent appreciation, and hearty and sincere use of the world's great hymns. Like all good things it will demand time and effort; but the expenditure will be more than compensated by the gain” (\emph{The Religious Training in the School and Home}, p. 163).

F. W. Herzberger says:

\begin{quote}
“Our Lutheran Church is pre-eminently the Singing Church of Evangelical Christendom. No other Church can rival her in the rich, soulful music in which she sings her immortal hymns. Countless other songs and melodies have been composed in their day, delighted their audience for a short while, and then passed into hopeless oblivion. Our majestic Lutheran chorals, however, have survived the wrecks of time, and are still today the delight of all true lovers of sacred music, irrespective of creed or language. “The Lutheran Church,” says Dr. Schaff, the noted Reformed theologian, “draws the fine arts into the service of religion, and has produced a body of hymns and chorals, which, in richness, power, and unction, surpass the hymnology of all other Churches in the world.” The late Alexander Guilmant, a Frenchman and devout Catholic, the unrivaled master of the organ in his day, declared that the Lutheran chorals are the most heart-stirring and inspiring tunes in the whole realm of sacred music. Now what is it that gives to our Lutheran chorals or church tunes their imperishable charm? Knowing their history as we do, we must say that it is the spirit of heroic faith, singing in every note its profound adoration of the merciful and omnipotent God that makes these old Lutheran chorals so universally and solemnly impressive in their character. They are alive with pure and holy devotion. They thrill the very depth of the Christian heart because they are born from the deepest and holiest passion of their inspired singers. With few exceptions, they were composed in the heroic days of the Reformation and the Thirty Years’ War, days that called for heroic courage to believe and confess the truth as it is in Jesus; days that demanded heroic submission to the inscrutable ways of our God and Redeemer. The same spirit of sublime, God-given heroism that inspired the texts of our immortal hymns also inspired their heart-stirring tunes. Hence the tunes are an integral part of our hymns. Deprive our hymns of their historic musical setting, sing them to a newer, modern tune, and you have deprived the rose of the fragrance she alone possesses, you have robbed the nightingale
\quellseite{400}
of her most rapturous note. You may then have a sorry hybrid of a poem and some sort of tune, but nevermore the original, forceful, edifying, compact hymn! For in our Lutheran hymns the text and the tune are welded together as inseparable as body and soul in man. The reason is that one and the same spirit of holy devotion gave birth to the texts as well as the chorals, or tunes, of our Lutheran hymnology. Broadly speaking then, our Lutheran chorals are pre-eminently devotional in character.

“It is different with the hymns and tunes of the eighteenth century. That was the time of decaying orthodoxy, and it witnessed the rise of Pietism in Germany and of Methodism in England. Speaking of English tunes in particular, it is a well-known fact that the Reformed Churches of Great Britain at first possessed no chorals of their own. Some of them (e.g., the Episcopal Church) originally borrowed their sacred tunes and even many hymns from the Lutheran Church of Germany. Others (e.g., the Presbyterians) contented themselves with chanting the Psalms of the Bible. They declared all “man-made” tunes and hymns to be inventions of the devil. When, however, Methodism swept over the British Islands, it produced the two greatest hymn-writers of the English-speaking world, Isaac Watts (1674—1748) and Charles Wesley (1708—1788). They were followed by other hymn-writers, both in England and America, whose songs have been set to original tunes. But what is their character? Like the emotional spirit that fostered them, they are, with a few classical exceptions, shallow, insipid, and lacking in that deep reverence of feeling that solemn harmony of tone which characterizes our old Lutheran chorals. It is true, they call themselves Gospel-hymns, but upon closer inspection you will find that very many of them contain very little Gospel and much less of true choral music. Many of them are unevangelical in text, urging and exhorting the sinner to consecrate himself to God by his own powers. Others are so silly and meaningless that sincere Christians in these churches, among them President Woodrow Wilson, have publicly protested against their further use. In their musical setting, particularly, these sensational Gospel-hymns are but little removed from the degenerate and discordant “rag-time” tunes with which the Salvation Army fills the streets of our large cities at night. How much these decadent church-tunes of a more recent date have served to vitiate the popular taste for sacred music needs no further comment.
\quellseite{401}
“We know that these old Lutheran hymns and tunes are not popular with the English-speaking people. But how can they be? Our English populace does not know them and therefore has still to learn them. And they can be learned by English people just as readily as they are learned by German or Norwegian folks. For the last ten years we are conducting a mission school in the tenement district of St. Louis, and we invite everybody and anybody to convince himself if our old Lutheran tunes cannot be learned by children of almost every nation under the sun! Again we can point to the negroes in our Colored Mission, who are originally neither German nor Norwegian, and who sing our Lutheran hymns with a vim as though they had learned them at their mother's breast. We know it takes time and patience to teach our rising English-speaking generations these noble hymns, but the joy at hearing our English youths finally singing the grand old hymns of the Reformation and glorifying their God and Redeemer in them richly repays you for all the trouble. When we hear of a Lutheran pastor who studiously avoids giving out our historic hymns and chorals in public worship, we cannot help thinking that he is either very vain and chasing after cheap popularity, or that he is very ignorant concerning the nature of true church-music, or that he is reprehensibly indolent and shirks the labor of teaching these grand tunes to his people. We hold it to be one of missionary duties of the Lutheran Church in America to acquaint the American public not only with the saving doctrine of our Church, but also with its sacred hymnology. If we Lutherans fail to do this, if we prefer the light, emotional operatic tunes of the present day to the devout, edifying tunes of our fathers, and thus suffer our historic hymns with their chorals to be forgotten, we are depriving our own posterity of the sweetest choral music this side of heaven. Therefore: ‘Lutheran Tunes for Lutheran Congregational Singing!’”\end{quote}

In order to assure intelligent singing of hymns, the teacher should read them to his class with good expression, and briefly explain them (but his explanation must not destroy the poetic perfume!). He should set forth the main thoughts, sketch the historic situation from which it took its rise, and try to enter sympathetically into the feelings of the poet. Some of the hymns, either in part or as a whole, will spontaneously fit in with other elements of religious
\quellseite{402}
instruction, especially with Biblical History and instruction in the Catechism so that from such connection the necessary light is cast upon them. For instance, in connection with the First Article, hymns such as these would be apposite: “Now thank we all our God”; “Praise to the Lord, the Almighty, the King of creation”; “All nations that on earth do dwell”; “O would, my God, that I could praise Thee”; “I sing to Thee with voice and heart”; “My soul, now bless thy Maker”; “If thou but suffer God to guide thee”; in connection with the Second Article: “All hail the power of Jesus’ name”; “Jesus, Jesus, only Jesus”; “With all my heart I love Thee, Lord”; “My dear Jesus, I’ll not leave”; “O power of love, all else transcending”; “Jesus, Thy blood and righteousness”; “Jesus, priceless treasure”; in connection with Abraham’s call: “Thy way and all sorrows”; in connection with the history of Christ’s passion: “O holy Jesus, how hast Thou offended”; “Christ, the Life of all the living”; “O sacred Head now wounded”; “Five wells I know”, etc.; compare also the following chapter. Finally a number of hymns should be committed to memory; here careful explanation is required, or else there will be confusion worse confounded. Such hymns should be the very flowers of the Church Book.

The following hymns are qualified to serve the purpose: Abide in grace, Lord Jesus; My life is hid in Jesus; Thy way and all thy sorrows; If thou but suffer God to guide thee; Now thank we all our God; Praise to the Lord, the Almighty; Good news from heav’n the angels bring; All praise to Jesus’ hallowed name; Praise God the Lord, ye sons of men; Lamb of God most holy; O sacred Head now wounded; Five wells I know; Jesus lives, no longer now; I know that my Redeemer liveth; O Holy Spirit enter in; O enter, Lord, Thy temple; Who knows how near my life’s expended; Why, my soul, thus trembling ever; God, who madest earth and heaven; Now rest beneath night’s shadows; My dear Jesus, I’ll not leave; Jesus, still lead on; Father, Son and Holy Spirit, I’m baptized in Thy dear name; A mighty fortress is our God; Take Thou my hand,
\quellseite{403}
O Father; Jesus, Lover of my Soul; Rock of Ages, cleft for me; Just as I am, without one plea; My faith looks up to Thee; Rest of the weary; The Lord my faithful Shepherd is; The Church’s one foundation. Compare § 31.

In a good Hymnal are to be found also the most important parts of the liturgy. The first task in this connection will be the introduction of the children to the \emph{structure and meaning of the common service and the communion service}. Let the teacher make clear to the pupils that we here deal with things sacred through use for the purpose of worship for over fifteen centuries; that here the adoration of God is the object in view and that we join the throng of suppliants of every time and clime. Let him make clear everything that requires elucidation, carefully practise the musical parts, arouse enthusiasm for the participation of all members of the congregation in the divine service, and especially in the liturgical part, in contrast to Rome which condemns its laity to silence. How should an active and spontaneous participation in these things be possible, if there has never been any introduction to them? From the communion service it is easy to go back to Confession and its liturgical setting. Especially the instruction of the catechumens with its pastoral character affords occasion to cast the necessary light on these acts, to discuss the salutary character of Confession and the custom of announcement for Communion, the principles of Lutheran worship, of Lutheran practice, and of Lutheran customs and Lutheran usage in general. The development of Lutheran customs and churchly usage is, on the one hand, of so much more importance in this country, as the old traditions of the home Church, in part for good reasons, have fallen into disuse; while, on the other hand, churchly usage is not only a protecting barrier for many but also a feature of healthy life in general. Among the remaining liturgical acts the baptismal ceremony requires special attention if the young at a later day, in connection with the
\quellseite{404}
duties of sponsorship, are to take an intelligent part in it. If, in connection with instruction on Baptism or Confirmation, light is cast upon the function of sponsorship with its duties and privileges, it may be possible once more to infuse life into this old institution.

\quellseite{405}
\chapter{§ 29. SECONDARY MATERIAL FOR INSTRUCTION}

A. Eckert, pp. 49. 164–175. *J. H. Schueren, pp. 97–104. O. Baumgarten, pp. 78; 82; 85–107. J. Gottschick, pp. 159–162. J. Berndt, pp. 56–67; 96. A. Rude, pp. 84–90. A. and F. Falcke, \emph{Kirchengeschichte; Praeparationen}, 1902. K. Just, \emph{Kirchengeschichtl. Unterricht}, 1903. K. Just, \emph{Kirchengeschichtl. Lesebuch}, 1903. P. Speer, \emph{Wie in unsern Schulen die Kirchengeschichte behandelt werden sollte}, 1903. F. Heyn, Kirchengeschichte in Reukauf and Heyn, \emph{Ev. Religionsbuch x 2}, 1907. Reukauf and Heyn, \emph{Kirchengeschichtliches Lesebuch}, 1907. Thraendorf and Meltzer, \emph{Kirchengeschichtliches Lesebuch}, 1907. H. Meltzer, \emph{Skizzen z. Behandlung d. Kirchengesch.}, 1909. *F. Zange, \emph{Kirchengeschichtl. Lesebuch}, 1913. Ch. F. Kent and H. B. Hunting, \emph{Witnesses for Christ}, 1913. H. W. Gates, \emph{Heroes of the Faith}, 1913. H. K. Rowe, \emph{Landmarks in Christian History}, 1914. G. H. Trabert, \emph{Church History for the People}. M. Reu, \emph{Life of Luther for the Christian Home}, 1917. C. P. Harry, \emph{Protest and Progress in the 16th Century}, 1917. *A. T. W. Steinhaeuser, \emph{Luther Primer}, 1917. K. Heilmann, \emph{Der Missionsunterricht nach Theorie und Praxis}, 1895. M. Henning, \emph{Taten Jesu in unsern Tagen}, 1906. Hemprich, \emph{Die Mission i. d. Erziehungsschule}, 1909. *G. Warneck, \emph{Die Mission i. d. Schule}, ¹⁴1911. *Th. Schaefer, \emph{Die Innere Mission i. d. Schule}, 71912. M. Hennig, \emph{Quellenbuch d. Inneren Mission}, 1912. Fred Beard, \emph{Graded Missionary Education in the Church School}. *J. F. Ohl, \emph{The Inner Mission}. *E. Pfeiffer, \emph{Mission Studies}, ²1912. *L. B. Wolf, \emph{Missionary Heroes of the Lutheran Church}. I. Boone, \emph{The Conquering Christ}, 1914. P. A. Nordell, \emph{The Modern Church}, 1914. *E. Singmaster, \emph{The Story of Lutheran Missions}, 1917. O. Koenig, \emph{Die Mission im Katechismusunterricht}, 1913. J. Richter, \emph{Die evangelischen Missionen}. J. and P. Richter, \emph{Saat und Ernte}. O. Fries, \emph{Geschichten und Bilder aus der Mission}.

Much of what has been said in connection with the Church Book constitutes material of secondary importance. Other material of secondary importance is instruction in Church History, the labors of the Church, and the constitution of the Church.

\quellseite{406}
The conviction that \emph{Church History} supplies suitable material for the instruction of our adolescent youth has been gaining ground. Of course, there is no room in such instruction for intricacies, no need of stressing dates. What is of importance is to picture the main epochs, and the men who were instruments of God in a superlative degree; above all, to describe the origin of our own Church, and of her most important symbolical book; to exhibit the sacrifices and courage of the fathers as a type for all times. What is most important of all is vividly to impress upon the young that Christ lives and reigns; that He has in His hands the threads of history, both of the Church and the world; that He will not permit His Church ever to go down in total darkness, but that He will guide her deeper and deeper into all truth through His Word; that He leads His disciples, in conformity to His own course, through grief and suffering, not to final ruin, but to eternal glory. If rightly taught, i.e., if emphasis is laid upon the biographical element, \emph{Church History} constitutes material decidedly productive of sympathetic, ethical, and religious interest, and for that reason is a great help in the attainment of the aim of religious instruction.

In point of content, such instruction readily joins itself to that of Biblical History for the Church was founded through the preaching of the gospel to Jews and Gentiles related in Biblical History. Subjects to be pictured are the spread of the Church and the persecution of Christians. The latter especially lends itself to detailed presentation, with individual pictures as object lessons. It is the history of the persecutions of the Church that brings out the power of faith to conquer the world; and at this day, it is able to arouse longing for a like faith. With the rise of Constantine, Christ celebrates His victory. Then the inner conflicts begin in the throes of which the confessions of the Church are born. Rise of monasticism. Rise of papacy. Retribution of God through the Saracens upon the lukewarm Church of Asia, Africa, and in part, of Spain. The Church among the German tribes; Charlemagne. Night in the Church, illumined by a few stars. The Holy Spirit seems to have suspended His activity; Jesus, the
\quellseite{407}
Saviour, has receded into the background, Mary and the Saints having replaced Him; instead of Christ, the pope appears to hold the sceptre; and steadily the features he exhibits become more like those of the Anti-Christ. But Christ lives. In men like Wichif and Hus, He makes the dawn to rise, with sunrise in store through the Reformation. The Spirit unfolds His power; the old promises of a Church to endure through the ages are still valid. Luther’s life. Augsburg Confession. The Catechism; divine services in the vernacular; the hymn. Lutheran Church, Reformed Church, Catholic Church; main points of difference. But hard conflicts are to follow: Thirty-years’ War. Paul Gerhardt. A. H. Francke. Era of Rationalism. But the Holy Spirit is mightier than the spirit of man, and the Church is renewed in faith. Not everything indeed is true faith which claims to be such; sectarianism and fanaticism are in evidence. Now fidelity toward the confession of the fathers is required. Fidelity needed especially in our land in which, through Muhlenberg, the Lutheran Church was founded. The Church which does not aim at outward splendor but clings to the Word in faith, is blessed by Christ; hers is the final victory.

These hints may serve as a meager outline for a course in Church History so far as the young have the right to look to the Church for information on the subject. The whole is to be pervaded by the principle: instruction for the purpose of training. For this reason it should indeed always be interesting, but never should it be sensational. As required by the principle of association, instruction in Church History should be joined to religious instruction in general, to the hymn especially. If the heroic faith of the martyrs has been the topic for discussion, the hymn is appropriate: “Onward, Christian soldiers!” If we have seen the fathers confess their faith at Nicea, the pupils may be directed to sing: “We all believe in one true God.” If the class has dwelt upon the tribulations of the Waldenses, the hymn will prove appropriate: “Thou little flock, be not afraid.” When the pupils surround in spirit the stake at which Hus is burning, let the hymn be sung: “Awake, thou Spirit who didst fire.” When the story of the Reformation flashes the Castle Church
\quellseite{408}
of Wittenberg into view, let the hymn be sounded: “Dear Christians, one and all rejoice.” The diets of Worms and Augsburg will call for their share of interest. Let it be signalled by calling for the battle hymn of the Reformation: “A mighty fortress is our God.” The story of Paul Gerhardt suggests the hymn: “Thy way and all thy sorrows.” Other hymns will suggest themselves as the occasion arises. Our \emph{Explanation of Luther’s Catechism} contains a few of the more salient features of Church History (pp. 156–167) which are merely intended to serve as basis for a mere detailed presentation to be given later. Cf. J. H. Schueren, \emph{Gedanken}, etc., pp. 107 f.

Of the various \emph{tasks of the Church}, Foreign and Inner Missions must be stressed especially. How much can be accomplished for the cause of Missions by arousing missionary enthusiasm in the young, has been shown preeminently by Herrnhut. The object must be to exhibit, through pictures from life, the desolation of the Christ-less world, and the marvelous changes wrought by the grace of God. Among the several spheres of Inner Missions, it is especially the Home Mission work, the deaconess cause, and similar works of mercy which require attention. An effective method would be to picture the lives of men who have achieved something extraordinary in those spheres; e.g., the life of Harms in the sphere of Foreign Missions, of Fliedner or Loehe in the sphere of the deaconess cause, that of A. H. Francke, or Wichern, or Geo. Mueller in the sphere of the orphan cause; or the young may be led in spirit into a deaconess home, orphans’ home, etc. Hennig’s Taten Jesu in unsern Tagen, Ohl’s Inner Mission, etc. can be of service in addition to the main factor, personal acquaintance. The labors of one’s own Church ought always to occupy the foreground; and here again those of the particular body with which one is affiliated, not merely because the brethren in the faith are first entitled to
\quellseite{409}
sympathy (Gal. 6:10), and because confederation with other Churches does not bring permanent blessing, but especially because the members of a communion should become so thoroughly identified with some specific task that they are able to say: “This is our work.” Then they will commence to work and to pray for the cause.

So far as the constitution is concerned, the chief elements of the synodical constitution can be supplied in connection with the history of one’s own synod, and the government of the local congregation may be discussed on basis of the congregational constitution.

\quellseite{410}
\chapter[§ 30. THE EDUCATIONAL AGENCIES]{§ 30. THE EDUCATIONAL AGENCIES OF THE CHURCH}

K. Buchrucker, pp. 66–115. *F. Zange, pp. 14–27. O. Baumgarten, pp. 50–85. J. Berndt, pp. 29–33. R. Kabisch, pp. 6–15. E. Chr. Achelis, pp. 338–343. H. Keferstein, \emph{Religionsunterricht u. Erziehg. zur Religion}, 1892. F. Adler, \emph{The Moral Instruction of Children}, 1892. J. McCann, \emph{The Making of Character}, 1900. Ch. W. Rishell, \emph{The Child as God's Child}, 1904. G. N. Coe, \emph{Education in Religion and Morals}, 1904. G. E. Dawson, \emph{The Child and His Religion}, 1909. E. P. St. John, \emph{Child Nature and Child Nurture}, 1911. G. Hodges, \emph{The Training of Children in Religion}, 1911. *Theo. Schmauk, \emph{In Mother's Arms}. Ch. H. Heathcote, \emph{Elements of Religious Education}, 1914. *H. Sneath, G. Hodges, and H. T. Tweedy, \emph{Religious Training in the School and Home}, 1917.

On Home Education: K. Loewe, \emph{Wie erziehen und belehren wir unsere Kinder waehrend der Schuljahre?} 1899. K. Loewe, \emph{Wie erziehe und belehre ich mein Kind bis zum 6. Lebensjahre?} ²1904. A. Matthias, \emph{Wie erziehen wir unsern Sohn Benjamin?} ⁶1907. R. Kabisch, \emph{Das neue Geschlecht}, 1913. E. Lyttleton, \emph{The Cornerstone of Education}, 1914.

On Parochial Schools compare ch. 20. On Sunday Schools compare ch. 20 and add: R. Emlein, \emph{Der Kindergottesdienst}, 1914. J. Steinbeck, pp. 224–238. P. Zauleck, \emph{Theorie und Praxis des Kindergottesdienstes in Vortraegen}. 1914. Ch. S. Albert, The School and the Church. A. H. Smith, The Lutheran Church and Child Nurture.

On Instruction of Catechumens: R. A. Kohlrausch, Der Konfirmandenunterricht, 1898. E. Simons, Konfirmation und Konfirmandenunterricht, 1900. E. Chr. Achelis, \emph{Die Bestrebungen zur Reform der Konfirmationspraxis und des Konfirmandenunterrichts} in Theol. Rundschau, 1901. Freie kirchlich-soziale Konferenz, xv and xvi, 1901. M. Gebhardt, Moderner Religions- und Konfirmandenunterricht, 1906. W. Bornemann, Der Konfirmandenunterricht u. der Religionsunterricht in der Schule in ihrem gegenseitigen Verhaeltnis, 1907. *O. Hardeland, ⁵² Konfirmandenstunden, ⁴1910. R. Steinmetz, Die Bereitung zur Konfirmation in Lehre und Leitung, 1910. O. Lorenz, Der Konfirmandenunterricht, ²1911. *K. Vogel, Seelsorgerlicher Konfirmandenunterricht, 1911. J. Steinbeck, Der Konfirmandenunterricht nach Stoffwahl, Charakter und Aufbau, ²1913.

On Instruction of the Confirmed: *H. Beck, Die Kirchliche Katechisation, 1905. *E.
\quellseite{411}
Siedel, \emph{Lebensphilosophie fuer Juenglinge}, ⁴1896. Mielke, \emph{Die religioese Fortbildung der schulentlassenen Jugend}, 1908. *E. Siedel, \emph{Der Weg zur ewigen Jugend}, ²²1909. *E. Siedel, \emph{Der Weg zur ewigen Schoenheit}, ²¹1911. A. Luettke, \emph{Unterredungen mit der konfirmierten Jugend}, ²1912. M. Reu, \emph{Die Nachpflege der Konfirmierten} in Kirch. Zeitschr., 1912. K. Knoke, \emph{Recht und Pflicht der Kirche hinsichtlich der Unterweisung der konfirmierten Jugend}, 1912. J. Schaller, \emph{The Book of Books}, 1918. M. Reu, \emph{Wartburg Lesson Helps, Senior Department}, 1918 ff. ON YOUNG PEOPLE’S SOCIETIES: Schwanbeck, \emph{Die Juenglings- und Jugendvereine}, 1890. K. Krummacher, \emph{Die evangelischen Jugendvereine und verwandte Bestrebungen}, ²1895. *P. Wurster, \emph{Die Lehre von der Inneren Mission}, 1895, pp. 380 ff. G. Hoezel, \emph{Die kirchliche Vereinsarbeit}, 1906. U. v. Hassel, \emph{Wer traegt die Schuld?} 1907. Hasse, \emph{Leitfaden fuer weibliche Jugendpflege}, 1910 f. H. Weicker, \emph{Der Jugendverein}, 1911. J. Schoell, \emph{Evangelische Gemeindepflege}, 1911. J. Eger and L. Heitmann, \emph{Die Entwicklungsjahre}, 1912 ff. Page, \emph{Evangelische Jugendpflege}: 1914. J. Steinbeck, pp. 280–315. E. Pfennigsdorf, \emph{Christus im modernen Geistesleben}, ¹⁴1914. Luther League Handbook and L. L. Topics. F. G. Detweiler, \emph{Baptist Young People at work}. W. D. Murray, \emph{Principles of the Organization of the Young Men’s Christian Association}, 1910. C. A. Barbour, \emph{Principles and Methods of Religious Work for Men and Boys}, 1912. R. C. Morse, \emph{History of the North American Y.M.C.A.}, 1913. E. H. Engelbrecht, \emph{Manual for Y.P.S.}, 1920. Augustana Luther League Council, \emph{Manual for Luther Leagues}, 1925.

But what agencies are at the disposal of the Church for the inculcation of this material?

The first agency for this purpose is the \emph{Christian home}. Although the home can be called an educational agency only in a wider sense, its significance is of greatest moment nevertheless. Not only is the home a constant factor in the educational work performed by the other agencies, but its own distinctive educational function also must continue after those of the other agencies have set in. In the home the foundation is laid for the training of the child; the home, likewise, must be the faithful custodian of that which is imparted by the other agencies to the adolescent youth. The home, accordingly, has an indispensable place in the group of educational
\quellseite{412}
agencies for the Church’s youth. He who has at heart the training of the young will not fail to give due consideration to the home whence the Church derives the material to be molded. He will rather arouse in the home an understanding for the true interests of the adolescent youth; he will remind it of its duty—doubly sacred by reason of infant Baptism—of rearing the young in the nurture and admonition of the Lord; he will invite it to co-operation with himself; and he will come to its aid with practical suggestions, counsel, and his own example.

McCunn does not exaggerate when he says in The Making of Character: “The vital matter in moral and religious education is the home as it normally is in its habitual preferences, its predominant interests, its settled estimates of persons and pursuits, its ordinary circle of associates, its standard of living, its accepted ideals of work and of amusement. For it is not only from the family but with the family eyes, that we all begin to look out upon the world. And if this first outlook is to see things for which men live in something like their true perspective, and not as distorted through the deluding medium of the home that is idle, frivolous, sordid, grasping, quarrelsome, or sentimental, this will be due far less to what is done of express educational design, far more to the ideal of life which the family consistently embodies. For it is only thus that the scale of moral valuation which the family has wrought into its life, will be likely, as the years go round, to reflect itself in the habitual feelings, estimates, and activities of its members. This kind of influence is, moreover peculiarly effective because it is made easier by the tie of natural affection. Without this, and the trustful confidence which goes with it, comparatively little can be done. And many a parent in whom the qualities which win it have been lacking, even though he may have been masterful and reasonable, has been compelled to realize his impotence. Yet, normally, the parent has a manifest advantage. That confidence which a stranger has to gain with difficulty, he finds either ready to hand, or at most less hard to win. This is double gain. It prompts a spontaneous trustfulness which opens the way for influence, and, as lesser adjunct, it invests a father’s or a mother’s disapprobation with a power to restrain and chasten such as cannot be found when love and trust are absent. In this the family is pre-eminent. No teacher, however kindly, no public authority, however

\quellseite{413}
paternal and mild, can rival it here. And if this be lost, whether by aloofness of parents, or wreck of family life, or by decay of the family as an institution, one of the purest springs of moral influence will be frozen at its source” (pp. 84 ff.).

Since it is the Church of God into which the mature as well as the immature youth has been grafted since Holy Baptism, she merely takes care of her own when she seeks to equip the home for constant co-operation in the work of child-training. A home that refuses to perform this task, cannot expect from the Church the privilege of Baptism for its children. Where, however, the home is willing to perform its duty, but for one reason or another, incapable of performing it properly, the Church must not only see to a better equipment of the home, but also provide an appropriate substitute. In the establishment of infant schools managed by some venerable matron or a deaconess and pervaded by the Christian spirit (hence quite different from the infant schools according to Froebel’s plan), also in the establishment of infant classes in Sunday School, the Church has made at least a beginning in this direction.

An educational agency additional to the home is the \emph{divine service}, however true it is that this is far more than a mere educational factor. Although it is the exception when children understand the sermon as a whole, they do understand many a section of it. With the sight of the worshipping congregation, with the liturgy, the power of song, the lesson conveyed by some striking picture here and there, it cannot be denied that the Church’s worship is a powerful educational factor in the life of the child. While there are biographies in which no wholesome effect is attributed to the attendance upon worship, there are others in which the contrary is asserted and permanent impressions are traced to the factor here stressed.

Rightly says Hodges in the chapter treating of “Sunday and the Children” in \emph{The Training of Children in Religion}: “On Sunday
\quellseite{414}
morning the children are to be taken to church. It is true that there is the peril of having the children disturb the service, and annoy the congregation; and there is also the peril of making the children hate the service. But if, in avoiding these dangers, the children are left at home, or are encouraged to consider the Sunday school a substitute for the church, there is a possibility that they never begin to go. Perhaps the best thing to do, in the midst of these perplexities, is to keep church-going as a special privilege and reward for good behavior until the age of reasonably steady habits. Then insist on attendance at church like attendance at school as a normal part of decent living and subject to precisely the same excuses. Commonly, children who are old enough to go to school are old enough to go to church. The two introductions to a wider experience may properly be made at the same period. Then nothing should be allowed to keep the child from church which would not validly keep him from school. Thus the habit is formed, and a solid contribution is made to the child's religious education. Indeed it depends much upon the character of the church service. When this, besides the sermon, provides responses and singing in which the children may join, then such a service is varied, with downsitting and uprising, and is constructed according to some understanding of human nature, the common nature not only of youth but of maturity. Even the children are not wearied.—If a choice must be made between the church and the Sunday school, on the ground that the two together are too long and wearisome, a wise preference will probably select the church. In that case, special attention will need to be given to the instruction of the children at home. One of the unedifying sights of our Christian Sunday is the spectacle of troops of children, dismissed from the church school, passing the church door, to spend the rest of the day idly. It is a prophecy of a coming generation of non-churchgoers. The habit of churchgoing, which counts for so much in the orderly religious life, is thus omitted. The boys and girls do not begin to go to church. Then, when they get through the Sunday school, in the midst of their teens, they are likely to turn their backs upon the whole system of organized religion. To this result, such training leads. Properly, the Sunday school should be a preparation for the church. The children should there be instructed to take part intelligently in the church services. They should learn there the words and music of the hymns which the congregation sing. There is opportunity,
\quellseite{415}
indeed, among the younger children for the singing of hymns which set forth the appropriate emotions of early childhood. But the older children are to learn the classic poetry of the hymnal just as they learn in the day school the classic poetry of the great masters. Some of it will far exceed their present experience, but no more so than the lines of Shakespeare and Milton, no more so than the Psalms. They will get enough for present use, and will store up treasure for the time to come. And in the church they will be able to take part with their elders. The Sunday school, however, is not a fair substitute for the church. It gives but little time to the great religious exercise of worship, and it makes no place for the sacramental side of the spiritual life. It does not bring the child into accord with the ancient, venerable, and universal expression of corporate devotion. It does not initiate him into that society in which he should have his membership all the rest of his life. It does not put him where he ought to be, and start him right.

“We probably underestimate both the endurance and the interest of children. No doubt, our forefathers demanded somewhat too much, with their protracted prayers and sermons and their two Sunday schools, one in the morning and the other in the afternoon. We are in danger of going into extreme over-anxiety as to the children’s comfort. They seem to get along pretty well with a good measure of time on week days at their schools. What is asked on Sunday is after all not excessive. The Sunday school lasts an hour, and the morning service rarely exceeds an hour and a half and is often shorter than that. It is true that at times most normal children rebel against both these forms of confinement in-doors. The best plan is to deal with their feeling about the church as we deal with their feelings about the school. We make them go to school, whether they like it or not. Sometimes they hate the school, but they almost always come into the right mind about it, and are grateful for the discipline. They may, on unusually pleasant Sundays, hate the church. No matter. If they are easily permitted to stay at home, they will despise it; and that is worse. A little wholesome temporary rebellion is better than contempt” (pp. 266 ff.).

A further noteworthy truth is found in what the same author says in \emph{Religious Training in the School and Home}: “The very building of the church is a very helpful factor in the moral and religious development of the child. . . The dignity and beauty of the architecture, the loftiness and spaciousness of the walls, the suggestion of pulpit and communion table, of font and altar, all beget a sense of solemnity
\quellseite{416}
and awe, of wonderment and of hushed expectancy that awakens and develops the moral obedience and the intelligent worship that are to be. . . In the poorest service the child gets something. His spirit is bathed in the awe and silenced in the hush. He visions the great and the good — or at least those whom he believes to be the great and the good—bound in penitence, standing jubilant in praise, instructed in the Being infinitely above him in wisdom and in goodness and in love. All this is educative. Where the minister has trained himself for this important part of his task, and is awake to the little faces that look up into his own, the hour becomes a memorable one. Many a restless and apparently thoughtless member of the junior congregation grows up to testify to the power which was exercised over him by hymn and sermon, little as they seemed to do for him at the time” (p. 7).

To be sure, attendance at divine service alone is not enough. If religious knowledge is to be something more than a group of disjointed fragments; if the will is to be habituated to that which is good, \emph{the school must be drawn into service} as an educative agency additional to home and divine worship. We do not think in this connection of the public school of our country. For though the attempt has been made again and again to incorporate in its curriculum moral and religious training in some form, such attempts are in opposition to the constitution of our country and to the character of the public school as determined by it. Whatever religion is tolerated in the public school is one from which the very heart has been torn, nor has it ever proved tenable in actual experience.

In several States it has been attempted to make the public school an agency for moral training. While, for instance, in North Carolina the teachers are merely requested to “encourage morality,” West Virginia has gone so far as to charge “all teachers, boards of education, and other school officers with the duty of providing that moral training for the youth of this State which will contribute to securing good behavior and manners, and furnish the State with exemplary citizens.” So far as religious instruction is concerned, in distinction from moral, the view has been entertained in some States that it might be rendered possible, at least in an elementary
\quellseite{417}
degree, by introducing the reading of the Bible in the school. The present status of this question in the Union is as follows: Bible reading is required by law in Alabama, Delaware, Florida, Georgia, Idaho, Kentucky, Maine, Massachusetts, New Jersey (only the Old Testament), Pennsylvania, and Tennessee. Bible reading is specifically permitted by law in Indiana, Iowa, Kansas, North Dakota, Oklahoma, and South Dakota (the only State which allows “unsectarian comment” on the portion read). State Supreme Court decisions favorable to the use of the Bible were handed down in Michigan and Texas (also in Ga., Ia., Kan., Ky., Me., Mass.); the decisions handed down in Nebraska, Wisconsin, and Ohio while not specifically forbidding, are generally considered adverse to Bible reading. The Bible is used in (some) schools under general terms of the law or by reason of its silence on the subject in Arkansas, California, Colorado, Connecticut, Maryland, Mississippi, Missouri, New Hampshire, New Mexico, New York, North Carolina, Oregon, Rhode Island, South Carolina, Utah, Vermont, Virginia, and West Virginia. The Bible is excluded from the schools by court decisions or by the Attorney General in Illinois, Louisiana, Minnesota, Washington (and apparently Nebr. and Wis.). The Bible is excluded by a ruling of the State Board of Education or the State Superintendent of Public Instruction in Arizona, Nevada, Wyoming, and Montana. In Indiana and North Dakota high school credits are given for the study of the Bible as literature outside of school hours. New York City furnishes religious teaching without charge outside of school hours in school buildings, skilled public school teachers giving the courses in religion. Cf. Wilbur F. Crafts, \emph{The Bible in School Plans of Many Lands}; Ch. N. Lischka, \emph{Private Schools and State Laws}, 1925; also the \emph{Report of the Commission on Christian Education} to the Quadrennial Meeting of the Federal Council of the Churches of Christ in America, 1916.

What are we to think of this? In order to substantiate more fully the opinion expressed above, we shall refer purposely to the opinions of non-Lutheran pedagogues. Especially noteworthy is the stand taken by Charles H. Heathcote in The Essentials of Religious Education. A decided advocate of religious education, he makes the following characteristic statement: “In theory, religion and education may be separated from each other; but in reality such a thought is impossible. The aim and goal of education and religion are virtually the same. The basis of true education is religion, and any effort to make education independent of religion narrows its scope,
\quellseite{418}
aim, and goal” (p. 7). Notwithstanding, he continues: “According to the interpretation of our constitution, religion cannot be taught in our public schools. . . We recognize the importance of reading the Bible in the schools, but we cannot call this exercise religious instruction, nor would we want it designated by such a term. When this reading is done with discrimination and without comment on the choice and splendid portions of Scriptures which should always be read, we believe great good can be accomplished in many ways. This reading may not be instructive in the analytic sense, but the mind of the pupil is impressed with the beauty and simplicity of God’s Word.—There are many educators and religionists who advocate the study of religion in the public schools. They make a contradistinction between religion and denominationalism. They advocate that religious instruction based on broad general terms of religious concepts, free from doctrinal, creedal, and denominational interpretations, could be put into the school curriculum. We realize there is much force and consideration given to their arguments, but we cannot see the feasibility of the plan. We believe in the broad interpretation of religion for it is to be thoroughly adaptable to all classes and conditions of humanity. When we speak of religion, we are, of course, referring to the broad principles of Christianity upon which the advocates of this theory agree, as it is the only religion which can give a positive civilization to the world. We do not believe that the introduction of religious instruction, even upon the basis of the broadest interpretation of Christian teachings, would work in practice. It is very evident that such a plan would not be an acceptable one to the Hebrew, Catholic, and a majority of the Protestants and many other forms of religious life which are represented in our public schools. All these conditions must be borne in mind in advocating this theory. At the same time it is well-nigh impossible to interpret religion on the broadest basis to eliminate every iota of denominational and doctrinal view-point. The public school is not a religious nor antireligious school, but it is a secular institution, and we want to see it remain such”. And afterward, on pp. 11 ff.: “We cannot agree with Dr. Seeley (\emph{Foundation of Education}, p. 248) in advocating the study of the Bible as a religious book. The Hebrew father would have every right to object, according to his religious belief and traditions, to the principles of the Christian religion being brought to his child. . . Religions cannot consistently be taught in our public schools.” Nothing remains therefore as bearer of religious education but home and Church.

\quellseite{419}
Similar conclusions, though from different premises and with different aims in view, are reached by G. A. Coe, in his work, \emph{Education in Religion and Morals}. After animadverting upon the introduction of the Bible in the public schools and upon the notion of teaching certain fundamental religious truths upon a biblical basis, he continues: “What then should be the next move toward improving the relation of the State schools toward religion? Without hesitation it may be said that the next move should be to induce the family and the Church consciously to assume their proper share of the responsibility for the character of the rising generation. Let us remove the beam that is in our own eyes. If I let weeds go to seed in my door-yard, they spread to my neighbor’s door-yard; but if I make my door-yard beautiful with flowers, I make it easier for my neighbor to beautify his own premises. As soon as the family and the Church are sufficiently aroused to begin to do their own duty, the public school question will grow wondrously simple. Strong purpose is contagious, and it has a remarkable way of finding methods. Our trouble is that we have not reached the point of giving ourselves to this reform. We are giving, instead, advice and criticism to the public schools, and in various ways we are hoping that organizations, methods, and schemes will do what only personal consecration can accomplish. We neglect the children in our homes; we do shilly-shally work in the Sunday school, and then shift to the State school the blame for the results! It is well, to be sure, to adopt at once feasible means of improving the State school, and depend upon it—any large and thorough improvement therein will wait until, through striving to build, each over against his own house, the churches and the homes have developed a proper educational consciousness among the people” (pp. 360 ff.).

When we speak of the school as an educational agency, we think primarily of the regular \emph{church, or parish, school}. We give it this name for the reason that the individual church, or parish, establishes and supervises it. Where the parish school is managed according to the basal idea that has given it birth, instruction will be imparted not only in religion but also in all other branches that have a place in the curriculum of a Christian public school. Such a school will preserve its unique character by teaching everything that is taught, according to the mind and in the spirit of Jesus
\quellseite{420}
Christ, and in harmony with the confession of the Church which sustains it. It will be arranged in grades, as required by the development of the children; it will be managed by teachers who have been specially trained for this work. Where the teachers understand their duty, faithfully fulfil it, and are themselves imbued with the spirit of Christ; where the pastor visits the school diligently, is in closest touch with its life, supervises all school activities, sees to it that the curriculum is carried out; where he pleads for the school in his sermon, in his pastoral work, and in the assembly of the congregation as a vital factor in congregational development, and where by such loyal co-operation of pastor and flock the congregation learns to recognize the great blessings of its school and freely entrusts the children of school age to it—there a constant stream of blessing will flow from such a school into the congregation, especially where the Christian home works hand in hand with the school. Beware of silencing the proposition to establish such a school with the phrase that it is an \emph{“exotic growth”}! This it cannot be for the very reason that, unlike the state school of Germany, it is a Church school (cf. pp. 188 f.). There is another reason: It was in actual existence in our country for a long time before there was such a thing as a public school, even apart from Lutheran congregations. The fact of its gradual extinction in the eastern part of our country, while partly due to other causes, is largely explained by a diminution of that spiritual energy and the lack of that vigorous Lutheran consciousness which should have pervaded the whole church life, the sphere of education included. Where no parish school was established in the West, or having been established, was permitted to pass away, the same causes are likely to have been in operation in one way or another. There is no doubt that the establishment and maintenance of such a school becomes more difficult where
\quellseite{421}
the maintenance of the mother tongue brought across the ocean is no longer one of its aims. But should the language which our children speak make us indifferent to the spirit in which geography, history, natural science, etc., are taught? Is it immaterial now whether these branches are taught in the spirit of the Word of God, which always discloses the wonder-working hand of God, or in the spirit of materialism and of the evolution theory, which pervades the textbooks of our public schools openly or in disguise; in the spirit which knows of no consummation of all things until the present process of development shall have been terminated by the judgment of God, or in the spirit of Spencer’s philosophy and pedagogy, according to which the history of humanity is a path constantly climbing higher—up to the highest heaven, which mankind is to create for itself by the labor it performs in its own behalf? Shall the natural man and an earth-born morality have the last word in the premises? Let us try to understand the idea in all its bearings: For six days in the week our children are to be exposed to the views of the natural man, and—perhaps—perceive a little of eternity’s breath in Sunday School on the seventh! Which is bound imperceptibly to gain the upper hand in these circumstances, especially where the home and its spirit does not rear a strong protecting wall around the soul? The spirit from the abyss, the spirit of lies, which utterly confuses all Christian concepts in the soul, and in a masterly manner represents as Christian what never can be such, is at work more busily than ever. Hence it is more necessary than ever to protect the souls of the rising generation from their earliest youth by instruction emanating from the spirit of Christ. And the \emph{Christain parish school} is the first place where that should be expected to be found. It is true that today when the public school system has reached a high degree of efficiency in equipment and method,
\quellseite{422}
the establishment and maintenance of a parish school is much more difficult than formerly, and that it offers no prospect of permanence in the cities unless several teachers are in charge; but we have quite a number of opulent city congregations which are equal to the financial demands thus made upon them, the more so as healthy co-operation shall tend to supersede in our Church the wretched competition of the present.

While it is therefore the holy duty of the Church in general and of the individual congregation in particular to establish \emph{regular parish schools} and to equip them with the requisite lesson helps and teaching force, conditions are not rare in which the establishment of \emph{regular parish schools} is precluded. Wherever that is the case, it is desirable, especially in country districts, to make all possible legitimate efforts to have \emph{Christian, Lutheran teachers} called to the public schools where—without in the least modifying the non-religious character of the school, but in a manner comporting with their Christian character—they may teach at public expense for as many months as is required by State law. Rome knows what she is doing when she sees to it that a disproportionately large number of teachers of her faith are given charge of the public schools. While the motive, which cannot be any other than the ultimate domination of the State by the Church, is questionable, the breadth of vision and the energy in the execution of the plan may serve as an example. In the months not taken up by public duty, such Lutheran school teachers, then, may teach parish school at congregational expense. — In cities where the establishment of a regular parish school should prove to be impossible, efforts should be made to send the children to a parish school, at least from the sixth to the ninth year. When they subsequently have to attend the public school, the spiritual influence of the first school years can be perpetuated
\quellseite{423}
by means of the Saturday and Sunday School. Where the distance from the parish school is so great that children from six to nine years of age are unable to cover it, such a partial parish school should be attended by the children during the years of preparation for confirmation.

Should even this be precluded by the shortsightedness of the leaders in the congregation or by a lack of training in the spirit of sacrifice and by unwillingness to be trained, efforts should be made to have the local or, still better, the State authorities \emph{set aside for instruction in religion a certain number of hours of the time claimed by the public school}. In proportion as the necessity for religious training shall impress itself upon the public conscience, such a plan will meet with the sympathy of the school authorities, unless the spirit of syncretism, now in the ascendant throughout the country, and aiming at a universal religion, shall foil the plan. Such instruction in religion by the Church should be supervised by a pastor conversant with teaching methods, and be imparted by teachers of a decidedly Christian character whose general equipment for teaching has been supplemented by a special course in religion and Christian pedagogy, arranged by the Church. The effect upon the child will be much deeper and more enduring if the teacher of religion is at least the equal of his colleagues in the public schools in point of technical training; if the teaching facilities are no less adequate than those obtaining in public instruction; if religious instruction is imparted in the rooms of the public school building, and if no hours are appointed for such instruction when the children are worn out by fatigue. Although the Christian character of the teacher is bound to be a decisive factor in religious instruction, no matter how adverse the conditions, there are psychological and pedagogic reasons that favor the measures above recommended. By
\quellseite{424}
systematic effort, such arrangements will generally, though not always, prove feasible. Cf. § 20.

Should the public authorities refuse to defer to the request of the congregation in the matter, it would not be advisable to wait for a change of sentiment. In that case, or if a parish school, either in the regular or partial form, should prove impracticable, \emph{a combination of Summer, Saturday, and Sunday School} will largely solve the problem. At least in the place where the pastor resides, this combination is feasible and has long since passed the experimental stage. It was subjected by me to a practical test as early as thirty-eight years ago; and since that time it has been found practicable in bilingual as well as exclusively English-speaking congregations. \emph{The Summer School} should be conducted for eight weeks at least three hours each day. Not a few pedagogues have stated as their opinion that a vacation of three months is accompanied by great disadvantages. Not a few local school authorities have established public summer schools; in some cities vacations have even been abolished, the time of daily instruction being shortened instead. Such measures constitute an irrefutable answer to the assertion that a religious summer school is a measure of cruelty against the children and an infringement upon the leisure to which they are entitled. \emph{Saturday School} is conducted throughout the year, with the possible exception of the eight weeks occupied by the Summer School; from three to five hours should be devoted to it. All the children of the congregation from the sixth or the eighth year respectively should attend these schools. Since only religious instruction, in its various ramifications, is given (Biblical History, Biblical Geography, Church History), much can be attained for the cause of education and training of children, provided the schools are characterized by system and conscience. Results such as may be attained in a properly conducted regular church school, however, are indeed beyond the scope of partial
\quellseite{425}
parish schools or these combination schools, because the influence of the latter is in part paralyzed by that of the non-religious State school, and because the factor of habit, which looms large in the sphere of child training, is not sufficiently operative. The Saturday School should supplement the Sunday School. The “presentation” of the Bible story given in the Sunday School is to be followed by the “penetration” and “application” in the Saturday School (cf. § 31). At any rate, the Saturday and the Sunday Schools must adjust themselves to each other. Any school whatever which aims at respectable achievements, the Summer School and the Saturday School included, should have a regular curriculum. So also \emph{the Sunday School should have a regular curriculum} embracing all grades. But such a curriculum cannot be established unless one has clear views in regard to the aim and tasks of the Sunday School. Where a regular parish school is in existence, and really attended by the whole congregational youth from the sixth year upward, the Sunday School is hardly a necessity. In that case it would be better to change it into a junior congregation, provided this would not affect attendance at the regular services, especially so far as the upper classes are concerned (cf. p. 391). The situation is different where the Sunday School can be made a missionary agency for the purpose of gathering religiously neglected children. Where there is no parish school, or where the parish school is not attended by all the children of the congregation, the Sunday School is a necessity, and its session is a school rather than a divine service. There must be teaching there, reciting, the covering of a definite amount of ground; but all this must be under the control of the paramount idea: the instruction given is religious, and the object is training; an impression is to be made not upon the intellect alone, but upon the whole man, who, through such instruction, is to be trained for something higher. The curriculum, however,
\quellseite{426}
can be carried out only where the several assistants or Sunday School teachers, regularly meet the pastor for their own instruction first, in order to discuss the subjects to be taught, both in regard to content and their bearings upon the several grades, and thereupon make thorough preparation for their classes. This involves the other proviso that, with the exception of the infant classes, the school is occupied not with diverse but with the same material (one and the same Bible story).

To the educational agencies mentioned—the home, the divine service, the church school in one form or another—must be added the \emph{instruction of catechumens}, which immediately precedes confirmation and dare never be omitted. Where there is a properly maintained regular parish school, three hours per week for six months will suffice for this instruction, especially if the pastor himself imparts religious instruction in the upper grade of the parish school. Where there is a partial parish school, so constituted that it is attended by the children of twelve and thirteen years of age for two years, the same number of hours will prove adequate. If there is neither a full nor a partial parish school nor “week day religious instruction,” catechumenal instruction should be given four times a week for two years unless incalculable injury is to be inflicted upon the congregation. If the catechumenal lesson could be given from eight to nine, instead of from four to five, a stronger effect would be achieved inasmuch as the children would be fresh and the whole day would come under the influence of the initial lesson. If the catechetical material is offered at a time of fatigue, the danger arises that it will be identified with ennui and fatigue. If such instruction is imparted as a mere appendage to public school instruction, the soul is in danger of falling a victim to the fallacy that religion is of less importance for the tasks of life than arithmetic and rhetoric. Should anyone
\quellseite{427}
even think that he can prepare his pupils for confirmation and the Christian life with a few “lectures”, he is guilty of most conscienceless conduct both from the pedagogic and the religious standpoint. The instruction of catechumens is the most important instruction; he who is not faithful here is not fit to hold the office which preaches reconciliation with God; and the more any congregation should insist upon shortening the time devoted to instruction, the more decided the pastor should be in opposition. — Should there be a second congregation in the parish, Sunday School should be maintained there regularly and, if possible, Saturday School; and wherever possible, the children of such congregation should join those of the main congregation for instruction during a term of from four to five months at least. Only self-denial and conscientious purpose can produce order and efficiency.

With confirmation, however, religious instruction is not brought to a conclusion. While an important goal has been reached, it is, after all, but preliminary to one beyond. Confirmation should be followed by further care. \emph{The Christian home} serves the purpose of helping the young to retain and exploit the treasures acquired. It has been shown above how much depends in just these years upon the example of father and mother: their mature faith, their firmness of character ought to inspire, influence, strengthen the youths’ religious life (cf. pp. 389–391). Family devotion both morning and evening; regular attendance at divine services; common reading of good literature; common pleasure and recreation; constant self-criticism and self-control in expressing judgments, in conversation, in actions; sympathy for questions engendered by youthful doubt; toleration for everything lawful; truthfulness, firmness, and freedom in noble union — these are requisites of the Christian home. It is to be regretted that the home, but too often, betrays too little
\quellseite{428}
understanding, willingness, and readiness for such ideals, and requires a vast amount of patient, toilsome training. Ofttimes it is even an enemy of such Christian training, not only by reason of its indifferent religious and moral attitude, but also because the axiom, “Confirmation is the close of religious education,” not without the fault of the Church, has been ingrained into its very life. Therefore, if the confirmed youth is to receive adequate care and the work begun is to be continued, another educational agency must take its place beside the home. We refer to the Bible Class of the Sunday School and the Young People’s Society. Attendance at Bible Class should be the self-evident duty of all the confirmed down to the dangerous sixteenth or seventeenth year (cf. p. 264 ff.). Those who have attended it during those years, are likely to do so also afterward. The Young People’s Society exists not only as an external tie for the confirmed youth, but also, and preeminently, as a means of their further mental and spiritual culture. It adopts its own constitution; the general direction is committed to the pastor (cf. \emph{Rules for Young People’s Societies}, Waverly, 1912). Either it is divided into two sections (from 13 to 16, and 17 to 21 respectively), or it is recruited from the ranks of those who have gone through the Bible Class of the Sunday School. If attendance upon the Bible Class (also upon the meetings of the lower section of the Y.P.S.) is a matter of obligation and congregational enactment, membership in the upper section is left free. Properly conducted, the Young People’s Society is an important educational agency, a source of blessing for all participants, and a society from which later the ranks of the most useful and mature members of the congregation recruit themselves. But prospect of permanence exists only where, from the very beginning, the work conforms to a preconceived plan; otherwise it will not survive the charm of novelty. It stands to reason that it should be, in part,
\quellseite{429}
a school of training for the future assumption of the burdens and duties of active membership in the congregation. The main stress, however, must be laid upon safeguarding and experiencing the things learned during the years of school, though under new aspects (cf. § 31), and always with the practical life as background. Thus viewed and conducted, the Young People’s Society requires indeed labor; but what is a faithful pastor not ready to do when he perceives that his labors can be a blessing to the congregation! Also here Luther’s word holds true: “Christ Himself shall be our reward!”

The grading of these educational agencies is precluded from the outset so far as the home, the instruction of catechumens, and the Young People’s Society (upper section) are concerned. It is indispensable in the parish school, as well as in the Saturday, Summer, and Sunday Schools. According to what principles the grading should be done, is shown in § 22 with sufficient clearness. If no more than two grades are possible, the division line will be the ninth or tenth year; if three grades are convenient, they will extend from the sixth to the eighth, from the ninth to the eleventh, and from the twelfth to the thirteenth or fourteenth year.

\quellseite{430}
\chapter[§ 31. THE DISTRIBUTION OF THE MATERIAL]{§ 31. THE DISTRIBUTION OF THE MATERIAL OVER THE DIFFERENT EDUCATIONAL AGENCIES AND THEIR SEVERAL GRADES}

Compare the literature given in §§ 20 and 30. *F. Zange, pp. 64–142. A. Eckert, pp. 62–78. J. Berndt, pp. 68–73. R. Kabisch, pp. 157–181. R. Seyfert, \emph{Versuch eines Lehrplans fuer den Religionsunterricht}, 1890. J. H. Schueren, \emph{Gedanken ueber den Religionsunterricht}, 1900. W. Armstroff, \emph{Die einheitliche Gestaltung des Religionsunterrichts in Kirche und Schule}, 1896. H. Brammer, \emph{Neue Bahnen fuer den Religionsunterricht}, 1900–1903. *A. Reukauf, \emph{Zur Lehrplantheorie der geschichtlichen Stoffe im Religionsunterricht der Volksschule}, 1906. Baumann, \emph{Lehrplan f. d. ev.-luth. Religionsunterricht der achtklassigen Volksschule}, 1906. S. Bang, \emph{Lehrplan f. d. 2–8 klassige einfache Volksschule}, 1906. A. Reukauf, \emph{Didaktik des evang. Religionsunterrichts}, 1906. *G. W. Pease, \emph{An Outline of a Bible School Curriculum}, 1909.

So far as the home is concerned, a distribution of the catechetical material is out of the question. While the home, as the child grows up, becomes largely the custodian of blessings imparted to the child in successive stages, any communication of religious truth in the childhood stage is of necessity spontaneous and occasional. Of what radical and decisive moment the infant period is for the later development and the formation of character, has already been shown on pages 245 ff., where also the needed hints as to the material to be employed are given. Here we merely desire to emphasize the point that the education dealing with the infant dare not disturb the quiet sanctuary of the awakening inner life; its aim will be above all to clarify, to guide, to achieve results through the living example. Here above everything the pedantic tone of the schoolmaster should be avoided and all cold reflection likewise! Cautious and chaste nurture of the nascent life imbedded within through Baptism should be the aim.

When we come to speak of the second educational agency, however, the school, a distribution of the material is necessary.

\quellseite{431}
We concern ourselves first of all with the regular parish school. Since Biblical History is to be taught in all its grades, the principle governing the distribution of this material should be clearly understood first of all. For a long time, the curriculum was devised on the \emph{concentric circle plan}, according to which originally all Bible stories were gone through with all the classes year after year (the Old Testament in the summer, the New Testament in the winter). Later this was changed in that a few, easily comprehended stories were taken up in the Lower Course, a review of these, with a few added ones, was undertaken in the Intermediate Course, and the full number of stories contained in the Biblical History was reached in the Upper Course. This method of distribution was made an object of attack by Ziller and his followers who advocated with all their might a curriculum based on the successive culture epochs, \emph{kulturgeschichtliche Stufen}. According to their plan the children should be made to live over again in their own experience the history of human development from the Patriarchs to the life of faith typified by the Evangelical Church. In pursuance of this plan, a preliminary course, ranging over two or three years, was arranged, with tales and Robinson Crusoe or with religious subjects as material. This was followed by the history of the Patriarchs in the third school year, the story of Moses and the Judges in the fourth, the story of David and Solomon in the fifth, the story of Jesus together with selections from the Prophets in the sixth, the Acts of the Apostles together with parts of the Epistles in the seventh, and the history of the Reformation together with a final resume of the Catechism in the eighth. Against the curriculum patterned upon concentric circles, as originally devised, the argument was advanced that, in carrying it out, too much was gone over in one year, which prevented thoroughness and for this reason made impossible the inner experience and aroused no speculative interests. Against its first form and its subsequent
\quellseite{432}
modifications the argument was raised that, with a repeated treatment of the same stories, ennui resulted for the children—an outcome unfavorable to the arousing of the religious and sympathetic interests. These arguments really disclose the weaknesses of the curriculum based upon the concentric circle plan. But the curriculum based on the culture epochs is not qualified to take its place. A crushing argument against it is the untenableness of the underlying idea, which, ignoring the fact of Baptism, does not introduce the child to religious materials until the third or fourth school year. It can also be urged that, without a repeated treatment of the subject matter of Biblical History, it is impossible to bring about such conversance with it as the young stand in need of upon entering life, even the socalled immanent review\footnote{By immanent review we mean the review connected with the treatment of the new material; I employ immanent review either when I proceed from the known material to discover unknown truths, or when I combine the new material with the old and fuse them into a unit. Compare the subject of “penetration”, § 32.} not being sufficient for the purpose. What is, however, eminently worthy of attention is the effort to present the Bible stories in large groups, for only in this way a vital penetration into the history of the several biblical characters, is made possible and, likewise, the arousing of a many-sided interest, sympathetic, religious, moral. Accordingly, a curriculum that possesses the advantages of the two mentioned, without being marred by their faults, will be found to be the most efficient one. For this reason simple, but stirring stories chosen from the New Testament as well as the Old, should be introduced already in the Lower Course; with the proviso that the stories selected cover not a great number of periods but only one or two, and, in any case, constitute a natural unit. In the Middle Course there should be a review of the material previously gone through; but these stories, now known, should at this time be placed under a higher aspect and not be
\quellseite{433}
treated at the same length, while the new stories should by far preponderate. But whatever material is taken up for treatment should constitute a compact group or a number of compact groups. All the lesson material hitherto used will be coordinated in the Upper Course into a coherent Sacred History. In the Lower Course the selection of material should conform as much as possible to the church year; in the Middle Course at least the festive stories will be reviewed in due time, while the church year receives no further consideration in the Upper Course than is necessary to read and explain the pericopes.

Carrying out the principles here laid down, we offer the following curriculum, sketched upon the supposition that the three grades of the lower course have a common instructor, likewise the middle course of two grades and the upper course, also consisting of two grades. The material assigned to the first or second grade of each course is so shifted that the pupils entering in 1927 and those who have entered in 1926 begin with the material selected for the second grade and then, in the second year take up the material designed for the first grade. Since, in the lower course, all instruction should be fused into an organic unit with Biblical History, we show at the hand of several examples from the lower course how such fusion is to be accomplished. At the same time we refer to our \emph{Wartburg Lesson Helps} and our \emph{Biblical History} (Chicago, 1918) where this fusion has been effected.

\emph{First or Primary Course} (comprising three school years, age 6–9). \emph{First Grade:} (1) How God created the world. A (=first year): Ps. 115:3; B (=second year): Ps. 104:24; C (=third year): Ps. 33:9; Morning prayer; Catechism: First Article without the explanation.—(2) How God made man happy. A: Ps. 118:1; B: 1 John 4:19; Catechism: First Commandment without explanation; Hymn: Awake, my heart rejoicing, 1st stanza; C: 1 John 5:3; Catechism: First Commandment with explanation; Hymn: Praise to the Lord, the Almighty, stanzas 1 and 2.—(3) How man sinned. A: Ps. 5:4; Prayer: Abide in grace, Lord Jesus, 1st stanza; B: Prov. 1:10; Prayer: Lord Jesus, who dost love me; C: Ps. 37:37; Catechism: Conclusion of the Commandments without explanation; Hymn: Abide in grace, Lord Jesus, stanzas 1, 5, 6.—(4) The Saviour
\quellseite{434}
came into the world. A: Luke 2:10–11; Hymn: As each happy Christmas; B. John 3:16; Hymn: Praise God the Lord, ye sons of men, stanzas 1, 2, 8. C: Luke 19:10; Hymn: Good news from heav'n, stanzas 1–4.—(5) Wise men from the East came to the child Jesus. A: Matt. 7:7; Hymn: Brightest and best of the sons of the morning; B: Prov. 23:26; Hymn: Thou, whose almighty word; C: 1 Tim. 2:4; Acts 4:12; Hymn: Thou whose almighty word; Catechism: Go ye and teach all, etc.—(6) How the child Jesus was saved from great danger. A: Ps. 34:7; Prayer: Lord Jesus, who dost love me; B: Ps. 91:11; Hymn: If thou but suffer God to guide thee, 1st stanza; C: Ps. 23:3; Hymn: Thy way and all thy sorrows, stanzas 1 and 2.—(7) The child Jesus is brought to the temple. A: Luke 11:28; Catechism: Third and Fourth Commandments without explanation; B: Ps. 26:8; Hymn: Blessed Jesus, at thy word, 1st stanza; C: Third and Fourth Commandments with explanation.—(8) How Jesus blessed Peter’s draught of fishes. A: All depends on our possessing God’s true love, and grace, and blessing; B: Ps. 127:1; Hymn: Take thou my hand, O Father, 1st stanza; C: Matt. 10:37; Hymn: Jesus, still lead on, 1st stanza; Abide in grace, 4th stanza.—(9) How Jesus saved His disciples during a storm. A: Ps. 50:15; Hymn: Ever is a peril near me, 1st stanza; B: Matt. 28:18; Hymn: Abide in grace, 6th stanza; C: Rom. 8:31; Hymn: How the wind in fury blind, 1st stanza.—(10) Jesus feeds the hungry multitude. A: Come, Lord Jesus, be our guest, etc. B: Ps. 103:13; Catechism: Prayer before meat; C: Ps. 145:15. 16; Hymn: If thou but suffer, stanzas 1–3; Catechism: Prayer after meat.—(11) Jesus heals the impotent man at the pool of Bethesda. A: Luke 1:37; B: Ps. 118:8; C: Ps. 103:8; Hymn: Thy way and all thy sorrows, stanza 6.—(12) How Jesus comforted a sorrowing father by restoring his daughter to life. A: I am the resurrection and the life; B: Mark 5:36 (Fear not; only believe); C: Ps. 25:3; Hymn: When children young and tender, 1st stanza.—(13) How Jesus restored a son to his mother. A: Rom. 12:5; B: 2 Tim. 1:10; C: John 5:28. 29.—(14) Jesus rejoices with the rejoicing and helps them out of their difficulty. A: Rom. 12:5; B: Ps. 34:8; Hymn: I am Jesus’ little lamb, 1st stanza; C: Ps. 77:14; Hymn: Thy way and all, 4th stanza.—(15) Jesus desires that little children should be brought to Him. A: Mark 10:14; Hymn: The truest Friend abides in heaven; B: Hymn: Jesus now thine own forever; C: Hymn: I am Jesus’ little lamb, stanzas 1–3.—(16) How Jesus entered Jerusalem
\quellseite{435}
on Palm Sunday. — (17) How Jesus prayed before His suffering. — (18) How Jesus was betrayed by one of His disciples. — (19) How Jesus died for us on the cross. — (20) How Jesus was buried on Good Friday. — (21) How the risen Saviour on Easter day appeared unto His disciples (John 20). — (22) How Jesus ascends into heaven in order to be with us always.

\emph{Second Grade} (1) Why Abraham, the man of faith, left his home. (2) How Abraham preferred to be separated from Lot rather than live in strife and contention. (3) How Abraham proved to be a friend in need. (4) What pleased God most in Abraham. (5) Abraham invites three guests and thus receives a new promise. (6) Abraham intercedes for the people of Sodom. (7) How Sodom and Gomorrah were destroyed because of their sin. (8) God gives Abraham a son, but Abraham loves God more than his son. (9) God provides a wife for Isaac. (10) Abraham dies and is buried by his son. (11) Jacob deceived his father and his brother. (12) Jacob fled from his brother. (13) Jacob saw the heavens open in his dream. (14) Of Jacob’s life in the home of Laban. (15) Jacob flees from Laban. (16) How Jacob meets his brother Esau. (17) Of Joseph, Jacob’s favorite son. (18) Of Joseph’s life in the house of Potiphar. (19) What happened to Joseph in prison. (20) How Joseph was exalted and made lord of all Egypt. (21) Joseph’s brothers come to Egypt. (22) Joseph’s brothers come to Egypt for the second time. (23) Joseph makes himself known to his brothers and treats them kindly. (24) Jacob again sees his son Joseph. (25) Jacob blesses his children and dies. (26) Joseph’s death. — To these, the festival stories are again added at the proper time, so that at least the most important facts of the life of Christ are touched upon also in this grade.

\emph{Third Grade:} In the third year the stories given above under grade 1 are again gone through. As a rule, one Bible story is gone through for each week of instruction; sometimes two weeks will be necessary. Allowing for the proper number of weeks for the review, which must set in after groups of stories belonging together have been taught, the usual number of weeks of school (about 36) is just sufficient for the material given above. The Bible verses, prayers, and hymns, which must be brought in relation to the Bible story should be assigned for study also on other days of the week, but never before light has been cast upon them from the Bible story to which they belong, and their recitation must always be introduced by questions bringing
\quellseite{436}
out their relation to the Bible story. The most important memory material for the three grades of the first course is given in the following, and should be brought in connection with the instruction in Bible History in the manner suggested above (compare \emph{Wartburg Lesson Helps, Primary Grade}):

\emph{First Year or Grade:} Bible verses: Ps. 115:3; 1 John 4:19; Ps. 5:4; 133:1; Matt. 5:9; Ps. 37:5; 37:37; 50:15; John 3:16; Prov. 23:26; Luke 11:28; Mark. 10:14; Prov. 11:1; John 1:29. Prayers: Morning prayer, prayer before meat, prayer after meat, evening prayer. Hymns: As each happy Christmas; Good news from heav’n, stanzas 1, 2; O sacred head now wounded, 6th stanza; Abide in grace, Lord Jesus, 1st stanza.—For the second year or grade: Bible verses: 1 John 5:3; Ps. 26:8; 1 John 4:16; Gen. 17:1; Ps. 33:9; 36:5; 34:7; 103:2; Luke 19:10; 1 John 1:7; Ps. 51:10; Gen. 39:9; Prov. 1:8; Matt. 5:8; Eph. 4:25; 1 Cor. 6:14; Romans 14:8; 8:28; Ps. 118:8; Ps. 91:11; 127:1; Mark 5:36; Romans 12:5. Prayers and Hymns: Now to gain a night’s repose, stanzas 1, 2; Ah, dearest Jesus, holy child, make Thee a bed soft, undefiled, within my heart, that it may be a quiet chamber kept for Thee; Awake my heart rejoicing, 1st stanza; Praise God the Lord, ye sons of men, stanzas 1, 2; Thou whose almighty word; If thou but suffer God to guide thee, stanzas 1, 3; I could not do without thee, 2nd stanza; I am Jesus’ little lamb, stanzas 1, 2; Abide in grace, Lord Jesus, stanzas 5, 6; Jesus, still lead on, 1st stanza; Glory be to Jesus, who in bitter pains, stanzas 1, 2; My faith looks up to Thee, stanzas 1, 2. Catechism: The Ten Commandments without explanation.—For the third year or grade: Bible verses: Ps. 119:105; John 5:39; Ps. 23:1; Ps. 23:2; Ps. 23:4; Ps. 104:24; Ps. 103:8; Gen. 8:22; Gen. 32:10; 8:21; John 14:6; Matt. 28:20; 1 Tim. 2:4; 2 Tim. 1:10; 2 Peter 3:9; 1 Peter 5:9; 1 Tim. 5:4; Hebrews 13:7; 1 Peter 3:9; Isa. 58:7; 1 John 2:15; Lev. 19:2; Gal. 6:7; Ps. 68:20; 2 Cor. 5:10; Isa. 28:29; John 5:28, 29. Prayers and Hymns: The Lord’s Prayer, Good news from heav’n, stanzas 1–6; Praise God, from whom all blessings flow; Jesus, Thy blood and righteousness; Jesus, now Thine own for ever, 1st stanza; Blessed Jesus, at Thy word, 1st stanza; Now the day is over; Glory be to Jesus; I am Jesus’ little lamb; He lives, my Lord has left the grave; O Holy Spirit, enter in, 1st stanza; Catechism: The first three Chief Parts without explanation. The explanation may be added to a few of the commandments. In each grade, the memory material previously learned is to be reviewed, and again brought into connection with the Bible story to which it belongs.

\quellseite{437}
\emph{Second or Intermediate Course} (comprising 2 school years or grades, age 10 and 11). \emph{First year or grade:} Bible History: History of creation and early history of man; Review of the history of the Patriarchs; stories of Moses; stories of Joshua; stories of the Judges; stories of Saul, David, and Solomon (the difficult stories may be omitted in this grade). Where time permits, a few stories of Elijah may be added. Review of the festival stories.—Hymns: Praise God the Lord, ye sons of men; Praise ye the Lord in simple joyous measure; My faith looks up to Thee; All hail the pow’r of Jesus’ name; Rest of the weary; Jesus loves me; Jesus lives, thy terrors now; If thou but suffer God to guide thee; Thy way and all thy sorrows; Now thank we all our God; Praise to the Lord, the almighty. Bible verses: One half of the verses marked with two stars in our “Catechism with Explanation”, and review of those learned in the First or Primary Course.—Catechism: The First and Second Chief Parts with explanation.

\emph{Second Year or Grade:} Biblical History: Stories from the life of Christ in groups f. i.: John the forerunner of Christ; The Savior’s youth; Jesus’ public ministry; Jesus the helper of the sick; Jesus the Savior in sorrow and death; Jesus the Savior of sinners; Jesus the teacher of the people; Christ’s success; Jesus disputes with the Pharisees; Jesus’ suffering and victory. Hymns: O how shall I receive Thee, stanzas 1, 2; My dear Jesus, I’ll not leave; Come, follow me, the Savior spake; My life is hid in Jesus, stanzas 1, 2; Jesus sinners doth receive; O Thou that hear’st when sinners cry; O sacred head now wounded, stanzas 1, 8, 10; Abide in grace, Lord Jesus.—Bible verses: One half of the verses marked with two stars in our catechism.—Catechism: The first three Chief Parts with explanation. Morning, evening and table prayers.

N. B. As in Course 1, the material given above is to be brought into connection with the instruction in Bible History; therefore the catechist must plan beforehand how this can best be done, in order to enable him to assign the memory material in the proper order. Compare \emph{Wartburg Lesson Helps, Second Course}. Intermediate Grade, and the Bible History published by the author.

\emph{Third or Upper Course} (comprising 2 years or grades, age 12 and 13). \emph{First Year or Grade:} Biblical History: The history of salvation in the Old Testament to Solomon and the division of the kingdom. History of salvation from the division of the kingdom to the restoration (compare also the Appendix to the Old Testament portion of the Bible History); with special emphasis on the labors and books
\quellseite{438}
of the Prophets.—Bible Reading.—Bible Study of the Old Testament.—Catechism: In this grade special periods for Catechism are set aside in which the Ten Commandments and the First Article are gone through according to the Catechism with explanation. Hymns and Bible verses are brought into connection with such instruction as much as possible, also Bible reading.—Bible verses: The verses marked with a cross in our catechism.—Hymns: O sacred head, now wounded; O Holy Spirit, enter in; Jesus Christ, my sure defense; A mighty fortress is our God.

\emph{Second Year or Grade:} Biblical History: The history of salvation in the New Testament to the ascension of Christ; The outpouring of the Holy Ghost and the establishing of the Christian Church by the Apostles, with special emphasis upon the Epistles. Sketches from Church History (see the appendix to the Catechism with explanation). Bible Reading. Introduction to the New Testament.—Catechism: Explanation of the Second and Third Article, and of the Third Chief Part. It may, however, be desirable, that the first three Chief Parts are completely treated in the second as well as in the first year. Memorizing of the Fourth and Fifth Chief Parts, also of the Part on Confession and Absolution, if this has not been done in the first year. Bible verses: Review of all Bible verses, and memorizing of longer portions of Holy Scriptures. Compare Ch. 27.—Hymns: Why, my soul, thus trembling ever; Who knows how near my life's expended; With all my heart I love Thee, Lord; What our Father does is well; The Church's one foundation.

Such distribution of material being calculated for a regular parish school, the question arises how the material is to be distributed where the establishment of a regular parish school has proved inexpedient. If there is a partial parish school, so arranged that children from the sixth to the ninth year are in attendance, that share of the material which, according to the preceding sketch, is allotted to the First Course of a regular parish school, falls to such a partial school, special attention to be paid to the simple, wholesome juvenile hymns, as found in the author's For Beginners (translated by H. Brueckner), or in a good juvenile hymnal. The material for the Middle Course is thereupon disposed of in the combination
\quellseite{439}
school (Summer, Sunday, Saturday), while the most important parts of the material assigned to the Upper Course of a regular parish school are gone over during catechumenal instruction, especially a survey of Biblical History, a number of Bible passages with explanation, and the essential facts relative to the church year, the hymn, and the order of service. If there is a partial parish school so arranged that children from the twelfth to the thirteenth year attend it, the combination school should assume in the preceding years the most important elements of the material assigned to the Lower and Middle Courses while that assigned to the Upper Course now falls to the one-course parish school.

If success has crowned the effort to win the public school authorities for the plan to allot some of the time claimed by the secular school to a representative of the Church trained in pedagogics, the material intended for the three courses of the regular parish school will constitute that for the instruction in religion thus made possible, while the material for catechumenal instruction will be the same in that case as that allotted to catechetical instruction in connection with a regular parish school.

Should, however, anyone draw the conclusion that the adjustment above proposed obviates the necessity for a regular parish school, we reply again: Without a regular, well conducted parish school, attended by all the children of the parish, many things are taught with too much haste and, therefore, cannot be made as productive for the whole man as they should; without it, especially in view of the frequent absence of the Christian spirit from the home, religious habit—a most important factor—does not come into its own; without it, the secular branches are not taught in Christ’s spirit, with the result that two conflicting views of the world compete for mastery in the soul, or even that the one prevailing in the public school may be victorious (cf. § 20). Notwithstanding,
\quellseite{440}
we concede that this perplexing feature can be, partly at least, removed in the same degree in which the Church succeeds in penetrating the home atmosphere and in securing for the catechetical instruction given during public school hours not only scientifically trained men, but also Christian characters of the Lutheran type. These latter, however, “do not fall from heaven”, but are the result of training, true as it is that they are a gift of God. Only the best human material produced by the Church is good enough for a function so fraught with possibilities for the future as this. If the ideal, a regular parish school, lies beyond the horizon of the attainable, let there be speedy, energetic, and joint efforts to make the best of the substitutes proposed by us. No complaints, no paper resolutions will here avail, but only purposeful, unwearying labor and prayer for an effective eye salve for oneself, the congregation, and the leading men of the Church. If despite all efforts and prayers, the obstacles in the way of a parish school, complete or partial, cannot be surmounted, let that be effected which is practicable in all circumstances, even in the mission field,—the combination of Summer, Sunday, and Saturday Schools with additional hours of instruction for catechumens; and let this opportunity, in a purposeful manner, be turned to account!

\trenner

\emph{The Sunday School}, in accordance with the mental development of the pupils, is divided into three departments of which the first comprises the pupils of infant age, the second those in the childhood stage, the third the youth of both sexes. Conditions render it advisable to assign all the children up to the seventh year to the First, or Beginners’, Department. The Second Department, comprising the children from the eighth to the thirteenth year, is divided into three, or at least two, grades. And the confirmed constitute the Third, or Senior, Department. Where there is no regular parish school
\quellseite{441}
or one not patronized by all the children of the parish, the Second Course requires the most thorough attention. That Biblical History, as also in the first, constitutes the principal material for instruction, is self-evident. “In all courses, the Bible story is fundamental; for the Bible story reveals the acts of God, in which our salvation, divinely wrought, is comprehended” (Schueren, Gedanken etc., p. 9; cf. above pp. 250, 252 f., 260). Scripture passages, the Catechism, and the hymn should be connected with it for the purpose of summarizing whatever truth the Bible story contains; but first rank must be assigned to Biblical History. Separate instruction in the Catechism, moreover, would be out of the question in view of the teaching force at the command of the congregation under present circumstances.

\emph{The Second Course—Childhood}. That a continuous instruction in Biblical History in which the narratives of the Bible are taken up seriatim from the creation of the world to the missionary journeys of Paul, is not the province of the Sunday School, follows from the fact that its pupils belong to the most divergent stages of life and development. Some selection has to be made. That could be done in such a way that the material designed for the regular parish school, with a few modifications, is allotted to the Second Course of the Sunday School, especially since the pupils of this department, as of the parish school, are about of the same age and require to be divided into Primary, Intermediate, and Junior grades. That done, each grade would have to finish a certain task calculated for two years, with the result that the Second Course would be divided into three independent units. In this way the educational facilities would best adjust themselves to the mental development. These reasons have prompted Dr. Schmauk, in his \emph{Graded System of Sunday School} to assign to the children from ten to fourteen years Bible Story, Bible Readings, Biblical History, Bible Geography, and Bible
\quellseite{442}
Biography as appropriate material, which has largely been handled in a masterly manner. Where for each one of these grades there is a separate teacher who knows his business and never forgets that he has fulfilled his task not merely by introducing the lesson material into the concept world of his pupils, but has, at the same time, stirred their emotional and volitional life, we believe the adoption of the lesson plan of the parish school (with modification) or that proposed by Schmauk, or of some similar one, quite practicable and promising of good results. But the question remains whether those postulates in most cases really exist. Not seldom the Sunday School uses the same room as the congregation; in that case all its members, at all events, those of the Second Course, which we now have in mind, have to be instructed in a room not seldom without partitions, often very crowded. It meets but once a week for the space of one hour, in contrast to a parish school whose three grades, even when crowded under one teacher into one room, yet receive two hours’ instruction a week in Biblical History. Above all, while there are a number of assistants of both sexes in the Sunday School, these are frequently rather helpless while there is usually but one competent, purposeful leader, unable in the short time at his disposal, to teach all three grades himself. To be sure, the teaching force can and should be improved (while the summer sessions of the teachers’ institutes, some of which are Lutheran, are not likely to be attended by most of the teachers, there remains the weekly teachers’ meeting by way of preparation for Sunday); notwithstanding, the postulates of those lesson plans, which are necessary to success, will be realized but seldom. How can a teaching force such as we have at our command, independently present the several Bible stories, adapt them to the different grades, elaborate from the individual story the inherent truth for faith and life, and train the class for mental co-operation? In these
\quellseite{443}
circumstances instruction in Sunday School is likely to be a mere cramming of the external events of the story; but there will be no arousing of the sympathies and of religious and moral promptings—so important for development of the inner life (cf. pp. 222, 224), and through it, a stirring of the will.

This was felt by Hunt; for this reason he deems it important that the pastor should address the whole school every Sunday, “to impress a spiritual truth upon the hearts of the pupils” (in \emph{The Lutheran Sunday School Handbook}, p. 62). True, only the remedy is not adequate to the evil. The school having been occupied with divergent, disconnected material, the pastor cannot link such “spiritual truth” to the subject matter with which the several classes had been occupied; there is no effective turning to account of the latter; the “spiritual truth” presented by the pastor is felt by the class to be merely superimposed from without, and the intended effect cannot be realized. However, the needed stimulation of the intellectual, sympathetic, and above all, of the ethical and religious interests, and the consequent prompting of the will are best effected where the whole school, or at least the Second Course (age 8–12) considers the same story on one and the same Sunday, and where this has been related and presented by the leader first, so presented—according to the intuitive principle and with reference to the Catechism, Scripture, and hymn material—that it is not only understood by the whole department, but that the religious truths imbedded in the story are also developed from it for direct application. But does the gradual development, so important for mental progress and certainly entitled to consideration, in this way really come into its own? This is the case indeed; for (1) in contrast to the former lesson plans of the International Sunday School Association, it is only for the Second Course that one and the same Bible story is required for the same Sunday; (2) an efficient leader will so formulate his presentation
\quellseite{444}
that the members of the Primary Grade can understand it and that at the same time it will not become dull and insipid for the pupils of the Junior Grade; and (3) the three grades of the Second Course are given a printed lesson help, in which the story under consideration with the related material from Scripture, the Catechism and the Church Book is analyzed for each class separately, and careful attention is paid to the progressive stages of the children’s mental development (cf. Wartburg Lesson Helps, Second Course).

Where all three grades of the Second Course always have the same story, such narratives have to be left out of consideration from the outset as possess no feature that makes them intelligible to the Primary Grade. The test of their acceptibility is not that no unassimilated fragment remains; for in that case few would be left for even the Third Course; but those should be excluded for which the apperceptive fulcrum is altogether lacking in the concept world of little children (cf. p. 253). Accordingly all stories of the Old Testament should be left out of consideration which deal with the people of Israel in general; for the term “people” is foreign to the concept world of little children. But the family stories and episodes from the lives of eminent men should occupy the foreground; for it is in families that children originate, and even for the comprehension of a biography the requisite apperceptive fulcrum is found in the child’s concept sphere (cf. Schueren, op. cit., p. 10). Accordingly certain sections of the New Testament, such as Zacharias’ hymn of thanksgiving, Jesus’ talk with Nicodemus, the conversation of Jesus with the Samaritan woman, Jesus’ discourses on the end of the world, etc., are excluded from the outset as material for the Second Course. We have come to the limit of the efficiency of this agency, and it is well that we should recognize that fact.

Out of the question for us is the haphazard selection
\quellseite{445}
of the lessons now from this period, now from that, after the manner of the old lesson plans of the I.S.S.A. The stories of the Bible should be presented in large groups, in coherent, integral masses of thought, derived from no more than two, or at the highest, three periods. Only thus can the children really penetrate into the material; only thus is it possible for those Bible characters to become factors in the child’s inner life and stir his varied feelings. Splendid is what Fankhauser says upon this subject: “When we have begun to tell of some personage, for instance, Abraham, the child has taken an interest in him if the proper method has been pursued. There is a man who has gone from home, away from father, mother, friends, and relatives, without knowing what is to come of it, blindly following the command of God. Something has been wrong in our presentation if the children do not desire to find out from us the next time how that man who has won their sympathy, fared further. It is a fine thing, then, to be able to meet the desire of the child; that means we have gained a great advantage for our instruction at the very outset. We do not have to explain at length whom we are going to talk about, in what circumstances he lived, etc. — things surely necessary where an unknown person is to be introduced. It is just like a piece of embroidery work: mother has already placed it on hoops and threaded the needle; a few feathers of a many-colored tropical bird begin to appear in the fabric; and wondering what the figure will look like when it is completed, we anxiously watch mother’s moves till all the gorgeous colors are blended in the finished product” (\emph{Die biblische Geschichte in Sonntagsschule und Religionsstunde}, p. 18). Another thing, but one new story is to be selected for each Sunday. It is true that there are still Biblical Histories which harness two or three stories mechanically together in order to produce 2x52 lessons, according to the example of old
\quellseite{446}
Huebner. From the pedagogic standpoint, this is illsuited even to the parish school, much more so to the Sunday School where the “driver’s stick” must be still less in evidence. “Not too much should be given at one time; but what is given should not be a fragment but a coherent, well-rounded unit” (Schueren, \emph{op. cit.}, p. 17). How could anyone possibly think of launching forth a second story immediately after the first while the first story still holds the mind of the child enthralled and is only just about to find its way to his heart, especially if the second one portrays a quite different picture and has a different aim? The effect of the second picture would merely be to obliterate or efface the first; it would nullify the intended effect upon the mind (cf. p. 207). That proceeding would tend to make the children superficial of mind and dull of soul. Is it not the imperative duty of the teacher to focus the attention of his pupils on one object until a clear intuition or a clear concept has been formed? And does he not conscientiously try to fulfill this duty even if the lightning like rapidity with which his pupils’ attention shifts from one object to another, renders it now and then very difficult for him to do it? Does he not at all times keep in mind the law of the narrowness of consciousness? How, then, should he suddenly become a merciless driver, chasing the poor children from one story to another and never permitting the concepts that have been aroused, to take firm root in the soul! (Cf. Fankhauser, \emph{op. cit.}, p. 21.) Only if the several stories constitute a catechetical unit, may they be presented in one lesson.

But what order is to be observed in the selection of the several stories? He who knows what a precious aid the church year is in the inculcation of the essential facts of salvation (cf. p. 375), and for this reason is concerned about introducing the young to the idea of the church year and preparing them, also in this direction, for intelligent participation
\quellseite{447}
in the worship of the mature congregation, will follow the order of the church year also in the Sunday School, and always present, at least during the festive half of the church year, stories from the New Testament, with special emphasis upon the outstanding events in the life of Jesus; while the other half of the church year may be devoted to the story of the Patriarchs, of Moses, etc., that is such units as are required by the curriculum based upon the successive cultural epochs. — Finally, because the Sunday School, especially where there is no parish school, has to be a school rather than a service of worship, and because we do not merely wish to arouse a transient emotion in the soul, but to impart a treasure for life, care must be taken that familiarity with the subject-matter is attained. To this end there will not only be a review every seventh or eighth Sunday, but there will be also a resumption of the same topics every third or fourth year. This, from another point of view, brings us back to the requirement: the three or two grades of the Second Course have the same lesson story though adapted to the various stages of development represented by the several grades. If the objection is raised that the recurrence of the same material every third year kills interest, we reply with Schueren: “He who thinks that the repeated relation of a story causes ennui, is not acquainted either with the character of the Bible story or of the child. Where the real story-telling note is struck even in a moderate measure, the biblical story will always appear new to the children; with increasing interest they will listen to the person telling a story that has caused them joy once before. How is a mother who has the story-telling gift to any degree whatever, tormented by the children to tell the same story for the tenth or twentieth time! How the children hang upon her lips when she tells the story! It impresses them as if it had never been told! Teachers whose experience is different should not accuse the children but themselves” (\emph{op. cit.}, p. 10).

\quellseite{448}
In the \emph{Wartburg Lesson Helps}, the following stories constitute the Second Course: First Year: 1. How God created the world.—2. How God made man happy.—3. How man sinned.—4. Through sin man was lost and condemned.—5. How the Savior came into the world.—6. The child Jesus is brought to the temple.—7. Wise men from the East come to the child.—8. How the child Jesus was saved from great danger.—9. Jesus, when a boy twelve years of age, again comes to the temple.—10. How Jesus changed water into wine at the wedding at Cana.—11. How Jesus blessed Peter’s draught of fishes.—12. How Jesus saved His disciples during the storm.—13. Jesus feeds the hungry multitude.—14. How Jesus healed the man sick of palsy both in body and in soul.—15. How Jesus comforted a sorrowing father by restoring his daughter to life.—16. Jesus desires that little children should be brought to Him.—17. Jesus is our Good Shepherd.—18. How Jesus died for us on Good Friday.—19. How Jesus arose from the dead and appeared to Mary Magdalene.—20. Jesus would help the greatest of sinners, if they would only come to Him.—21. Jesus tells who the neighbor is, whom we should help.—22. Jesus tells us to forgive those who sin against us.—23. Jesus tells us to pray.—24. Jesus ascends into heaven in order to be with us always.—25. Jesus sends the Holy Spirit, who shall lead us to Jesus.—26. Of godly Abraham who by faith left his home.—27. Abraham would rather separate from Lot than live in strife with him.—28. How Abraham proved to be a friend in need.—29. What pleased God most in Abraham.—30. How Abraham spake with God as a friend speaks with a friend.—31. How Sodom and Gomorrah were destroyed because of their sin.—32. God gives Abraham a son, but Abraham loves God more than his son.—33. How God provided a wife for Isaac.—34. How Jacob deceived his father and his brother.—35. How Jacob fled from his brother and in his dream saw the heaven open.—36. How God prospered Jacob in the strange land and safely brought him home again.—37. Of Joseph, Jacob’s dearest son.—38. How Joseph was sold by his brethren for twenty pieces of silver.—39. How Joseph was humbled and imprisoned.—40. How Joseph was exalted and made lord of all Egypt.—41. Joseph’s brethren come to Egypt for the first time.—42. Joseph’s brethren come to Egypt for the second time.—43. Joseph makes himself known to his brethren.—44. Joseph cares for his father in Egypt.

\emph{The Second Year:} 1. How sin came into the world.—2. How sin spread so rapidly.—3. How God destroyed man because of sin.—4. How the Saviour’s messenger was born.—5. How the Saviour came into the world.—6. How John prepared the way for Jesus and baptized Him.—7. How Jesus resisted the devil’s attempt
\quellseite{449}
to dissuade Him from being our Saviour.—8. How Jesus healed the centurion’s servant, and encountered great faith.—9. How Jesus restored a son to his mother.—10. How Jesus heard the prayer of a heathen woman.—11. How Jesus entered the house of a sinner and made a new man out of him.—12. How Jesus entered Jerusalem on Palm Sunday.—13. How the Lord’s Supper was instituted.—14. How hard it was for Jesus to suffer and die for us.—15. How Jesus was betrayed by one of His disciples.—16. How one of His disciples denied Jesus.—17. How Jesus was carried from one court to another on account of our sin.—18. How Jesus died and was buried on Good Friday.—19. How Jesus arose from the dead and appeared to Mary Magdalene.—20. How Jesus comforted two disciples on their ways to Emmaus.—21. How Jesus appeared to His disciples and brought Thomas to believe in His resurrection.—22. How Jesus prepared salvation for all, though believers only partake in it.—23. How he only will be saved who perseveres in faith.—24. How Jesus ascends into heaven to be with us always.—25. How Jesus sends the Holy Spirit that He may lead us to Christ.—26. How a small boy was wonderfully preserved and received the name Moses.—27. How God appeared to Moses and called him to be the saviour of the people of Israel.—28. What Moses accomplished before proud Pharaoh.—29. How Moses led the people of Israel through the Red Sea.—30. How God cared for His people on their journey.—31. How God gave the Ten Commandments to the people of Israel.—32. How Moses took leave of the people of Israel and died.—33. How Joshua brought the people of Israel into the land of Canaan.—34. Of Eli’s wicked sons and pious Samuel.—35. How Samuel appointed Saul to become king.—36. How God rejected Saul because of his disobedience.—37. How God selected pious David to be king.—38. How David trusted in God and slew Goliath.—39. How David suffered much from Saul, though he found a good friend in Jonathan.—40. How David refrained from rendering evil for evil.—41. How God exalted David and made him king.—42. How Absalom sorely grieved David.—43. How Solomon, a king after God’s heart, built the temple in Jerusalem.—44. How God permitted a famine to come over Israel and miraculously fed the prophet Elijah.

The First Course or Beginners’ Department comprises the children from five to seven years of age. No regular lesson can be assigned them. They must be entertained rather than instructed, yet so entertained that they are led into an understanding of religious fundamentals, and the new life implanted
\quellseite{450}
in Baptism is nourished and made to become active. In the author’s book, For Beginners, it is practically worked out, how these children are to be instructed. It contains the following lessons:

1. God made all things.—2. God made Adam, the first man.—3. What God has given me.—4. God cares for me.—5. God protects us by His angels.—6. God protects us by His angels.—7. How God has made me His child.—8. How we may speak with our Heavenly Father.—9. How we may speak with our Heavenly Father.—10. God hears our prayers.—11. Jesus is our Good Shepherd.—12. God sees and knows everything.—13. God knows and sees everything.—14. How God punishes us if we do evil.—15. How God punishes us if we do evil.—16. God keeps His promise.—17. We should thank God when He has helped us.—18. God looks upon our hearts.—19. What we should do on Sunday.—20. We should help all those who are in need.—21. The Bible is the best book.—22. How beautiful it is in heaven.—1. Christmas is coming.—2. An angel comes to the virgin Mary.—3. God gives the child Jesus to the virgin Mary.—4. The angels come to the shepherds.—5. The shepherds come to see the child Jesus.—6. How the shepherds made known the birth of Jesus.—7. The wise men seeking Jesus.—8. The wise men find and worship Jesus.—9. God protects the child Jesus from Herod.—10. How God and man are pleased with the child Jesus.—11. The boy Jesus goes to celebrate Easter.—12. The boy Jesus in the temple hears the word of His father in heaven.—13. Mary finds the boy in the temple.—14. Jesus finds His first disciples.—15. How Jesus changes water into wine.—16. Jesus feeds the hungry.—17. The winds and the seas obey Jesus.—18. Jesus heals the sick man.—19. Jesus opens the eyes of two blind men.—20. Jesus raises the daughter of Jairus to life.—21. How the Lord Jesus blessed children.—22. Jesus enters Jerusalem on Palm Sunday.—23. How Jesus prayed in the Garden of Gethsemane.—24. Jesus is taken prisoner.—25. The Lord Jesus is crowned with thorns.—26. Jesus bears His cross to Golgotha.—27. Jesus dies on the cross for us.—28. The Lord Jesus is buried.—29. Jesus is raised from the dead and shows Himself to Mary Magdalene.—30. Jesus is taken up into heaven.

The Third Course, or Senior Department (early and later adolescence) of the Sunday School is chiefly concerned with introducing the pupils to the Bible, to explain it, and to train for a personal use of it. Where introduction to the history of the Church and her labors is not assigned to the Young
\quellseite{451}
People’s Society, the place for that, too, would be here. At the very outset it should be borne in mind that a dry introduction to the extent of the canon, the series of books composing it, details as to the origin of the canon, and the history of the Bible—all necessary things to be known—accomplishes very little after all. The young must be led into the Scriptures themselves, they should read them again and again. While the necessary explanations should be forthcoming in connection with the reading of the Scriptures, they should not bear the character of an exegesis, such as would be in place in a theological seminary; but it should most carefully bring out the essentials and fundamentals, show their bearing upon the lessons of the Catechism, and always aim at practical results in the lives of the young people. Whether it may not be expedient here and there to take a shorter route through the whole Bible at the outset, to be followed by a more thorough course afterward, depends upon circumstances.

The Saturday School should go hand in hand with the Sunday School, paying particular attention to the “penetration” and “application” of the Bible stories which have been “presented” in the Sunday School. How that is to be done, can be seen from § 34. Here time is afforded for an analysis of the Catechism, for the necessary explanation of its terms, and a preparation by these means for more thorough subsequent study of the Catechism, always, be it remembered, in connection with Biblical History. Here is also the occasion for practicing church hymns and for their preliminary explanation. We know from practical experience how, in this way, the pastor gets the Sunday School really under his control and effectually supplements its work.

The great value of the Summer School is found in this that it affords a daily opportunity for handling religious material with the pupils. The teacher is enabled to introduce his pupils to larger complexes; to delineate with them characters,
\quellseite{452}
that of Abraham, Jacob, Joseph, Moses, David; to survey the unfolding of the prophetic ministry of Christ: only thus is it possible for the child to identify himself with another age and another character; only thus it is possible for him to enter upon an “ideal association” with the persons of sacred history. Separate instruction in the Catechism should be given in this school to the older pupils. The Ten Commandments especially need to be stressed in this connection, in order to acquaint the children with the order which the Father in heaven wants observed in His house (cf. p. 313). If it be true that the explanations of the Catechism should never smack of doctrinal barrenness, but always throb with life, this is true of the Commandments in particular. Instead of insisting upon definitions of the many conceptions occurring in the First Chief Part, the children should be made to visualize pictures from life; there they will be shown in what Christian life and conduct consists. Alongside the First Chief Part, the Second might receive attention, but only in outline; it is sufficient so to explain it that it can be intelligently memorized. Narratives from the history of the Church and Missions, the fundamentals of Biblical Geography, the practise and discussion of the more important church hymns—these are subjects that will be used to supplement the lesson material so far mentioned.

The instruction of catechumens is bound to have its center in the Catechism. If a parish school is available and thorough instruction has been given in the first three Chief Parts, it is practicable in these circumstances to abandon the regular order of the Five Chief Parts and to select another way. Confirmation, for which the catechumens are to be prepared, may be made the starting point; from there one may go back to Baptism, proceed to the Creed next, and finally, at the hand of the Decalogue and the Lord’s Prayer, describe the Christian’s walk, the whole process to issue in Holy Communion and Confession. Where there is a thorough knowledge of Biblical History as
\quellseite{453}
well, the following process, too, is feasible: Creation and the fall of man may be made the starting point, with which the First Article readily connects itself, whereupon, after a brief glance at the redemptive purpose (cf. Gen. 3:15), the story of salvation down to the resurrection of Christ may be gone through with the pupils (cf. pp. 299–307). Here a stop is made and the Second Article taken up and the story of salvation pursued down to Pentecost, whereupon the Third Article is explained and the Sacrament of Baptism. Then, at the hand of the Decalogue and the Lord’s Prayer, the conduct of the Christian may be described, with special reference to the example of Christ, and instruction would conclude with an explanation of Confession, the Holy Supper, and Confirmation. Where neither of these postulates exist, the instruction of catechumens will closely follow the parts of the Catechism: it will bring out particularly the essentials, with a constant view to the practical Christian life and the spirit permeating it, a pastoral atmosphere pervading the whole process. Such instruction will further make a beginning of introducing the pupil to the Scriptures, cast light upon Lutheran usage, the church year, the order of service, and give practical information regarding the tasks of the Christian congregation (especially missions). Compare O. Hardeland, \emph{52 Konfirmandenstunden}, ⁴1910; R. Steinmetz, \emph{Das gute Bekenntnis, Hilfsbuch zur Vorbereitung auf die Konfirmation}; Steinmetz, \emph{Die Bereitung zur Konfirmation in Lehre und Leitung}, 1910; Vogel, \emph{Seelsorgerlicher Konfirmandenunterricht}, 1911.

In the Bible Class of the Sunday School (cf. what was said above) and in the Young People’s Society, the object must be an introduction to the Scriptures and to the secondary material of instruction as it was called in § 29. How to introduce the pupils to the Holy Scriptures, has already been indicated (p. 372); it is practically demonstrated in the Third Course of the Wartburg Lesson Helps (Book of Books; Book
\quellseite{454}
of Life), composed by the author of this textbook. How the secondary material is to be used in Young People’s Societies we show by adding several lesson plans. In regard to presentation compare Part V. In part also this material has been worked out practically, for instance in the author’s book, \emph{\emph{The Life of Dr. Martin Luther, Sketched for Young People’s Societies, and the Necessary Directions for General Discussion Appended}}, 1917, and in the department \emph{Luther (or Wartburg) League Topics} in the Lutheran Herald.

1. The Church and Her History: The Heathen World and its Moral Condition at the Time of Christ; The Preaching of the Gospel Among the Heathen Through the Apostles; The Persecution of Christians under Nero; Polycarp and Justin; The Persecution in Gaul; Perpetua and Felicitas; Persecution under Diocletian; Constantine and the Victory of Christianity; Julian the Apostate; Pictures from the Life of the Early Church; Athanasius and the Doctrinal Conflicts of his Time; Antonius and Monasticism; Chrysostom and his Congregation in Constantinople; The Origin of Papacy; Mohammed and the Judgment of God upon the Church; The Gospel in Germany; Charlemagne; Missions North and East; Gregory VII and Henry IV; The Crusades; The Monastic Orders; The Decay of the Church; The Waldensians; Wiclif and Hus.—Luther; Melanchthon; Zwingli and Calvin; The Gospel in England; The Gospel in France; The Counter-Reformation and the Thirty Years’ War; Paul Gerhardt; Duke Ernest of Gotha; Phil. Spener; Rationalism; The New Spring; Origin of the Union.—Muehlenberg and the Founding of the Lutheran Church in America; Rationalism in the Lutheran Church of America and the Founding of the General Synod; Founding of the Missouri Synod; Founding of the Iowa Synod; Founding of the General Council; Task of the Lutheran Church of America.

—2. The Church and Her Confessions: The Duty of the Church to Confess her Faith and the Three Oecumenical Confessions; The Lutheran Church and her Confessions; The Confessions of the Lutheran Church and the Reformed Church; The Lutheran Confession and the Union (with special reference to its representatives in America); The Lutheran Confession and Methodism; The Lutheran Confession and the Baptists; The Lutheran Confession and the Adventists; The Lutheran Confession and Christian Science.—3. The Church and her Problems: Church and World; Church and Politics; The Church and Political
\quellseite{455}
Parties; The Church and the Secular Press; The Church and the Labor Movement; The Church and Capital; The Church and Prohibition; The Church and the Lodge; The Church and the Public School; The Church and the Language Question.—4. The Church and Her Tasks: a) Foreign Missions: The Duty of Carrying on Foreign Missions; Ziegenbalg and Schwarz, and Lutheran Missions in India; Hans Egede the Apostle of Greenland; Zinzendorf and the Moravian Missions; Williams the South Sea Apostle; W. Carey the Missionary to the Bengalese; Guetzlaff the Founder of the Chinese Missions; The Leipzig Missions in India; L. Harms and the Hermannsburg Missions in Africa; Livingstone the Explorer and Missionary in the Dark Continent; Madagascar, the Grief and the Joy of Missions; The New Dettelsau Missions in Australia and New Guinea; Baierlein and Lutheran Missions Among the Indians; A Present Survey of Foreign Missions.—b) Inner Missions: Inner Missions the Duty of the Church of Jesus; A. H. Francke and Inner Missions; Muehlenberg’s Labors Among the American Dispersion; Wyneken’s Pioneer Labors; W. Loehe and His Efforts on behalf of the Lutherans in America; The Pioneers of the Iowa Synod; The Organization of a Mission Congregation.—W. Loehe and the Deaconess Work; Theodore Fliedner and Kaiserswert; Wichern and the Rauhe Haus in Hamburg; G. Mueller the Orphan Father of Bristol; Survey of Inner Missions in Germany; Passavant the Father of Lutheran Inner Missions in America; A Day in the Deaconess Home in Philadelphia or Milwaukee; A Day in the Orphans Home of Mount Vernon; A Day in the Orphans Home of Waverly; City Mission in Minneapolis; Our Homes for the Aged; Our Higher Institutions.—5. The Church and Her Usages: The Necessity for Fixed Usages in the Church; The Church Year; The Pericopes; The Celebration of the Sunday; The Order of the Regular Sunday Service; The Order of the Communion Service; The Order of the Baptismal Service; Announcement for Confession; Sponsorship; System in Giving.—6. The Church and Her Constitution: The Proper Organization of the Local Lutheran Congregation; The Ideal Constitution for the Local Congregation; The Necessity for Membership in a Local Congregation; Church Discipline; The Synod and Its Organization; Connection Between the Synod and the Congregation; The Correct Relation to other Synods and Its Congregations and Members.

In addition to the several series given above, such as these would be feasible: Bible Characters: Adam and Eve; The Sons of Adam
\quellseite{456}
and Eve, or The First Stage of Culture; Noah and the Flood; Nimrod and the Origin of World Power and of Paganism; Abraham; Isaac; Jacob; Joseph; Moses; Joshua; Samuel; David; Solomon; Elijah; Isaiah; Jeremiah; Daniel; Ezra; The Maccabees; Herod.—Jesus’ Life in the Light of the Ten Commandments; Mary the Mother of the Lord; Mary the Sister of Lazarus; Mary Magdalene; Peter; John; Judas; Stephen; James the Just; The First Christian Congregation; Ananias and Sapphira; Paul on the Way to Damascus; Paul in Lystra; Paul in Athens and Corinth; Paul in Caesarea; Paul in Rome; Timothy (here the sections of Scripture bearing upon the subject should be read by the youth by way of preparation).—Church Hymns Explained from the Aspect of Their Origin: The Church Hymn and Its Significance in Reformation Times: “A Mighty Fortress” and Parallels from Luther’s Life; Speratus and “To Us Salvation Now Has Come”; Ph. Nicolai and “Wake, Arise, the Voice Is Calling”; M. Rinkart and “Now Thank We All Our God”; P. Gerhardt and “Thy Way and All Thy Sorrows”; Neumark and “If Thou but Suffer God to Guide Thee”; Fleming and “In All that may Betide me, The Most High God Shall Guide Me” (In allen meinen Taten); John Heermann the Singer in the School of Suffering; John Frank the Godly Burgomaster and his Hymn “Deck Thyself, O Soul, with Gladness”; Gellert and “Jesus Lives! No Longer Now”; J. Sebastian Bach the Godly Organist; “Jesus Christ, My Sure Defense” and the Blessings in Its Wake.—Biblical, Ethical, and Apologetic Materials: The Old Testament and Its Origin; The New Testament and Its Origin; The History of the Bible in the Church; The Psalms and Their Significance for the Inner Life; The Prophetic Writings and their Significance for the Inner Life; How the Bible Should Be Read to Build up the Inner Life.—The Blessing of True Friendship; Youth’s Happy Time; The Special Dangers of Youth; The Necessity of Conversion; Prayer and Its Blessing; The Highest Blessings of All; Conscience and Its Cultivation; Truthfulness a Fundamental Condition of True Piety; Healthy Recreation; Getting Ready for the Last Journey; Our Duty Toward the Dead; the Three Articles of Our Christian Faith and the Intellectual Movements of the Present; Is It True that Only Fools Believe?—Pictures from the Holy Land might similarly do service.

